% ===========================================================================
%  LYRA — GUIDA DELLO SVILUPPATORE
% ===========================================================================
%
%  Uno dei due documenti del progetto. L'altro è user-guide.tex (il manuale
%  utente, in inglese) e condivide con questo tutto ciò che sta in comune/.
%
%  Compilazione:      make sviluppatore        (oppure: latexmk -pdf guida-sviluppatore)
%  Entrambi:          make
%  Pulizia:           make pulisci
%
%  STRUTTURA. Un capitolo per cartella sotto sviluppatore/, e dentro ogni
%  cartella un file per sezione. Il file 00-capitolo.tex contiene SOLO il
%  titolo del capitolo e l'elenco delle sue sezioni: nessun testo. Così si
%  sposta una sezione spostando un file, e due persone scrivono due sezioni
%  dello stesso capitolo senza toccare lo stesso file.
% ===========================================================================
\documentclass[11pt, a4paper, twoside, openany]{book}

% ===========================================================================
%  comune/preambolo.tex — impostazioni comuni ai due documenti
% ===========================================================================
%
%  Lo includono sia guida-sviluppatore.tex sia manuale-utente.tex. Ciò che
%  vale per uno solo dei due NON va qui: va nel suo file principale, altrimenti
%  fra sei mesi nessuno sa più perché una cosa è dove è.
%
%  Motore: pdfLaTeX. Non lualatex, non xelatex, e la scelta è deliberata —
%  questo preambolo usa solo pacchetti di una TeX Live di base, così chi riceve
%  la cartella compila senza installare nulla. Nessun font di sistema, nessun
%  --shell-escape, nessuna dipendenza da Python.
% ===========================================================================

% --- lingua e codifica ------------------------------------------------------
\usepackage[T1]{fontenc}
\usepackage[utf8]{inputenc}

%  L'italiano in due modi, perché non tutte le installazioni hanno il pacchetto
%  babel-italian (le immagini Docker di TeX Live «base» tipicamente no, e una
%  documentazione che non compila lì è inutile metà delle volte che serve).
%    - c'è italian.ldf  -> si usa quello, è la via normale;
%    - non c'è          -> babel >= 3.75 costruisce la lingua dal proprio file
%                          locale/it/babel-it.ini, che è sempre presente.
%  In entrambi i casi si ottengono «Capitolo», «Indice», «Figura», «Tabella».
%  La sillabazione italiana è un'altra cosa: sta nel FORMATO, non nel pacchetto,
%  e se manca si vede subito (parole spezzate all'inglese). Il controllo è in
%  fondo a questo file.
\IfFileExists{italian.ldf}
  {\usepackage[italian]{babel}}
  {\usepackage[italian, provide = *]{babel}}

\usepackage{csquotes}

% --- caratteri --------------------------------------------------------------
%  Palatino per il testo e per la matematica (mathpazo): la documentazione ha
%  formule — vol targeting, restringimento della covarianza — e un testo con la
%  matematica di un'altra famiglia si legge come due documenti incollati.
%  Helvetica ridotta per i titoli, monospazio di Latin Modern per il codice.
\usepackage{mathpazo}
\usepackage[scaled=0.90]{helvet}
\usepackage{lmodern}
\renewcommand{\ttdefault}{lmtt}

\usepackage{microtype}

% --- composizione del paragrafo ---------------------------------------------
%  Un documento pieno di nomi di file e di simboli in monospazio ha molte
%  parole che non si possono spezzare: `diagnose_regime_signals.py` e' lungo
%  quattro centimetri e non ha punti di sillabazione. Con i valori standard
%  TeX preferisce far sbordare la riga nel margine piuttosto che allargare
%  gli spazi, e il risultato si vede.
%
%  Questi tre valori invertono la preferenza: meglio una riga un po' larga di
%  spazi che una riga che esce dalla giustezza. `hbadness` alza la soglia di
%  segnalazione di conseguenza, cosi' il log continua a riportare i problemi
%  veri e non le righe lente che abbiamo appena autorizzato.
\tolerance = 1500
\emergencystretch = 3em
\hbadness = 2000

% --- impaginazione ----------------------------------------------------------
\usepackage[
  a4paper,
  inner = 28mm, outer = 24mm, top = 26mm, bottom = 28mm,
  headsep = 7mm, footskip = 12mm,
]{geometry}

\usepackage{emptypage}      % le pagine lasciate bianche restano davvero bianche

% Niente pagine bianche. I due documenti si leggono a schermo: la pagina
% vuota che il libro inserisce per far cominciare capitoli, parti e
% frontespizio sulla pagina di destra lì è solo un buco. Con l'opzione
% openany della classe i capitoli e le parti cominciano sulla pagina
% seguente; qui lo stesso vale per \frontmatter, \mainmatter e i
% \cleardoublepage scritti a mano (frontespizi), che la classe book
% eseguirebbe comunque. Il layout resta a due facciate (margini interno ed
% esterno, intestazioni alternate).
\let\cleardoublepage\clearpage
\usepackage{setspace}
\setstretch{1.06}

% --- colori, codice, figure -------------------------------------------------
\usepackage{etoolbox}
% ===========================================================================
%  comune/colori.tex — la tavolozza
% ===========================================================================
%
%  Sono i colori del PRODOTTO, non colori scelti per la documentazione: gli
%  stessi valori che stanno in dashboard/static/css/ui-v2.css e nella landing.
%  Documentare un prodotto con i colori di qualcun altro costringe chi legge a
%  tradurre due volte — una figura verde acqua qui e un pulsante verde acqua
%  nell'app sono la stessa cosa, e si deve vedere.
%
%  Chi cambia la tavolozza del prodotto cambia QUESTO file e nient'altro:
%  nessun altro file della documentazione nomina un colore per esteso.
% ===========================================================================
\RequirePackage{xcolor}

% --- i colori di base della landing -----------------------------------------
\definecolor{lyraInchiostro}{HTML}{111318}  % testo, titoli
\definecolor{lyraPannello}{HTML}{F5F5F2}    % fondo dei riquadri
\definecolor{lyraBordo}{HTML}{ECECEA}       % filetti e bordi
\definecolor{lyraAccento}{HTML}{2BB59A}     % verde acqua: l'unico accento
\definecolor{lyraPositivo}{HTML}{0F7A62}    % numeri in utile
\definecolor{lyraNegativo}{HTML}{9A3412}    % numeri in perdita

% --- i pastelli, usati per distinguere gli strati nelle figure --------------
\definecolor{lyraCielo}{HTML}{D9F1F7}
\definecolor{lyraMenta}{HTML}{DDF3E4}
\definecolor{lyraPesca}{HTML}{FFE4D4}
\definecolor{lyraLavanda}{HTML}{E6DEFF}
\definecolor{lyraCrema}{HTML}{FFF1C9}

% --- derivati, per bordi e testi sopra i pastelli ---------------------------
%  Si ricavano dai precedenti invece di essere scritti a mano: cambiando un
%  pastello, il suo bordo lo segue.
\colorlet{lyraCieloBordo}{lyraCielo!55!lyraInchiostro}
\colorlet{lyraMentaBordo}{lyraMenta!55!lyraInchiostro}
\colorlet{lyraPescaBordo}{lyraPesca!55!lyraInchiostro}
\colorlet{lyraLavandaBordo}{lyraLavanda!55!lyraInchiostro}
\colorlet{lyraCremaBordo}{lyraCrema!55!lyraInchiostro}

\colorlet{lyraGrigio}{lyraInchiostro!55}
\colorlet{lyraGrigioChiaro}{lyraInchiostro!28}
\colorlet{lyraAccentoScuro}{lyraAccento!70!lyraInchiostro}
\colorlet{lyraAccentoTenue}{lyraAccento!14!white}

% --- i colori del codice ----------------------------------------------------
\colorlet{codFondo}{lyraPannello}
\colorlet{codBordo}{lyraBordo}
\colorlet{codCommento}{lyraInchiostro!45}
\colorlet{codParola}{lyraAccentoScuro}
\colorlet{codStringa}{lyraNegativo!85}
\colorlet{codNumero}{lyraInchiostro!70}
\colorlet{codRiga}{lyraInchiostro!35}

% ===========================================================================
%  comune/stile-tikz.tex — il vocabolario grafico delle figure
% ===========================================================================
%
%  Ogni figura del progetto disegna con QUESTI stili e non con parametri
%  propri. La ragione e' leggibilita': se in una figura il riquadro di un
%  modulo e' azzurro e in quella dopo e' dello stesso azzurro un archivio su
%  disco, chi legge impara due volte e si fida zero. Qui un colore significa
%  una cosa sola in tutto il documento:
%
%     cielo     il MOTORE di calcolo (Python, core/)
%     menta     l'INTERFACCIA (Flask, dashboard/)
%     pesca     il SITO (Next.js) e tutto cio' che sta in rete
%     lavanda   i dati PERSISTENTI (disco, Postgres)
%     crema     cio' che viene da FUORI (broker, Stripe, provider dati)
%     bianco    strati astratti (porte, contratti, protocolli)
%
%  Questo file e' incluso sia dai documenti sia dalle figure compilate da sole
%  (figure/tikz/*.tex con la classe standalone), percio' non deve dipendere da
%  niente che esista solo dentro un documento.
% ===========================================================================
\RequirePackage{tikz}
\usetikzlibrary{
  arrows.meta, positioning, fit, backgrounds, calc, shapes.geometric,
  shapes.multipart, shapes.symbols, decorations.pathreplacing,
  decorations.pathmorphing, patterns, shadows.blur, matrix, chains, scopes
}

\tikzset{
  % --- impostazioni generali ----------------------------------------------
  every picture/.style = {
    line width = 0.7pt,
    font = \sffamily\small,
  },
  %
  % --- RIQUADRI: un modulo, un componente, una responsabilita' -------------
  riquadro/.style = {
    rectangle, rounded corners = 3pt, draw,
    minimum height = 9.5mm, minimum width = 30mm,
    inner xsep = 6pt, inner ysep = 5pt,
    align = center, font = \sffamily\small,
    text = lyraInchiostro,
  },
  motore/.style    = {riquadro, fill = lyraCielo,   draw = lyraCieloBordo},
  interfaccia/.style = {riquadro, fill = lyraMenta,  draw = lyraMentaBordo},
  sito/.style      = {riquadro, fill = lyraPesca,   draw = lyraPescaBordo},
  persistenza/.style = {riquadro, fill = lyraLavanda, draw = lyraLavandaBordo},
  esterno/.style   = {riquadro, fill = lyraCrema,   draw = lyraCremaBordo},
  astratto/.style  = {riquadro, fill = white, draw = lyraGrigio, dashed},
  %
  % Il riquadro piccolo: per gli elenchi fitti dove il nome basta.
  minuto/.style = {riquadro, minimum height = 7mm, minimum width = 22mm,
                   font = \sffamily\scriptsize, inner xsep = 4pt},
  %
  % --- FORME PARTICOLARI ---------------------------------------------------
  archivio/.style = {
    cylinder, shape border rotate = 90, aspect = 0.22, draw = lyraLavandaBordo,
    fill = lyraLavanda, minimum height = 13mm, minimum width = 22mm,
    inner xsep = 0pt, inner ysep = 4pt, align = center,
    font = \sffamily\scriptsize, text = lyraInchiostro,
  },
  decisione/.style = {
    diamond, draw = lyraInchiostro!60, fill = white, aspect = 1.9,
    inner xsep = 2pt, inner ysep = 0pt, align = center,
    font = \sffamily\scriptsize, text = lyraInchiostro,
  },
  persona/.style = {
    rectangle, rounded corners = 10pt, draw = lyraInchiostro!70,
    fill = white, minimum height = 9mm, minimum width = 24mm,
    align = center, font = \sffamily\small\itshape,
  },
  %
  % --- STRATI: il contenitore che raggruppa -------------------------------
  strato/.style = {
    rectangle, rounded corners = 5pt, draw = lyraGrigioChiaro,
    fill = lyraPannello!55!white, inner sep = 9pt,
  },
  stratoTitolo/.style = {
    font = \sffamily\scriptsize\bfseries, text = lyraInchiostro!65,
    anchor = north west,
  },
  %
  % --- FRECCE --------------------------------------------------------------
  %  Tre tipi, e la differenza va vista senza leggere la legenda:
  %    flusso      pieno    -> i dati passano di qui
  %    dipendenza  tratteg. -> questo conosce quello (ma nulla scorre)
  %    ritorno     grigio   -> la risposta che torna indietro
  flusso/.style = {
    -{Stealth[length = 2.4mm, width = 1.8mm]},
    draw = lyraInchiostro!75, line width = 0.8pt,
  },
  dipendenza/.style = {
    -{Stealth[length = 2.2mm, width = 1.6mm]}, dashed, dash pattern = on 2.4pt off 2pt,
    draw = lyraInchiostro!55, line width = 0.6pt,
  },
  ritorno/.style = {
    -{Stealth[length = 2.2mm, width = 1.6mm]},
    draw = lyraGrigio, line width = 0.6pt,
  },
  accento/.style = {
    -{Stealth[length = 2.6mm, width = 2mm]},
    draw = lyraAccento, line width = 1pt,
  },
  bloccato/.style = {
    -{Bar[width = 2.4mm]}, draw = lyraNegativo, line width = 0.9pt,
  },
  %
  % --- ETICHETTE SULLE FRECCE ---------------------------------------------
  etichetta/.style = {
    font = \sffamily\scriptsize, text = lyraInchiostro!70,
    fill = white, inner xsep = 2.5pt, inner ysep = 1pt,
  },
  etichettaAccento/.style = {etichetta, text = lyraAccentoScuro},
  %
  % --- NOTE A MARGINE ------------------------------------------------------
  nota/.style = {
    font = \sffamily\scriptsize\itshape, text = lyraInchiostro!60,
    align = left, inner sep = 2pt,
  },
  notaAccento/.style = {nota, text = lyraAccentoScuro, font = \sffamily\scriptsize},
  %
  % --- SEGNAPOSTO NUMERICI (1) (2) (3) sulle figure a passi ---------------
  passo/.style = {
    circle, draw = lyraAccento, fill = white, text = lyraAccentoScuro,
    font = \sffamily\scriptsize\bfseries, inner sep = 1.2pt,
    minimum size = 4.6mm,
  },
}

% --- la cornice con il titolo, usata dalle figure a strati -------------------
%  #1 opzioni  #2 nodi da contenere  #3 titolo
%  Il titolo sta SOPRA la cornice, non dentro: dentro finisce addosso al primo
%  nodo ogni volta che il nodo è vicino al bordo, e lo si scopre solo guardando
%  il PDF. Il nome del nodo è contato, così due strati nella stessa figura non
%  si sovrascrivono a vicenda.
\newcounter{lyrastratoN}
\newcommand{\lyrastrato}[3][]{%
  \stepcounter{lyrastratoN}%
  \begin{scope}[on background layer]
    \node[strato, fit = #2, #1] (cornice\thelyrastratoN) {};
    \node[stratoTitolo, anchor = south west]
      at ([shift={(1pt,1.5pt)}]cornice\thelyrastratoN.north west) {#3};
  \end{scope}}

% --- legenda compatta, da mettere sotto una figura ---------------------------
%  Si usa solo quando la figura introduce un colore nuovo: ripeterla sotto ogni
%  figura la fa diventare rumore che nessuno legge piu'.
\newcommand{\lyralegenda}[1]{%
  \begin{tikzpicture}[baseline]
    \matrix[matrix of nodes, nodes = {anchor = west, font = \sffamily\scriptsize},
            column sep = 3mm, row sep = 0.6mm] {#1\\};
  \end{tikzpicture}}

% ---------------------------------------------------------------------------
%  DUE TRAPPOLE, scritte qui perché costano mezz'ora ciascuna e l'errore che
%  producono indica sempre il punto sbagliato.
%
%  1. `\\` dentro un gruppo del contenuto di un nodo.
%        {testo\\{\scriptsize una riga\\e un'altra}}   <-- SI ROMPE
%        {testo\\{\scriptsize una riga}\\{\scriptsize e un'altra}}  <-- giusto
%     TikZ raccoglie le righe guardando i `\\` al primo livello: quello dentro
%     le graffe è il `\\` di LaTeX, che in un box orizzontale non ha senso.
%     L'errore che ne esce è «\tikzscope@linewidth non definito», segnalato sul
%     nodo DOPO quello colpevole.
%
%  2. Una tcolorbox dentro un nodo. Vale per \cod, \pulsante, \campo: sono
%     tcbox, cioè figure TikZ annidate, e lasciano la figura esterna in uno
%     stato incoerente. Stesso errore, stessa confusione. Nelle figure si usa
%     \codT (qui sotto).
% ---------------------------------------------------------------------------

% --- il codice DENTRO una figura --------------------------------------------
%  Nel testo un simbolo si scrive \cod{...}, che è un \tcbox: una tcolorbox,
%  cioè una tikzpicture. Infilarla dentro il contenuto di un nodo TikZ rompe lo
%  stato della figura che la contiene, e l'errore salta fuori sul nodo
%  SUCCESSIVO (`\tikzscope@linewidth` non definito) — mezz'ora buttata a
%  cercarlo nel posto sbagliato. Dentro le figure si usa questo, che è solo un
%  cambio di carattere e non apre niente.
%  \upshape serve davvero: dentro un nodo `nota`, che è in corsivo, un
%  monospazio corsivo si legge male e sembra un'altra cosa.
\newcommand{\codT}[1]{{\ttfamily\upshape\small\color{lyraInchiostro!85}#1}}

\newcommand{\lyravoce}[2]{%
  \tikz[baseline=-0.55ex]\node[#1, minimum width = 4.2mm, minimum height = 2.6mm,
        inner sep = 0pt, rounded corners = 1.4pt]{};~#2}

% ===========================================================================
%  comune/stile-codice.tex — come si mostra il codice
% ===========================================================================
%
%  PERCHE' listings E NON minted. minted colora meglio, ma per farlo chiama
%  Pygments durante la compilazione: serve Python installato e il permesso
%  --shell-escape, che molti editor LaTeX negano per difetto e che le
%  compilazioni automatiche rifiutano. Una documentazione che non compila sulla
%  macchina di chi la legge non e' documentazione. listings sta dentro LaTeX e
%  compila ovunque, e la differenza di resa non vale quel prezzo.
%
%  Tre modi di mostrare codice, e vanno usati per cose diverse:
%
%    \begin{codice}[python]{didascalia}{percorso/file.py}  blocco con cornice
%    \inputcodice{python}{didascalia}{file.tex}{percorso}  da file esterno
%    \cod{nome_simbolo}                                     dentro una frase
%
%  Il PERCORSO non e' decorazione: chi legge la documentazione ha il sorgente
%  aperto accanto, e un frammento senza il suo file costringe a cercarlo.
% ===========================================================================
\RequirePackage{listings}
\RequirePackage{tcolorbox}
\tcbuselibrary{listings, skins, breakable}

% --- stile comune ----------------------------------------------------------
\lstdefinestyle{lyraBase}{
  basicstyle      = \ttfamily\footnotesize\color{lyraInchiostro},
  keywordstyle    = \color{codParola}\bfseries,
  commentstyle    = \color{codCommento}\itshape,
  stringstyle     = \color{codStringa},
  numberstyle     = \ttfamily\scriptsize\color{codRiga},
  numbers         = left,
  numbersep       = 9pt,
  stepnumber      = 1,
  showstringspaces= false,
  breaklines      = true,
  breakatwhitespace = false,
  % La freccia dice che la riga continua: senza, una riga spezzata si legge
  % come due istruzioni, ed e' il tipo di errore che si propaga nel copia-incolla.
  prebreak        = \raisebox{0ex}[0ex][0ex]{\color{lyraAccento}\ensuremath{\hookleftarrow}},
  tabsize         = 4,
  columns         = fullflexible,
  keepspaces      = true,
  upquote         = true,
  extendedchars   = true,
  literate        =
    {à}{{\`a}}1 {è}{{\`e}}1 {é}{{\'e}}1 {ì}{{\`\i}}1 {ò}{{\`o}}1 {ù}{{\`u}}1
    {À}{{\`A}}1 {È}{{\`E}}1 {Ì}{{\`I}}1 {Ò}{{\`O}}1 {Ù}{{\`U}}1
    {€}{{\texteuro}}1 {°}{{\textdegree}}1 {·}{{\textperiodcentered}}1
    {—}{{---}}1 {–}{{--}}1 {«}{{\guillemotleft}}1 {»}{{\guillemotright}}1
    {“}{{\textquotedblleft}}1 {”}{{\textquotedblright}}1
    {’}{{\textquoteright}}1 {‘}{{\textquoteleft}}1
    {→}{{$\rightarrow$}}1 {←}{{$\leftarrow$}}1 {✓}{{$\checkmark$}}1,
}

% ---------------------------------------------------------------------------
%  REGOLA: dentro un blocco di codice si scrive SOLO ASCII.
%
%  Niente lettere accentate, niente caratteri di disegno, niente frecce.
%  `literate` non basta, e nemmeno il pacchetto listingsutf8: tcblisting
%  scrive il contenuto in un file intermedio e lo rilegge, quindi i byte
%  arrivano al kernel prima che listings possa intercettarli. L'errore e'
%  «Invalid UTF-8 byte sequence», segnalato nel file .listing e senza dire
%  quale blocco l'ha causato.
%
%  Quando un valore del sorgente contiene un accento, si translittera e si
%  mette un commento accanto. Gli alberi di cartelle si disegnano in ASCII
%  (|-- e `--), oppure in TikZ, dove non passano da listings.
%  `make controlla` segnala le violazioni prima che fermino la compilazione.
% ---------------------------------------------------------------------------

% --- i linguaggi che compaiono nel progetto ---------------------------------
\lstdefinelanguage{JavaScriptLyra}{
  keywords = {break, case, catch, class, const, continue, debugger, default,
    delete, do, else, export, extends, finally, for, function, if, import, in,
    instanceof, let, new, of, return, super, switch, this, throw, try, typeof,
    var, void, while, with, yield, async, await, static, get, set},
  keywordstyle = \color{codParola}\bfseries,
  ndkeywords = {true, false, null, undefined, document, window, console,
    fetch, JSON, Promise, Array, Object, Math},
  ndkeywordstyle = \color{lyraInchiostro!75}\bfseries,
  sensitive = true,
  comment = [l]{//},
  morecomment = [s]{/*}{*/},
  morestring = [b]',
  morestring = [b]",
  morestring = [b]`,
}

\lstdefinelanguage{TypeScriptLyra}{
  morekeywords = [1]{break, case, catch, class, const, continue, default,
    delete, do, else, enum, export, extends, finally, for, function, if,
    implements, import, in, instanceof, interface, let, new, of, private,
    protected, public, readonly, return, super, switch, this, throw, try,
    type, typeof, var, void, while, yield, async, await, as, satisfies},
  morekeywords = [2]{string, number, boolean, any, unknown, never, Promise,
    Record, Array, Partial, Required},
  keywordstyle = [1]\color{codParola}\bfseries,
  keywordstyle = [2]\color{lyraInchiostro!70},
  sensitive = true,
  comment = [l]{//},
  morecomment = [s]{/*}{*/},
  morestring = [b]',
  morestring = [b]",
  morestring = [b]`,
}

\lstdefinelanguage{Console}{
  basicstyle = \ttfamily\footnotesize,
  morecomment = [l]{\#},
  sensitive = false,
}

% --- l'ambiente principale --------------------------------------------------
%  #1 linguaggio, con il nome CORTO (python, js, ts, sql, json, yaml, html,
%     css, console): \lstlingua lo traduce nel nome che listings conosce, così
%     chi scrive non deve ricordare che JSON non esiste e che gli somiglia Java
%  #2 didascalia
%  #3 percorso del file nel progetto (vuoto = nessuno)
\newtcblisting{codiceInterno}[3]{
  listing only,
  % Sia lo stile sia la lingua devono stare QUI dentro: `listing options`
  % sostituisce `listing style` invece di aggiungersi, e il sintomo e' che
  % meta' delle impostazioni (breaklines, numeri di riga) sparisce senza
  % dirlo. Le righe lunghe escono dal riquadro, ed e' l'unico segno.
  listing options = {style = lyraBase, language = #1},
  breakable,
  enhanced,
  colback   = codFondo,
  colframe  = codBordo,
  boxrule   = 0.5pt,
  arc       = 3pt,
  left      = 3pt, right = 3pt, top = 4pt, bottom = 4pt,
  toptitle  = 2pt, bottomtitle = 2pt,
  title     = {\footnotesize\sffamily\bfseries\color{lyraInchiostro}#2%
               \ifstrempty{#3}{}{\hfill\normalfont\ttfamily\scriptsize
                 \color{lyraInchiostro!60}#3}},
  colbacktitle = lyraPannello!60!white,
  coltitle  = lyraInchiostro,
  fonttitle = \footnotesize,
  titlerule = 0pt,
  before skip = 10pt, after skip = 10pt,
}

% Nomi corti per i linguaggi, cosi' chi scrive non deve ricordare i nomi
% interni di listings ne' che JSON e YAML non esistono.
%
% Si usa \lstalias e non una macro che traduce: l'opzione `language` di
% listings viene letta senza espandere, quindi una macro li' dentro produce
% «Missing \endcsname» alla fine del blocco di codice — errore che non dice
% niente di cio' che e' successo.
\lstalias{console}{Console}
\lstalias{python}{Python}
\lstalias{js}{JavaScriptLyra}
\lstalias{ts}{TypeScriptLyra}
\lstalias{sql}{SQL}
\lstalias{json}{Java}      % JSON non esiste in listings: Java gli somiglia
\lstalias{yaml}{Python}    % idem per YAML: commenti con #, due punti, indentazione
\lstalias{html}{HTML}
\lstalias{css}{HTML}
\lstalias{testo}{}

% --- codice da file esterno -------------------------------------------------
%  #1 linguaggio  #2 didascalia  #3 percorso del file su disco  #4 percorso da
%  mostrare. Si usa per gli estratti lunghi tenuti fuori dal testo.
\newcommand{\inputcodice}[4]{%
  \begin{tcolorbox}[
    breakable, enhanced,
    colback = codFondo, colframe = codBordo, boxrule = 0.5pt, arc = 3pt,
    left = 3pt, right = 3pt, top = 4pt, bottom = 4pt,
    title = {\footnotesize\sffamily\bfseries\color{lyraInchiostro}#2%
             \hfill\normalfont\ttfamily\scriptsize\color{lyraInchiostro!60}#4},
    colbacktitle = lyraPannello!60!white, coltitle = lyraInchiostro,
    fonttitle = \footnotesize, titlerule = 0pt]
    \lstinputlisting[style=lyraBase, language=#1]{#3}
  \end{tcolorbox}}

% --- codice dentro una frase ------------------------------------------------
%  Fondo tenue invece del grassetto: dentro un paragrafo il grassetto chiama
%  l'occhio e spezza la lettura, mentre un nome di funzione va riconosciuto
%  senza essere gridato.
%  Le due penalita' attorno alla scatola non sono decorazione. Un \tcbox e'
%  una scatola indivisibile: se capita dopo una parentesi o un trattino --
%  cioe' senza spazio davanti -- TeX non ha dove spezzare la riga e la lascia
%  sbordare nel margine. Sono state la causa di quasi tutte le righe troppo
%  lunghe della prima stesura. Con queste, la scatola puo' sempre andare a
%  capo per conto suo.
%
%  Resta un limite: una scatola piu' larga della giustezza sborda comunque.
%  Per un percorso lungo si usa \file{} (testo normale, si spezza) oppure lo
%  si mette in un blocco di codice.
\DeclareRobustCommand{\cod}[1]{%
  \penalty50
  \tcbox[on line, boxsep = 0.6pt, left = 1.4pt, right = 1.4pt, top = 0pt,
         bottom = 0pt, arc = 2pt, boxrule = 0pt,
         colback = lyraPannello, colframe = lyraPannello]
        {\ttfamily\small #1}%
  \penalty0 }

% Percorso di file nel testo.
%  Usa lo stesso ciclo di \percorsoHTTP (definito in comune/macro.tex): un nome
%  come `diagnose_regime_signals.py' e' lungo quattro centimetri e senza punti
%  di sillabazione, quindi in una colonna stretta sborda. Con i punti di
%  interruzione dopo / _ . e - va a capo dove si legge comunque.
%  Il ciclo richiede ASCII e niente spazi: per il resto c'e' \cod{}.
\DeclareRobustCommand{\file}[1]{{\color{lyraInchiostro!78}\percorsoHTTP{#1}}}

% ===========================================================================
%  Un COMMENTO del sorgente citato per esteso
% ===========================================================================
%  Diversa da `codiceInterno`, e la differenza non e' estetica. Molti commenti
%  di questo progetto sono prose italiane di dieci righe che spiegano una
%  decisione: sono la parte piu' utile da citare, e dentro un blocco di
%  listings non ci stanno — listings vuole ASCII puro (vedi sopra), e
%  translitterare «perche'» e le lettere accentate significa citare qualcosa
%  che nel sorgente non c'e'.
%
%  Qui il testo e' testo LaTeX normale: accenti, virgolette e trattini lunghi
%  si compongono come nel resto del documento, e la citazione resta fedele.
%
%    \begin{commento}{percorso/del/file.py:120}
%    Il testo del commento, a capo come nel sorgente.
%    \end{commento}
%
%  Il corsivo e il filetto a sinistra dicono «questo l'ha scritto chi ha
%  scritto il codice, non chi ha scritto la guida».
\newtcolorbox{commento}[1]{
  enhanced, breakable,
  colback   = codFondo, colframe = lyraInchiostro!30,
  boxrule   = 0pt, leftrule = 2.2pt, arc = 2pt,
  left = 9pt, right = 9pt, top = 6pt, bottom = 6pt,
  fontupper = \small\itshape,
  title     = {\normalfont\ttfamily\scriptsize\color{lyraInchiostro!60}#1},
  colbacktitle = codFondo, coltitle = lyraInchiostro,
  titlerule = 0pt, toptitle = 2pt, bottomtitle = 0pt,
  before skip = 10pt, after skip = 10pt,
}


% --- numeri -----------------------------------------------------------------
%  In italiano il separatore delle migliaia sarebbe il punto, ma in un
%  documento pieno di valori numerici (soglie, quantili, prezzi) il punto
%  decimale inglese del codice e il punto delle migliaia del testo si
%  confondono. Spazio fine: non si confonde con niente.
\usepackage{siunitx}
\sisetup{group-separator = {\,}, group-minimum-digits = 4, detect-all}

% --- tabelle ----------------------------------------------------------------
\usepackage{booktabs}
\usepackage{longtable}
\usepackage{multirow}
\usepackage{array}
\usepackage{tabularx}
\usepackage{colortbl}

\newcolumntype{L}[1]{>{\raggedright\arraybackslash}p{#1}}
\newcolumntype{C}[1]{>{\centering\arraybackslash}p{#1}}
\newcolumntype{R}[1]{>{\raggedleft\arraybackslash}p{#1}}
\newcolumntype{M}{>{\ttfamily\footnotesize}l}

% Righe alternate tenui: su una tabella di venti righe l'occhio perde la riga,
% e un colore appena percepibile costa meno di una griglia piena.
\newcommand{\rigaTenue}{\rowcolor{lyraPannello!55!white}}

% --- elenchi ----------------------------------------------------------------
\usepackage{enumitem}
\setlist{itemsep = 1pt, parsep = 2pt, topsep = 4pt}
\setlist[itemize]{label = \textcolor{lyraAccento}{\textbullet}}
\setlist[enumerate]{label = \textcolor{lyraAccentoScuro}{\arabic*.}}

% --- figure -----------------------------------------------------------------
\usepackage{graphicx}
\usepackage{float}
\usepackage[font = small, labelfont = {bf, sf}, textfont = sf,
            labelsep = period, skip = 7pt]{caption}
\usepackage{subcaption}
\usepackage{adjustbox}
\graphicspath{{figure/}{figure/tikz/}{figure/schermate/}}

% --- intestazioni e piè di pagina -------------------------------------------
\usepackage{fancyhdr}
\usepackage{titlesec}
\usepackage{titletoc}

% --- riferimenti ------------------------------------------------------------
\usepackage[
  colorlinks = true,
  linkcolor  = lyraAccentoScuro,
  citecolor  = lyraAccentoScuro,
  urlcolor   = lyraAccentoScuro,
  filecolor  = lyraAccentoScuro,
  breaklinks = true,
  bookmarksnumbered = true,
  bookmarksopen = true,
  bookmarksopenlevel = 1,
]{hyperref}
\usepackage{url}
\urlstyle{sf}

% --- indice analitico -------------------------------------------------------
\usepackage{imakeidx}

% ===========================================================================
%  Stile dei titoli
% ===========================================================================
%  I capitoli si aprono con il numero grande a sinistra e un filetto verde
%  acqua: è il segno che si gira pagina in una sezione nuova, e lo si vede
%  sfogliando senza leggere.
\titleformat{\chapter}[display]
  {\sffamily\LARGE\bfseries\color{lyraInchiostro}}
  {\raggedright\normalfont\sffamily\fontsize{52}{52}\selectfont
   \color{lyraAccento!35!lyraPannello}\thechapter}
  {-6pt}
  {\raggedright}
  [\vspace{4pt}{\color{lyraAccento}\rule{\textwidth}{1.2pt}}]
\titlespacing*{\chapter}{0pt}{-24pt}{26pt}

\titleformat{\section}
  {\sffamily\large\bfseries\color{lyraInchiostro}}
  {\color{lyraAccentoScuro}\thesection}{0.7em}{}
\titlespacing*{\section}{0pt}{18pt}{7pt}

\titleformat{\subsection}
  {\sffamily\normalsize\bfseries\color{lyraInchiostro!88}}
  {\color{lyraAccentoScuro}\thesubsection}{0.6em}{}
\titlespacing*{\subsection}{0pt}{13pt}{4pt}

\titleformat{\subsubsection}[runin]
  {\sffamily\small\bfseries\color{lyraInchiostro!80}}
  {}{0pt}{}[\;---\;]
\titlespacing*{\subsubsection}{0pt}{9pt}{0.5em}

\setcounter{secnumdepth}{3}
\setcounter{tocdepth}{2}

% ===========================================================================
%  Indice
% ===========================================================================
%  L'indice è composto con titletoc, non con tocloft: i due pacchetti si
%  contendono le stesse macro e caricarli insieme rompe l'indice in modi che
%  si vedono solo a documento finito.
\titlecontents{chapter}[0pt]
  {\addvspace{9pt}\sffamily\bfseries}
  {\color{lyraAccentoScuro}\contentslabel{1.6em}}
  {}
  {\color{lyraInchiostro!55}\titlerule*[0.6pc]{.}\contentspage}
\titlecontents{section}[1.9em]
  {\sffamily\small}
  {\color{lyraInchiostro!70}\contentslabel{2.3em}}
  {}
  {\color{lyraInchiostro!45}\titlerule*[0.6pc]{.}\contentspage}

% ===========================================================================
%  I riquadri del testo
% ===========================================================================
%  Quattro soli, e ciascuno ha un lavoro preciso. Un documento con nove tipi di
%  riquadro è un documento in cui nessun riquadro significa qualcosa.
%
%    \begin{nota}{titolo}        un chiarimento utile, saltabile
%    \begin{attenzione}{titolo}  qui si rompe qualcosa
%    \begin{perche}{titolo}      perché è così e non altrimenti
%    \begin{pratica}{titolo}     cosa fare, passo per passo
% ===========================================================================
\newtcolorbox{nota}[1]{
  enhanced, breakable,
  colback = lyraCielo!35!white, colframe = lyraCieloBordo,
  boxrule = 0pt, leftrule = 2.2pt, arc = 2pt,
  left = 8pt, right = 8pt, top = 6pt, bottom = 6pt,
  fonttitle = \sffamily\bfseries\small,
  coltitle = lyraInchiostro,
  title = {#1}, titlerule = 0pt, colbacktitle = lyraCielo!35!white,
  toptitle = 1pt, bottomtitle = 1pt,
  before skip = 10pt, after skip = 10pt,
}

\newtcolorbox{attenzione}[1]{
  enhanced, breakable,
  colback = lyraPesca!40!white, colframe = lyraNegativo!70,
  boxrule = 0pt, leftrule = 2.2pt, arc = 2pt,
  left = 8pt, right = 8pt, top = 6pt, bottom = 6pt,
  fonttitle = \sffamily\bfseries\small,
  coltitle = lyraNegativo,
  title = {#1}, titlerule = 0pt, colbacktitle = lyraPesca!40!white,
  toptitle = 1pt, bottomtitle = 1pt,
  before skip = 10pt, after skip = 10pt,
}

\newtcolorbox{perche}[1]{
  enhanced, breakable,
  colback = lyraPannello, colframe = lyraGrigioChiaro,
  boxrule = 0pt, leftrule = 2.2pt, arc = 2pt,
  left = 8pt, right = 8pt, top = 6pt, bottom = 6pt,
  fonttitle = \sffamily\bfseries\small,
  coltitle = lyraInchiostro!75,
  title = {#1}, titlerule = 0pt, colbacktitle = lyraPannello,
  toptitle = 1pt, bottomtitle = 1pt,
  before skip = 10pt, after skip = 10pt,
}

\newtcolorbox{pratica}[1]{
  enhanced, breakable,
  colback = lyraMenta!45!white, colframe = lyraAccento,
  boxrule = 0pt, leftrule = 2.2pt, arc = 2pt,
  left = 8pt, right = 8pt, top = 6pt, bottom = 6pt,
  fonttitle = \sffamily\bfseries\small,
  coltitle = lyraAccentoScuro,
  title = {#1}, titlerule = 0pt, colbacktitle = lyraMenta!45!white,
  toptitle = 1pt, bottomtitle = 1pt,
  before skip = 10pt, after skip = 10pt,
}

% ===========================================================================
%  Macro del progetto
% ===========================================================================
% ===========================================================================
%  comune/macro.tex — le abbreviazioni del progetto
% ===========================================================================
%
%  Servono a scrivere una cosa in un posto solo. Il nome di un pacchetto
%  Python compare duecento volte nella guida: se un giorno `approaches/`
%  diventa `overlays/`, si cambia qui e non duecento volte a mano — e
%  soprattutto non si finisce con metà documento che dice una cosa e metà
%  un'altra, che è il modo in cui la documentazione smette di essere creduta.
% ===========================================================================

% --- la lingua del documento -----------------------------------------------
%  \iflyraInglese la definisce comune/preambolo.tex. Qui la si crea se manca,
%  perché `make figure' include questo file senza il preambolo. Il nome del
%  condizionale si scrive con \csname: scritto per esteso dentro un ramo
%  saltato, TeX lo conterebbe come un \if aperto e sbaglierebbe i \fi.
\ifcsname iflyraInglese\endcsname\else
  \expandafter\newif\csname iflyraInglese\endcsname
\fi

% --- nomi del prodotto ------------------------------------------------------
\newcommand{\Lyra}{\textsc{Lyra}}
\newcommand{\LyraCapital}{\textsc{Lyra Capital}}
\newcommand{\app}{\iflyraInglese the application\else l'applicazione\fi}
\newcommand{\App}{\iflyraInglese The application\else L'applicazione\fi}
\newcommand{\sito}{\iflyraInglese the site\else il sito\fi}
\newcommand{\Sito}{\iflyraInglese The site\else Il sito\fi}

% --- i pacchetti del motore -------------------------------------------------
%  Si scrivono \pkgCore e non \file{core/}: così il giorno in cui una cartella
%  cambia nome, cambia qui.
\DeclareRobustCommand{\pkg}[1]{{\ttfamily\small\color{lyraInchiostro!78}#1}}
\newcommand{\pkgCore}{\pkg{core/}}
\newcommand{\pkgStrategie}{\pkg{strategies/}}
\newcommand{\pkgApprocci}{\pkg{approaches/}}
\newcommand{\pkgPesi}{\pkg{weights/}}
\newcommand{\pkgOverlay}{\pkg{overlays/}}
\newcommand{\pkgRegime}{\pkg{regime/}}
\newcommand{\pkgDati}{\pkg{data/}}
\newcommand{\pkgBacktest}{\pkg{backtest/}}
\newcommand{\pkgLive}{\pkg{live/}}
\newcommand{\pkgEsecuzione}{\pkg{execution/}}
\newcommand{\pkgReport}{\pkg{reporting/}}
\newcommand{\pkgDashboard}{\pkg{dashboard/}}
\newcommand{\pkgDominio}{\pkg{domain/}}
\newcommand{\pkgPorte}{\pkg{ports/}}
\newcommand{\pkgApplicazione}{\pkg{application/}}
\newcommand{\pkgAdattatori}{\pkg{adapters/}}
\newcommand{\pkgComposizione}{\pkg{composition/}}
\newcommand{\pkgApi}{\pkg{api/}}

% --- termini del dominio, con la prima occorrenza indicizzata ---------------
\newcommand{\term}[1]{\emph{#1}\index{#1}}
\newcommand{\termdef}[1]{\textbf{#1}\index{#1}}

% --- riferimenti a classi, funzioni, file -----------------------------------
%  \simbolo{Nome}            una classe o una funzione
%  \sorgente{file.py}{120}   un punto preciso del sorgente
\DeclareRobustCommand{\simbolo}[1]{\cod{#1}}
\newcommand{\sorgente}[2]{{\ttfamily\footnotesize\color{lyraInchiostro!60}#1:#2}}

% --- tasti e voci dell'interfaccia, per il manuale utente -------------------
%  \voce{Overview}           una voce di menu o una scheda
%  \pulsante{Run backtest}   un pulsante da premere
%  \campo{Target vol}        un campo da compilare
%  \percorso{Settings > Data} un percorso di navigazione
\newcommand{\voce}[1]{\textsf{\textbf{#1}}}
\newcommand{\pulsante}[1]{%
  \tcbox[on line, boxsep = 1pt, left = 3pt, right = 3pt, top = 0.5pt,
         bottom = 0.5pt, arc = 3pt, boxrule = 0.5pt,
         colback = lyraInchiostro, colframe = lyraInchiostro,
         coltext = white]{\sffamily\scriptsize\bfseries #1}}
\newcommand{\pulsantesec}[1]{%
  \tcbox[on line, boxsep = 1pt, left = 3pt, right = 3pt, top = 0.5pt,
         bottom = 0.5pt, arc = 3pt, boxrule = 0.5pt,
         colback = white, colframe = lyraInchiostro!50,
         coltext = lyraInchiostro]{\sffamily\scriptsize\bfseries #1}}
\newcommand{\campo}[1]{\textsf{\itshape #1}}
\newcommand{\percorso}[1]{\textsf{\small #1}}
\newcommand{\freccia}{\;\textcolor{lyraAccento}{$\rightarrow$}\;}

% --- una rotta HTTP ---------------------------------------------------------
%  \rotta{GET}{/api/marketplace/<id>/perche-non-mio}
%
%  Esiste perche' \cod{} e' una scatola indivisibile: in una colonna di
%  tabella larga cinque centimetri un percorso lungo sborda, e non c'e'
%  niente che TeX possa fare. Qui il percorso e' testo normale e puo' andare
%  a capo dopo ogni barra, che e' il punto naturale.
%
%  Il ciclo prende un carattere alla volta: funziona perche' un percorso non
%  contiene spazi (un argomento non delimitato li mangerebbe) ed e' in ASCII
%  (una lettera accentata e' un carattere attivo e manda in errore il ciclo).
%  Per qualunque altra cosa si usa \cod{}.
\makeatletter
\def\lyra@finePercorso{\lyra@finePercorsoMARK}
\def\lyra@percorso#1{%
  \ifx#1\lyra@finePercorso
    \let\lyra@prossimo\relax
  \else
    #1%
    % Dopo questi quattro caratteri la riga puo' andare a capo. Sono i punti
    % in cui un percorso o un nome di file si leggono comunque bene spezzati.
    \ifx#1/\allowbreak\else
      \ifx#1\_\allowbreak\else
        \ifx#1.\allowbreak\else
          \ifx#1-\allowbreak\fi\fi\fi\fi
    \let\lyra@prossimo\lyra@percorso
  \fi
  \lyra@prossimo}
\DeclareRobustCommand{\percorsoHTTP}[1]{%
  {\ttfamily\small\color{lyraInchiostro!85}%
   \lyra@percorso#1\lyra@finePercorso}}
\makeatother

\DeclareRobustCommand{\rotta}[2]{%
  {\sffamily\scriptsize\bfseries\color{lyraAccentoScuro}#1}\,\percorsoHTTP{#2}}

% --- la figura standard -----------------------------------------------------
%  Centra, mette la didascalia sotto e l'etichetta. Esiste perché le tre righe
%  scritte a mano ogni volta divergono: una figura senza \label non si può
%  citare, e una citata per numero si rompe appena se ne aggiunge un'altra.
%  #1 file da includere (senza estensione)  #2 didascalia  #3 etichetta
%  `max width` e non `width`: una figura piccola resta della sua misura, una
%  piu' larga della giustezza viene rimpicciolita invece di sbordare nel
%  margine. Senza, ogni figura andrebbe dimensionata a mano e ridimensionata
%  a ogni modifica del testo.
\newcommand{\figuraTikz}[3]{%
  \begin{figure}[htbp]
    \centering
    \begin{adjustbox}{max width = \linewidth}
      \input{figure/tikz/#1}
    \end{adjustbox}
    \caption{#2}\label{fig:#3}
  \end{figure}}

%  Come sopra ma per un'immagine, con larghezza.
%  #1 larghezza  #2 file  #3 didascalia  #4 etichetta
\newcommand{\figuraImg}[4]{%
  \begin{figure}[htbp]
    \centering
    \includegraphics[width = #1\linewidth]{#2}
    \caption{#3}\label{fig:#4}
  \end{figure}}

%  Schermata con cornice: una schermata senza bordo si confonde con il fondo
%  della pagina, e non si capisce dove finisce l'interfaccia. La cornice sta
%  DENTRO la larghezza chiesta: altrimenti una schermata larga quanto la
%  pagina sborda di un punto per lato.
\newcommand{\figuraSchermata}[4]{%
  \begin{figure}[htbp]
    \centering
    \setlength{\fboxsep}{0pt}\setlength{\fboxrule}{0.5pt}%
    {\color{lyraGrigioChiaro}\fbox{\includegraphics[width = \dimexpr #1\linewidth - 2\fboxrule\relax]{#2}}}
    \caption{#3}\label{fig:#4}
  \end{figure}}

% --- riferimenti incrociati con la parola davanti ---------------------------
%  In inglese la parola vuole la maiuscola («see Chapter 4»), in italiano no.
\newcommand{\vediFig}[1]{\iflyraInglese Figure\else figura\fi~\ref{fig:#1}}
\newcommand{\vediTab}[1]{\iflyraInglese Table\else tabella\fi~\ref{tab:#1}}
\newcommand{\vediCap}[1]{\iflyraInglese Chapter\else capitolo\fi~\ref{cap:#1}}
\newcommand{\vediSez}[1]{\iflyraInglese Section~\else \S\,\fi\ref{sez:#1}}

% --- numeri e unità ---------------------------------------------------------
%  Ogni lingua con la sua convenzione, la stessa dell'applicazione e del sito:
%  in italiano «10\,%» e «49\,€», in inglese «10%» e «€49» (le migliaia, in
%  inglese, si scrivono con la virgola nell'argomento: \eur{1,490}).
\newcommand{\perc}[1]{\iflyraInglese #1\%\else #1\,\%\fi}
\newcommand{\eur}[1]{\iflyraInglese \texteuro#1\else #1\,\texteuro\fi}


% ===========================================================================
%  Le macro dentro i segnalibri del PDF
% ===========================================================================
%  hyperref deve trasformare il titolo di ogni capitolo e sezione in una
%  stringa semplice per il pannello dei segnalibri. Le macro che compongono
%  qualcosa — una scatola colorata, un ciclo sui caratteri — lo mandano in
%  tilt: l'errore e' «TeX capacity exceeded [input stack size]» su una riga
%  che non ha niente di strano, tipo
%      \subsection{Che cosa c'e' sotto \file{data/}}
%
%  Qui si dice a hyperref di considerarle, nei soli segnalibri, come
%  «stampa il contenuto e basta».
\pdfstringdefDisableCommands{%
  \makeatletter
  \let\file\@firstofone
  \let\cod\@firstofone
  \let\pkg\@firstofone
  \let\simbolo\@firstofone
  \let\percorsoHTTP\@firstofone
  \let\voce\@firstofone
  \let\campo\@firstofone
  \let\pulsante\@firstofone
  \let\pulsantesec\@firstofone
  \let\percorso\@firstofone
  \let\term\@firstofone
  \let\termdef\@firstofone
  \def\rotta#1#2{#1 #2}%
  \def\freccia{ > }%
  \def\Lyra{Lyra}%
  \def\LyraCapital{Lyra Capital}%
  \makeatother
}

% ===========================================================================
%  Controllo: la sillabazione italiana c'è davvero?
% ===========================================================================
%  I pattern di sillabazione si caricano quando si GENERA il formato, non
%  quando si compila: se mancano, babel non fallisce — ripiega sull'inglese, e
%  il documento esce con le parole spezzate nei punti sbagliati. È un difetto
%  che si vede solo guardando la pagina, quindi vale la pena dirlo nel log.
%
%  Per rimediare, su una TeX Live completa: `tlmgr install hyphen-italian`;
%  su Debian/Ubuntu: `apt install texlive-lang-italian`; poi `fmtutil-sys
%  --byfmt pdflatex`.
%
%  Il controllo guarda \language, che è un parametro intero del motore e si
%  può sempre confrontare: 0 è la lingua di partenza (inglese). Se dopo aver
%  selezionato l'italiano siamo ancora lì, i pattern non ci sono.
\AtBeginDocument{%
  \ifnum\language = 0
    \PackageWarningNoLine{lyra}{%
      Pattern di sillabazione ITALIANI assenti: il testo sara' spezzato\MessageBreak
      con le regole inglesi. Installare hyphen-italian e rigenerare il\MessageBreak
      formato (fmtutil-sys --byfmt pdflatex)}%
  \fi}


\makeindex[intoc = true, title = Indice analitico]

\hypersetup{
  pdftitle  = {Lyra — Guida dello sviluppatore},
  pdfauthor = {Lyra Capital},
  pdfsubject = {Architettura, motore quantitativo, estensioni e rilascio},
  pdfkeywords = {Lyra, quant, trading, architettura, Python, Flask, Next.js},
}

% --- intestazioni -----------------------------------------------------------
\pagestyle{fancy}
\fancyhf{}
\renewcommand{\headrulewidth}{0pt}
\renewcommand{\chaptermark}[1]{\markboth{\thechapter.\ #1}{}}
\renewcommand{\sectionmark}[1]{\markright{\thesection\ #1}}
\fancyhead[LE]{\sffamily\scriptsize\color{lyraInchiostro!55}\leftmark}
\fancyhead[RO]{\sffamily\scriptsize\color{lyraInchiostro!55}\rightmark}
\fancyfoot[LE,RO]{\sffamily\small\color{lyraInchiostro!70}\thepage}
\fancyfoot[RE,LO]{\sffamily\scriptsize\color{lyraInchiostro!40}Lyra · Guida dello sviluppatore}

\fancypagestyle{plain}{%
  \fancyhf{}%
  \renewcommand{\headrulewidth}{0pt}%
  \fancyfoot[LE,RO]{\sffamily\small\color{lyraInchiostro!70}\thepage}%
  \fancyfoot[RE,LO]{\sffamily\scriptsize\color{lyraInchiostro!40}Lyra · Guida dello sviluppatore}}

\begin{document}

\frontmatter
% ===========================================================================
%  comune/frontespizio-sviluppatore.tex
% ===========================================================================
\begin{titlepage}
\thispagestyle{empty}
\begin{tikzpicture}[remember picture, overlay]

  % La banda verde acqua in alto: è il segno del prodotto, e serve a
  % riconoscere il documento di lato in una pila di fogli.
  \fill[lyraAccento] (current page.north west)
    rectangle ([yshift = -9mm]current page.north east);

  \node[anchor = north west, inner sep = 0pt]
    at ([shift = {(28mm, -42mm)}]current page.north west)
    {\begin{minipage}{125mm}
       {\sffamily\fontsize{10}{12}\selectfont
        \color{lyraAccentoScuro}\bfseries
        LYRA CAPITAL\par}
       \vspace{3mm}
       {\sffamily\fontsize{40}{44}\selectfont\color{lyraInchiostro}
        Guida dello\\sviluppatore\par}
       \vspace{5mm}
       {\color{lyraAccento}\rule{34mm}{1.6pt}\par}
       \vspace{6mm}
       {\sffamily\fontsize{12}{17}\selectfont\color{lyraInchiostro!62}
        Architettura, motore quantitativo,\\
        estensioni e rilascio della piattaforma.\par}
     \end{minipage}};

  % Il riquadro in basso dice a chi serve il documento. Sta in copertina
  % perché chi lo apre per sbaglio deve accorgersene subito: il manuale utente
  % (la User guide, in inglese) è un altro file, e cercare qui come si preme
  % un pulsante è tempo perso.
  \node[anchor = south west, inner sep = 0pt]
    at ([shift = {(28mm, 34mm)}]current page.south west)
    {\begin{minipage}{125mm}
       \begin{tcolorbox}[enhanced, colback = lyraPannello,
                         colframe = lyraBordo, boxrule = 0.6pt, arc = 4pt,
                         left = 10pt, right = 10pt, top = 9pt, bottom = 9pt]
         {\sffamily\footnotesize\color{lyraInchiostro}%
          \textbf{A chi serve.} A chi deve \emph{modificare} Lyra: aggiungere
          una strategia, cambiare il motore, toccare l'interfaccia, mettere
          mano al marketplace.\par
          \vspace{4pt}
          \textbf{A chi non serve.} A chi deve \emph{usare} Lyra: quello è la
          \emph{User guide}, il manuale utente in inglese, documento separato
          nella stessa cartella.\par}
       \end{tcolorbox}
     \end{minipage}};

  \node[anchor = south east, inner sep = 0pt, align = right]
    at ([shift = {(-24mm, 18mm)}]current page.south east)
    {\sffamily\scriptsize\color{lyraInchiostro!45}
     Documento generato da sorgente \LaTeX\\
     \today};

\end{tikzpicture}
\end{titlepage}

\cleardoublepage

% --- la pagina che dice come è fatto il documento ---------------------------
\thispagestyle{empty}
\vspace*{18mm}

{\sffamily\Large\bfseries\color{lyraInchiostro}Come è fatto questo documento\par}
\vspace{3mm}
{\noindent\color{lyraAccento}\rule{\textwidth}{1pt}\par}
\vspace{7mm}

\noindent
Quattro parti, in ordine di quanto a fondo si deve andare.

\begin{description}[leftmargin = 2.4em, style = nextline, itemsep = 5pt]
  \item[Parte I — Orientarsi]
    Che cos'è Lyra, come è diviso, e come si mette in moto una copia sulla
    propria macchina. Si legge una volta.
  \item[Parte II — Come è fatto]
    L'architettura e il motore di calcolo: dove stanno le decisioni e perché
    sono dove sono. È la parte da leggere prima di toccare qualunque cosa.
  \item[Parte III — Come si cambia]
    Aggiungere una strategia, una modalità di pesatura, una pagina,
    un'operazione sul database. Ogni capitolo è una procedura che si segue
    con il codice aperto accanto.
  \item[Parte IV — Tenerlo in piedi]
    Le prove automatiche, il rilascio, e l'elenco di ciò che si rompe più
    spesso con il motivo accanto.
\end{description}

\vspace{6mm}
{\sffamily\large\bfseries\color{lyraInchiostro}Le convenzioni\par}
\vspace{3mm}

\noindent
I riquadri sono quattro e hanno quattro significati distinti.

\vspace{3mm}
\begin{nota}{Nota}
Un chiarimento utile ma non necessario: saltandolo non ci si perde niente di
ciò che serve per lavorare.
\end{nota}

\begin{perche}{Perché così}
La ragione di una scelta che, vista da fuori, sembrerebbe potersi fare
diversamente. Serve a non «semplificare» qualcosa che è complicato per un
motivo.
\end{perche}

\begin{pratica}{In pratica}
Cosa fare, passo per passo. Se si cerca solo l'operazione, sono questi i
riquadri da seguire.
\end{pratica}

\begin{attenzione}{Attenzione}
Qui si rompe qualcosa: dati persi, denaro movimentato, o un errore che si
scopre solo settimane dopo. Non sono avvisi di cortesia.
\end{attenzione}

\vspace{4mm}
\noindent
Nel testo, \cod{questo} è un simbolo del codice, \file{questo/} un percorso nel
progetto, \voce{Questa} una voce dell'interfaccia e \pulsante{Questo} un
pulsante da premere. Ogni blocco di codice porta in alto a destra il file da
cui viene: la documentazione si legge con il sorgente aperto accanto, e un
frammento senza il suo file costringe a cercarlo.

\cleardoublepage


\tableofcontents

\mainmatter

% --- PARTE I — orientarsi ---------------------------------------------------
\part{Orientarsi}
% ===========================================================================
%  sviluppatore/01-introduzione/00-capitolo.tex
% ===========================================================================
%  SOLO il titolo e l'elenco delle sezioni. Il testo sta nei file 01-, 02-, ...
% ===========================================================================
\chapter{Che cos'è Lyra}\label{cap:introduzione}

% ===========================================================================
%  sviluppatore/01-introduzione/01-cosa-e.tex
% ===========================================================================
\section{Il problema che risolve}\label{sez:cosa-e}

\Lyra{} è una piattaforma per costruire strategie di investimento sistematiche,
provarle sui dati del passato e farle girare su un conto reale. Le tre cose
sono una sola: il motivo per cui il sistema è grande è che il \term{backtest} e
l'operatività \term{live} devono produrre la stessa decisione a partire dagli
stessi dati, e tenere allineate due strade che calcolano la stessa cosa costa
più codice che scriverne una.

Tre numeri, per dare la scala prima dei dettagli:

\begin{center}
\begin{tabular}{@{}L{44mm}R{20mm}L{62mm}@{}}
\toprule
\textbf{Misura} & \textbf{Valore} & \textbf{Dove sta} \\
\midrule
\rigaTenue
File, in tutto        & \num{806}      & dei quali \num{544} Python \\
Righe Python          & \num{119084}   & \num{28494} nella sola \pkgDashboard \\
\rigaTenue
Righe di sorgente     & \num{175225}   & Python, JavaScript, CSS, HTML,
                                         YAML, Markdown \\
Rotte HTTP            & \num{249}      & \pkgDashboard, undici moduli \\
\bottomrule
\end{tabular}
\end{center}

\noindent
Circa il \num{40}\,\% del sorgente è interfaccia: \num{28494} righe di Python
più \num{42072} fra template e file statici. Non è uno squilibrio. Il valore
del prodotto sta nel poter costruire una strategia senza scrivere Python, e
quella possibilità costa pagine, editor, validatori e un sistema a plugin.

\subsection*{Che cosa fa, in ordine di uso}

\begin{enumerate}
  \item \textbf{Scarica i dati.} Prezzi, fondamentali, tassi, notizie. Finiscono
        in file \file{parquet} su disco, in una gerarchia fissa.
  \item \textbf{Costruisce una strategia.} Scrivendo Python, disegnando un
        grafo di nodi, o descrivendola a parole a un modello locale.
  \item \textbf{La prova.} Un \term{walk-forward} che riaddestra e rivaluta a
        finestre successive, mai sull'intero campione in una volta.
  \item \textbf{Combina più strategie.} Ogni strategia produce uno \term{stream}
        di profitti e perdite; un \term{ensemble} li pesa, e come pesarli è esso
        stesso un plugin --- quattordici modalità, dall'equiponderato al
        meta-allocatore addestrato.
  \item \textbf{Riduce l'esposizione quando serve.} Gli \term{overlay} tagliano
        l'esposizione in regime di stress, e sono anch'essi plugin: sette.
  \item \textbf{Esegue.} Lo stesso ensemble, lo stesso codice di decisione,
        collegato a un broker invece che ai dati storici.
  \item \textbf{Pubblica.} Una strategia può essere messa nel marketplace:
        qualcun altro la compra e se la installa nel proprio sistema.
\end{enumerate}

\begin{perche}{Perché il marketplace cambia l'architettura}
Finché una strategia è un file sul proprio computer, può essere qualunque
cosa. Dal momento in cui può essere \emph{venduta e installata altrove},
diventa un artefatto con un contratto: deve dichiarare che cosa gli serve,
deve poter essere verificata senza essere eseguita, e deve poter essere
confrontata con ciò che il venditore aveva promesso. Buona parte dei vincoli
che nel codice sembrano eccessivi --- i descrittori obbligatori, i validatori
\simbolo{AST}, i \cod{META} fatti di soli letterali --- esiste per questo, e
non per rigore astratto.
\end{perche}

\subsection*{Che cosa non fa}

Non è un broker, non custodisce denaro e non esegue ordini per conto di
nessuno: si collega al conto che l'utente ha già. Non è un fornitore di dati:
usa le chiavi API dell'utente verso Polygon, FRED e simili. Non dà consigli
d'investimento, e nessun numero prodotto dal sistema è una previsione --- un
backtest dice come una strategia si \emph{sarebbe} comportata, con i costi
stimati, sui dati che esistevano.

% ===========================================================================
%  sviluppatore/01-introduzione/02-mappa.tex
% ===========================================================================
\section{Le due metà}\label{sez:due-meta}

La prima cosa da sapere di \Lyra{} è che non è un'applicazione: sono due, che
girano in posti diversi e si parlano attraverso tre soli canali.

\figuraTikz{mappa-sistema}
  {Le due metà del sistema. Tutto ciò che non attraversa uno dei tre canali
   numerati resta dalla parte in cui è nato.}
  {mappa-sistema}

\subsection{La metà che gira sul computer dell'utente}

È il progetto Python: \pkgCore, \pkgStrategie, \pkgBacktest, \pkgLive,
\pkgReport, e l'interfaccia Flask in \pkgDashboard. Gira in locale,
tipicamente su \url{http://127.0.0.1:5000}, e si impacchetta come applicazione
desktop (\file{desktop\_app.py} apre la stessa dashboard in una finestra
nativa invece che in una scheda del browser).

Qui stanno i dati di mercato e i risultati dei backtest, e qui restano. Non
c'è nessun percorso per cui un prezzo scaricato o un'equity line calcolata
finisca in rete.

\subsection{La metà che gira in rete}

È il sito, un'applicazione Next.js con un database Postgres. Tiene gli
account, le password, gli abbonamenti, i pagamenti e il catalogo del
marketplace. Non ha accesso ai dati di mercato e non esegue calcoli
quantitativi.

\subsection{I tre canali, e perché sono tre}

\begin{description}[leftmargin = 2.2em, style = nextline, itemsep = 6pt]
  \item[\textbf{1} --- identità e permessi]
    La dashboard chiede al sito «chi è questo utente, e che cosa può fare».
    La risposta è un \term{token} personale con dentro il livello di
    abbonamento e il flag di amministratore. Senza di esso l'applicazione
    funziona lo stesso: non si entra nel marketplace, e basta.
  \item[\textbf{2} --- chiavi e impostazioni]
    Le variabili d'ambiente --- chiavi Polygon, FRED, Stripe, SMTP --- sono
    gestite dal sito, non da un \file{.env} locale. È il verso opposto al
    canale 1, ed è l'unico che va da destra a sinistra.
  \item[\textbf{3} --- artefatti]
    Una strategia venduta sul marketplace viaggia come pacchetto di file. Va
    su quando si pubblica, viene giù quando si installa.
\end{description}

\begin{nota}{Dove cercare un bug che «sparisce»}
Quando un comportamento si vede in un ambiente e non nell'altro, la domanda da
farsi è quale dei tre canali è in mezzo. Un marketplace vuoto in locale e
pieno sul sito è il canale 1 (token mancante o scaduto); un download dati che
fallisce solo su una macchina è il canale 2 (chiave non arrivata); un
artefatto installato che non compare è il canale 3 più una delle trappole del
sistema a plugin (\vediCap{estendere}).
\end{nota}

\section{Le due applicazioni dentro la metà locale}\label{sez:strangler-intro}

Dentro la metà locale ci sono, a loro volta, due applicazioni web: la
dashboard Flask storica e un'API FastAPI più recente, avviate separatamente e
in ascolto su porte diverse. Non è un residuo da ripulire: è una migrazione
deliberata, il cosiddetto \term{strangler fig}, e il \vediCap{architettura} la
spiega per esteso.

\figuraTikz{strangler}
  {Le due applicazioni web della metà locale. Entrambe passano dalle stesse
   porte e usano lo stesso motore: ciò che cambia è il guscio.}
  {strangler}

% ===========================================================================
%  sviluppatore/01-introduzione/03-vocabolario.tex
% ===========================================================================
\section{Dieci parole da fissare subito}\label{sez:vocabolario}

Nel codice e nell'interfaccia queste parole hanno un significato preciso, e
alcune significano qualcosa di diverso da ciò che si direbbe leggendole. Vale
la pena leggerle una volta: metà delle domande su \Lyra{} sono domande di
vocabolario travestite da domande di architettura.

\begin{description}[leftmargin = 0pt, style = nextline, itemsep = 7pt]

\item[\termdef{Strategia} (\pkgStrategie)]
  Un oggetto che, a ogni barra, guarda gli indicatori e il portafoglio e
  restituisce una lista di \term{segnali}. Non decide quanto comprare: decide
  \emph{che cosa} e in che direzione. Vive dentro il motore e passa per
  l'allocatore, il rischio e l'esecuzione.

\item[\termdef{Approccio} (\pkgApprocci)]
  Un oggetto che \emph{salta} tutto il motore e produce direttamente una serie
  di rendimenti giornalieri. Serve a chi ha già una pipeline propria --- un
  modello ML con il suo feature store, per dire --- e vuole solo che il
  risultato entri nel sistema come uno stream qualsiasi. Una strategia e un
  approccio non sono due nomi per la stessa cosa.

\item[\termdef{Stream}]
  Una serie storica di profitti e perdite. È l'unità di misura comune: una
  strategia, un approccio e un overlay producono tutti uno stream, ed è su
  questo che l'ensemble lavora. Tre formati: \cod{equity\_curve},
  \cod{pnl\_daily}, \cod{pnl\_sparse}.

\item[\termdef{Ensemble}]
  La combinazione pesata di più stream. Il peso di ciascuno è deciso da una
  \term{modalità di pesatura}.

\item[\termdef{Modalità di pesatura} (\pkgPesi)]
  Il criterio con cui l'ensemble distribuisce il capitale fra gli stream.
  Quattordici installate, dall'equiponderato alla minima varianza con
  restringimento di Ledoit--Wolf al meta-allocatore addestrato sui regimi. È
  un plugin: si aggiunge senza toccare il motore.

\item[\termdef{Overlay} (\pkgOverlay)]
  Un moltiplicatore di esposizione fra \num{0} e \num{1}, applicato
  all'ensemble già composto. \cod{1.0} significa esposizione piena,
  \cod{0.0} fuori dal mercato. Non produce rendimenti: li scala. Gli overlay
  attivi si moltiplicano fra loro, e il risultato non dipende dall'ordine ---
  proprietà voluta, non accidente.

\item[\termdef{Regime} (\pkgRegime)]
  Uno stato di mercato fra tre: \cod{bull}, \cod{neutral}, \cod{bear}. Si
  costruisce con due artefatti distinti: una \emph{regola} (dalla serie al
  regime) e un \emph{mercato} (quali simboli usano quale regola, con quali
  parametri). Le etichette sono sempre quelle tre, perché le strategie
  dichiarano i propri \cod{regime\_weights} su quei nomi.

\item[\termdef{Walk-forward}]
  Il modo in cui si valuta una strategia: si addestra su una finestra, si
  valuta sulla successiva, si avanza. Mai una valutazione sull'intero
  campione, che misurerebbe quanto bene il modello ricorda invece di quanto
  bene prevede.

\item[\termdef{Artefatto}]
  Un plugin impacchettato per essere pubblicato: strategia, approccio,
  overlay, modalità di pesatura, metrica, indicatore. È l'unità che il
  marketplace vende e installa.

\item[\termdef{Porta} (\pkgPorte)]
  Un contratto Python (\cod{Protocol}) fra il centro del sistema e il mondo:
  un repository, una sorgente dati, un broker, un orologio. Trentadue, e
  i casi d'uso conoscono solo queste.

\end{description}

\begin{attenzione}{Due parole che tradiscono}
\textbf{Overlay} in \pkgOverlay{} è un moltiplicatore di esposizione; ma
esistono anche approcci con \cod{stream\_category = "ml\_overlay"}, che
\emph{non} sono overlay: sono approcci, con un contratto diverso. Il
\vediCap{estendere} tiene i due casi separati con cura, perché confonderli
produce un artefatto che si installa e non compare mai.

\textbf{Project} nell'interfaccia non è il progetto software: è una
composizione salvata di stream, pesi e overlay. Nella \emph{User guide}
è un capitolo a sé.
\end{attenzione}

% ===========================================================================
%  sviluppatore/01-introduzione/04-come-leggere.tex
% ===========================================================================
\section{Da dove cominciare, secondo cosa si deve fare}\label{sez:percorsi}

Questa guida non va letta tutta. Quattro percorsi, a seconda del motivo per
cui si è qui.

\begin{center}
\begin{tabular}{@{}L{48mm}L{98mm}@{}}
\toprule
\textbf{Si deve\ldots} & \textbf{Leggere} \\
\midrule
\rigaTenue
far partire una copia e
guardarsi intorno
  & \vediCap{ambiente}. Un'ora, e si ha il sistema che gira. \\

aggiungere una strategia,
un overlay, una pesatura
  & \vediCap{motore} per capire che cosa decide che cosa, poi il
    \vediCap{estendere}, che è una procedura da seguire con il codice
    aperto. \\

\rigaTenue
toccare l'interfaccia
  & \vediCap{dashboard}. Prima però \vediSez{strangler}: c'è un motivo
    se ci sono due applicazioni web e non una. \\

mettere mano a pagamenti,
abbonamenti, marketplace
  & \vediCap{marketplace}, e la parte sui permessi nel
    \vediCap{architettura}. È l'area dove un errore costa denaro vero. \\
\bottomrule
\end{tabular}
\end{center}

\begin{pratica}{Il minimo indispensabile}
Chi ha fretta e deve comunque toccare il codice legga due sole cose: la
separazione fra decisione ed esecuzione (\vediSez{decisione-esecuzione}) e le
due forme del sistema a plugin (\vediSez{due-forme}). Sono i due punti in cui
un errore non dà errore: produce un numero sbagliato, o un artefatto che non
compare.
\end{pratica}

\section{Che cosa questa guida non copre}\label{sez:fuori-perimetro}

\begin{itemize}
  \item \textbf{Il contenuto finanziario delle strategie.} Perché il momentum
        12-1 funzioni, se il PEAD sia ancora sfruttabile, come si tari un
        \emph{vol target}: sono domande di merito, e questo è un documento di
        ingegneria. Dove una scelta numerica è codificata nel sistema, la
        guida dice dov'è e che cosa cambia, non se sia giusta.
  \item \textbf{L'uso dell'applicazione.} Dove si clicca per lanciare un
        backtest sta nel \emph{Manuale utente}, documento separato nella
        stessa cartella.
  \item \textbf{Il codice del sito, riga per riga.} Il sito è documentato qui
        per ciò che espone alla metà locale: endpoint, formato dei token,
        regole dei permessi. Il resto è un'applicazione Next.js ordinaria.
\end{itemize}


% ===========================================================================
%  sviluppatore/02-ambiente/00-capitolo.tex
% ===========================================================================
\chapter{Mettere in moto una copia}\label{cap:ambiente}

% ===========================================================================
%  sviluppatore/02-ambiente/01-prerequisiti.tex
% ===========================================================================
\section{Che cosa serve avere installato}\label{sez:prerequisiti}

Per la metà locale: Python, e nient'altro di obbligatorio. Per il sito:
Node.js e un Postgres raggiungibile. Le due metà si avviano in modo
indipendente, e per lavorare sul motore il sito non serve --- senza di esso si
perde solo il marketplace.

\begin{center}
\begin{tabular}{@{}L{34mm}L{26mm}L{78mm}@{}}
\toprule
\textbf{Componente} & \textbf{Versione} & \textbf{Note} \\
\midrule
\rigaTenue
Python & \num{3}.10 o più & Il codice usa \cod{match}, \cod{X | None} e
  le \cod{dataclass} con \cod{slots}. \\
pandas & $<$\,\num{3} & Vincolo reale: il codice precede l'API di pandas~3. \\
\rigaTenue
numpy & $<$\,\num{2} & Stesso motivo. In caso di errori nello stack
  numerico, questi due vincoli sono la prima cosa da provare. \\
Node.js & \num{18} o più & Solo per il sito. \\
\rigaTenue
Postgres & \num{14} o più & Solo per il sito. In sviluppo va bene anche
  un'istanza gestita (Neon). \\
\bottomrule
\end{tabular}
\end{center}

\subsection{Installazione delle dipendenze Python}

\begin{codiceInterno}{console}{Dalla radice del progetto}{}
python -m venv .venv
source .venv/bin/activate        # Windows: .venv\Scripts\activate
pip install -r requirements.txt
\end{codiceInterno}

\noindent
Il file \file{requirements.txt} è diviso in blocchi commentati e dichiara
di che cosa è responsabile ogni dipendenza. Tre di essi vanno guardati:

\begin{description}[leftmargin = 1.6em, style = nextline, itemsep = 4pt]
  \item[Stack numerico] \pkg{numpy}, \pkg{pandas}, \pkg{scipy},
    \pkg{scikit-learn}, \pkg{statsmodels}, \pkg{pyarrow}. Senza
    \pkg{pyarrow} non si leggono i \file{parquet}, cioè non si legge niente.
  \item[Machine learning] \pkg{lightgbm}, \pkg{hmmlearn}, \pkg{shap}. Servono
    agli approcci ML e alle modalità di pesatura \cod{regime\_ml} e
    \cod{regime\_style}. Un sistema senza di esse parte, e quelle voci
    semplicemente non funzionano.
  \item[Sorgenti dati] \pkg{yfinance}, \pkg{requests}, \pkg{fredapi},
    \pkg{python-dotenv}.
\end{description}

\begin{nota}{Il modello locale è opzionale}
Lo Studio può generare il codice di un artefatto a partire da una descrizione
in italiano. Il modello gira in locale via Ollama, che non è nelle
dipendenze Python e va installato a parte. Senza, lo Studio funziona
ugualmente: si perde la generazione da testo, restano l'editor e il grafo a
nodi.
\end{nota}

\subsection{Con Docker, invece}

Nella radice ci sono \file{Dockerfile} e \file{docker-compose.yml}. Avviano la
sola dashboard, con i dati montati come volume.

\begin{codiceInterno}{console}{Dalla radice del progetto}{}
UID=$(id -u) GID=$(id -g) docker compose up --build
# poi: http://localhost:5000
\end{codiceInterno}

\begin{attenzione}{Su Linux servono UID e GID}
Senza passarli, i file che il contenitore scrive nei volumi --- i
\file{parquet} scaricati, i risultati dei backtest --- risultano di un altro
proprietario e dall'esterno non si riescono più a modificare. Su macOS e
Windows il problema non si pone.
\end{attenzione}

% ===========================================================================
%  sviluppatore/02-ambiente/02-avvio.tex
% ===========================================================================
\section{I quattro modi di avviarla}\label{sez:avvio}

Nella radice del progetto ci sono quattro punti d'ingresso, e fanno cose
diverse. Si lanciano tutti \emph{dalla radice}, mai da dentro una
sottocartella: gli import sono assoluti rispetto a quella.

\begin{center}
\begin{tabular}{@{}L{52mm}L{18mm}L{66mm}@{}}
\toprule
\textbf{Comando} & \textbf{Porta} & \textbf{Che cosa avvia} \\
\midrule
\rigaTenue
\file{python scripts/run\_dashboard.py}
  & \num{5000} & L'interfaccia completa. È questo che si usa nel
                 \num{99}\,\% dei casi. \\
\file{python run\_api.py}
  & \num{8000} & Solo l'API FastAPI del motore. Nessuna pagina, solo JSON. \\
\rigaTenue
\file{python desktop\_app.py}
  & \num{5000} & La stessa dashboard, in una finestra nativa invece che in una
                 scheda del browser. Richiede \pkg{pywebview}. \\
\file{npm run dev} (in \file{sito/})
  & \num{3000} & Il sito. Indipendente dai tre precedenti. \\
\bottomrule
\end{tabular}
\end{center}

\subsection{La dashboard}

\begin{codiceInterno}{console}{Avvio normale}{}
python scripts/run_dashboard.py
# http://127.0.0.1:5000
\end{codiceInterno}

\noindent
Host e porta si cambiano con le variabili d'ambiente \cod{FLASK\_HOST} e
\cod{FLASK\_PORT}, oppure nei campi corrispondenti della configurazione. Il
default è \cod{127.0.0.1}, cioè raggiungibile solo dalla macchina stessa, e
questo non è un dettaglio di comodità.

\begin{attenzione}{La dashboard non va esposta in rete}
Quasi nessuna rotta della dashboard richiede un login, ed è una scelta
deliberata: l'applicazione gira sul computer di chi la usa, e chiedere un
account per entrare in casa propria non avrebbe senso. Il prezzo è che chi
raggiunge la porta \num{5000} può fare tutto: lanciare backtest, leggere le
credenziali del broker, scrivere file. Mettere \cod{FLASK\_HOST=0.0.0.0} su
una macchina raggiungibile da fuori significa consegnarla. Le uniche aree
protette sono quelle che toccano il sito --- marketplace, profilo,
abbonamento, amministrazione --- e il \vediSez{permessi} spiega
esattamente quali e perché.
\end{attenzione}

\subsection{L'API del motore}

\begin{codiceInterno}{console}{Con ricaricamento automatico}{}
uvicorn api.app:app --reload
# http://127.0.0.1:8000/docs  -> la documentazione interattiva generata
\end{codiceInterno}

\noindent
Gira \emph{accanto} alla dashboard, non al suo posto. Per lavorare
sull'interfaccia non serve; serve quando si tocca un caso d'uso, perché
l'API lo espone senza passare da Flask ed è il modo più rapido per provarlo.

\subsection{Il sito}

\begin{codiceInterno}{console}{Dalla cartella del sito}{}
cd sito
npm install
npm run dev
# http://localhost:3000
\end{codiceInterno}

\noindent
Il sito ha bisogno di un Postgres e delle proprie variabili d'ambiente
(\cod{DATABASE\_URL}, \cod{NEXTAUTH\_SECRET}, le chiavi Stripe e SMTP). In
produzione queste non stanno in un file: si gestiscono dal pannello del
sito, ed è lo stesso meccanismo che poi le passa alla metà locale attraverso
il canale~\num{2} (\vediSez{due-meta}).

\begin{pratica}{Il giro minimo per vedere qualcosa}
\begin{enumerate}
  \item \file{pip install -r requirements.txt}
  \item \file{python scripts/download\_d1.py} --- scarica un po' di dati
        giornalieri (\vediSez{dati}).
  \item \file{python scripts/run\_dashboard.py}
  \item Nel browser: \voce{Run} \freccia{} una strategia \freccia{}
        \pulsante{Run backtest}. Alla fine compare uno stream in
        \voce{Streams}.
\end{enumerate}
Senza il passo~\num{2} l'applicazione parte e le pagine sono vuote: non è un
errore, è che non c'è niente da mostrare.
\end{pratica}

% ===========================================================================
%  sviluppatore/02-ambiente/03-dati.tex
% ===========================================================================
\section{Dove stanno i dati}\label{sez:dati}

Tutti i percorsi del sistema sono definiti in un file solo,
\file{core/paths.py}, e nessun altro modulo costruisce percorsi a mano. Chi
deve aggiungere una cartella la aggiunge lì.

\subsection{La cartella di lavoro}

La radice non è necessariamente quella del repository. \simbolo{PROJECT\_ROOT}
si decide a ogni avvio, in questo ordine:

\begin{enumerate}
  \item \cod{LYRA\_HOME} impostata nell'ambiente \freccia{} si usa quella. È il
        modo per tenere i dati su un disco diverso, o condividerli fra più
        installazioni.
  \item Applicazione impacchettata (PyInstaller, Electron) \freccia{} la
        cartella utente del sistema operativo:
        \file{\%LOCALAPPDATA\%/Lyra} su Windows,
        \file{\textasciitilde/Library/Application Support/Lyra} su macOS,
        \file{\textasciitilde/.local/share/Lyra} su Linux.
  \item Altrimenti \freccia{} la radice del repository.
\end{enumerate}

\begin{perche}{Perché i dati escono dal bundle}
In un'applicazione impacchettata il codice vive in una cartella di sola
lettura che ogni aggiornamento sovrascrive. Se i \file{parquet} e i risultati
dei backtest stessero lì dentro, un aggiornamento cancellerebbe mesi di
lavoro. \simbolo{CODE\_ROOT} --- dove sta il codice --- e
\simbolo{PROJECT\_ROOT} --- dove si scrive --- sono due cose distinte proprio
per questo, e vanno tenute distinte anche quando coincidono, come in sviluppo.
\end{perche}

\subsection{Che cosa c'è sotto \file{data/}}

\begin{center}
\begin{tabular}{@{}L{40mm}L{104mm}@{}}
\toprule
\textbf{Cartella} & \textbf{Contenuto} \\
\midrule
\rigaTenue
\file{data/assets/<TF>/}
  & OHLCV, un \file{parquet} per simbolo, divisi per timeframe:
    \file{D1/}, \file{H4/}, \file{H1/}, \file{M30/}, \file{M15/},
    \file{M5/}. Lo stesso simbolo può esistere a più timeframe, ciascuno
    con il proprio intervallo di date. \\
\file{data/driver/}
  & Le serie che guidano i regimi (SPY, DXY, VIX\ldots). Stanno a parte
    perché possono avere un timeframe diverso dagli asset. \\
\rigaTenue
\file{data/rates/}      & Tassi di politica monetaria per valuta, da FRED. \\
\file{data/news/}       & Archivio delle notizie. \\
\rigaTenue
\file{data/earnings/}   & Date e sorprese degli utili, per simbolo. \\
\file{data/fundamentals/} & Bilanci, per simbolo. \\
\rigaTenue
\file{data/mms/}        & Feature della heatmap di liquidità, per
                          (simbolo, timeframe). \\
\file{data/\_universe.json}
  & Che cosa il sistema sa di avere. Non è un indice generato: è il
    registro, e una cartella piena di \file{parquet} non censiti qui per il
    sistema è vuota. \\
\bottomrule
\end{tabular}
\end{center}

\noindent
I risultati stanno altrove, sotto \file{output/}: \file{output/strategies/}
per gli stream (è da qui che la dashboard li legge, seguendo il
\cod{output\_path\_pattern} dichiarato da ciascun artefatto),
\file{output/ml/} per feature e predizioni, \file{output/reports/} per i
report aggregati.

\begin{nota}{I percorsi legacy sono ancora letti}
\file{tests/outputs/}, \file{data/feature\_store/}, \file{data/ml\_outputs/}
sono posizioni storiche. L'helper \simbolo{discover\_parquet()} cerca prima
nei percorsi nuovi e poi in questi, così una installazione vecchia continua a
funzionare. Scrivere lì non è più previsto.
\end{nota}

\subsection{Portare dentro i primi dati}

Gli script di scarico stanno in \file{scripts/} e sono uno per sorgente:

\begin{codiceInterno}{console}{I quattro che servono per cominciare}{scripts/}
python scripts/download_d1.py              # azioni ed ETF, giornaliero
python scripts/download_regime_drivers.py  # le serie dichiarate dai mercati di regime
python scripts/download_fred_rates.py      # tassi (serve FRED_API_KEY)
python scripts/download_polygon_m5.py      # intraday (serve POLYGON_API_KEY)
\end{codiceInterno}

\noindent
Lo stesso si fa dall'interfaccia, in \voce{Data}, che chiama le rotte
\cod{/api/data/download/<modalità>} e mostra l'avanzamento. La differenza è
che dall'interfaccia il lavoro parte come \term{job} con un registro
consultabile, dalla riga di comando no.

\begin{attenzione}{\cod{download\_regime\_drivers.py} legge i file di mercato}
Non ha un elenco di simboli suo: scarica i \cod{driver} dichiarati nei file
sotto \file{regime/markets/}. Se si aggiunge un mercato che guarda una serie
nuova, questo script è il modo di procurarsela --- e finché non lo si lancia,
quel mercato non modula niente e lo dice solo in un avviso nel registro.
\end{attenzione}

% ===========================================================================
%  sviluppatore/02-ambiente/04-cartelle.tex
% ===========================================================================
\section{La mappa delle cartelle}\label{sez:cartelle}

Trentuno cartelle nella radice. Raggruppate per mestiere diventano sei
famiglie, e la famiglia dice già quasi tutto su che cosa ci si può aspettare
di trovarci dentro.

\subsection{Il guscio a strati}

Sono le cartelle introdotte con la migrazione all'architettura esagonale
(\vediSez{esagonale}). Sono piccole: la loro dimensione non misura quanto
lavoro fanno, ma quanto poco codice serve a tenere separato il resto.

\begin{center}
\begin{tabular}{@{}L{30mm}R{14mm}L{86mm}@{}}
\toprule
\textbf{Cartella} & \textbf{Righe} & \textbf{Che cosa contiene} \\
\midrule
\rigaTenue
\pkgDominio      & \num{478}  & I tipi puri: \cod{dataclass} congelate, nessuna
                                dipendenza, nemmeno da pandas. \\
\pkgPorte        & \num{239}  & I trentadue contratti. È il file
                                da leggere per primo di tutto il progetto. \\
\rigaTenue
\pkgApplicazione & \num{579}  & I casi d'uso. Conoscono solo le porte. \\
\pkgAdattatori   & \num{2593} & Le implementazioni concrete, e i validatori
                                \simbolo{AST} degli artefatti. \\
\rigaTenue
\pkgComposizione & \num{109}  & Centonove righe, ed è l'unico posto in tutto
                                il sistema dove si sceglie un concreto. \\
\pkgApi          & \num{548}  & L'applicazione FastAPI e i suoi tredici
                                router. \\
\bottomrule
\end{tabular}
\end{center}

\subsection{Il motore}

\begin{center}
\begin{tabular}{@{}L{30mm}R{14mm}L{86mm}@{}}
\toprule
\textbf{Cartella} & \textbf{Righe} & \textbf{Che cosa contiene} \\
\midrule
\rigaTenue
\pkgCore        & \num{7587} & Tipi del dominio quantitativo
                               (\cod{Signal}, \cod{Portfolio}, \cod{Fill}),
                               percorsi, manutenzione. \\
\pkgBacktest    & \num{1567} & Due file. Uno è il motore di simulazione, ed è
                               il punto in cui si decide che cosa è lecito
                               guardare. \\
\pkgEsecuzione  & \num{2428} & I runner: walk-forward, backtest singolo,
                               approcci, gestione dei job. \\
\rigaTenue
\pkgDati        & \num{6251} & Caricamento, universo, categorie, qualità,
                               sorgenti esterne. \\
\pkgReport      & \num{2096} & \simbolo{EnsembleEngine}: la composizione degli
                               stream, i pesi, gli overlay. \\
\rigaTenue
\pkgLive        & \num{4170} & Lo stesso calcolo, collegato a un broker. \\
\bottomrule
\end{tabular}
\end{center}

\subsection{Gli artefatti}

Sono le cinque famiglie di plugin, più tre cartelle che ospitano artefatti
minori. Ogni sottocartella è un artefatto, e il nome della cartella è il nome
canonico --- non quello dichiarato dentro il file.

\begin{center}
\begin{tabular}{@{}L{30mm}R{18mm}L{82mm}@{}}
\toprule
\textbf{Cartella} & \textbf{Installati} & \textbf{Forma del contratto} \\
\midrule
\rigaTenue
\pkgStrategie & \num{9} attive, \num{20} spente
  & \cod{DESCRIPTOR} in \file{strategy.py} \\
\pkgApprocci  & \num{3} attivi, \num{4} dormienti
  & \cod{APPROACH\_DESCRIPTOR} in \file{approach.py} \\
\rigaTenue
\pkgPesi      & \num{14}
  & \cod{META} + \cod{compute} in \file{weight\_mode.py} \\
\pkgOverlay   & \num{7}
  & \cod{META} + \cod{compute} in \file{overlay.py} \\
\rigaTenue
\pkgRegime    & \num{3} regole, \num{3} mercati
  & \cod{META} + \cod{compute}, in due cartelle distinte \\
\pkg{indicators/}, \pkg{metrics/}
  & ---
  & Artefatti minori, stesso schema \\
\bottomrule
\end{tabular}
\end{center}

\begin{pratica}{Accendere e spegnere un artefatto}
Il prefisso \cod{\_} davanti al nome della cartella è l'interruttore:
\file{strategies/\_value/} non viene caricata, \file{strategies/value/} sì.
Rinominare la cartella è l'unico gesto necessario --- non c'è nessun elenco da
aggiornare da nessuna parte. Delle \num{29} cartelle sotto \pkgStrategie,
\num{20} cominciano con il trattino basso e restano fuori.
\end{pratica}

\subsection{L'interfaccia, e il resto}

\pkgDashboard{} è l'applicazione Flask: \num{28494} righe di Python in
\num{36} file, più \num{42072} righe fra template Jinja, JavaScript e CSS. È
l'oggetto più grande del progetto e ha un capitolo suo
(\vediCap{dashboard}).

Restano: \pkg{features/}, \pkg{mms/} e \pkg{volume\_profile/} (calcoli
specialistici che le strategie possono richiedere dichiarandoli nel
descrittore), \pkg{genetic/} (la generazione di strategie per programmazione
genetica), \pkg{validation/} (controlli statistici sui risultati),
\pkg{scripts/} (gli scarichi e le utilità da riga di comando), \pkg{tests/},
\pkg{config/} (il solo \file{quantsys.yaml}) e \pkg{build/} (impacchettamento).

\begin{attenzione}{\pkg{brokers/} è morta}
Nella radice c'è una cartella \pkg{brokers/} di \num{402} righe che duplica
codice vivo in \pkgLive. Nessun modulo la importa. È un residuo: non va
estesa, e ciò che sembra doverci andare va in \file{live/broker.py}.
\end{attenzione}



% --- PARTE II — come è fatto ------------------------------------------------
\part{Come è fatto}
% ===========================================================================
%  sviluppatore/03-architettura/00-capitolo.tex
% ===========================================================================
\chapter{Come è fatto}\label{cap:architettura}

% ===========================================================================
%  sviluppatore/03-architettura/01-esagonale.tex
% ===========================================================================
\section{L'architettura esagonale, in concreto}\label{sez:esagonale}

Il guscio a strati occupa \num{4546} righe su \num{119084}: meno del
\num{4}\,\% del Python del progetto. Non ha riscritto niente. Avvolge il
codice che c'era --- \pkgStrategie, \pkgReport, \pkgEsecuzione, \pkgLive,
\pkgDati{} --- e gli mette davanti un confine.

\figuraTikz{esagonale}
  {I cinque strati. L'unica regola è la direzione delle frecce.}
  {esagonale}

\subsection{La regola, e perché è una sola}

Si dipende solo verso l'interno. \pkgApplicazione{} non conosce nessun
concreto: conosce solo le porte. I concreti si scelgono in un file solo,
\file{composition/container.py}.

\begin{perche}{Perché il confine vale la pena}
L'ecosistema che \Lyra{} vuole essere --- applicazione scaricabile, sito
community, marketplace --- è fatto di \emph{tre client dello stesso motore}.
Finché il motore è una libreria che ciascun client importa a modo suo, ogni
client ne vede una parte diversa e ogni cambiamento ne rompe due. Rendendolo
un servizio con un contratto stabile, i client tornano a essere consumatori, e
l'inversione delle dipendenze non è più una buona pratica da ricordarsi: è
imposta dal confine di processo.
\end{perche}

\subsection{\pkgDominio: i tipi, e basta}

Quattrocentosettantotto righe di \cod{dataclass} congelate. Nessuna
dipendenza --- nemmeno pandas. \simbolo{StrategyInfo}, \simbolo{StreamInfo},
\simbolo{RunReport}, \simbolo{OrderRequest}, \simbolo{SecurityReport},
\simbolo{LiveRunnerConfig} e una trentina d'altri.

Sono il vocabolario che attraversa gli strati. Un caso d'uso restituisce un
\simbolo{StrategyInfo}, non un \simbolo{DataFrame} e non un
\simbolo{StrategyDescriptor} --- che è un oggetto del motore, con dentro una
\cod{factory} richiamabile.

\subsection{\pkgPorte: trentadue contratti}

Il file da leggere per primo in tutto il progetto, e sono
\num{239} righe. Ogni porta è un \cod{Protocol} marcato
\cod{@runtime\_checkable}.

\begin{codiceInterno}{python}{Due porte intere, non un estratto}{ports/\_\_init\_\_.py}
@runtime_checkable
class StrategyRepository(Protocol):
    """Sorgente delle strategie disponibili (mappate a StrategyInfo)."""
    def list(self) -> list[StrategyInfo]: ...
    def get(self, name: str) -> Optional[StrategyInfo]: ...


@runtime_checkable
class WalkForwardRunner(Protocol):
    """Esegue un walk-forward per una strategia. Ritorna l'exit code."""
    def run(self, strategy_name: str) -> int: ...
\end{codiceInterno}

\noindent
Sono \cod{Protocol}, non classi base astratte, e la differenza conta:

\begin{description}[leftmargin = 1.6em, style = nextline, itemsep = 4pt]
  \item[Tipizzazione strutturale] Un concreto soddisfa la porta
    \emph{senza ereditarla}. Non deve importare \pkgPorte, non deve
    dichiarare niente: le basta avere i metodi giusti. L'accoppiamento è il
    minimo possibile, e si può avvolgere codice preesistente senza
    toccarlo.
  \item[Porte piccole e separate] \simbolo{DataSource} $\neq$
    \simbolo{Broker} $\neq$ \simbolo{StrategyRepository} $\neq$
    \simbolo{Clock}. Un adattatore implementa quello che gli serve. È
    l'\emph{Interface Segregation} preso sul serio: una porta
    «\simbolo{Everything}» avrebbe costretto ogni finto a implementare
    trenta metodi per provarne uno.
  \item[Sostituibilità verificata] Nei test i casi d'uso girano con
    repository e runner finti. Non è una proprietà sperata: è il modo in cui
    le prove sono scritte.
\end{description}

\begin{nota}{La porta \simbolo{Broker} è ri-esportata, non riscritta}
\simbolo{live.broker.LiveBroker} era già una classe base astratta, cioè già
una porta, scritta prima del guscio. \pkgPorte{} la ri-esporta invece di
duplicarla:

\begin{codiceInterno}{python}{}{ports/\_\_init\_\_.py}
try:
    from live.broker import LiveBroker as Broker  # noqa: F401
except Exception:  # pragma: no cover
    Broker = object  # type: ignore
\end{codiceInterno}

\noindent
È la mossa corretta: due definizioni della stessa interfaccia divergono, e
quando divergono nessuno se ne accorge finché un broker non si comporta in un
modo in un posto e in un altro altrove.
\end{nota}

\subsection{\pkgApplicazione: i casi d'uso}

Ogni caso d'uso è una classe con una responsabilità, riceve le porte dal
costruttore e non conosce nessun concreto. Sono brevi --- spesso tre righe --- e
devono restarlo: il caso d'uso orchestra, non calcola.

\begin{codiceInterno}{python}{Tre casi d'uso, per intero}{application/\_\_init\_\_.py}
class ListStrategies:
    def __init__(self, repo: StrategyRepository) -> None:
        self._repo = repo

    def execute(self) -> list[StrategyInfo]:
        return self._repo.list()


class RunWalkForward:
    """Esegue un walk-forward e (se c'e' un sink) ne registra il report."""
    def __init__(self, runner: WalkForwardRunner,
                 report_sink: Optional[ReportSink] = None) -> None:
        self._runner = runner
        self._sink = report_sink
\end{codiceInterno}

\begin{attenzione}{Il calcolo non va qui}
La tentazione, aggiungendo una funzionalità, è di metterne «un pezzetto» nel
caso d'uso: un filtro, una media, una conversione. Il giorno in cui succede,
quel pezzetto esiste solo nel percorso che passa per l'API e non in quello
che passa per la dashboard, e i due danno numeri diversi. Il calcolo sta nel
motore; il caso d'uso lo chiama e basta.
\end{attenzione}

\subsection{\pkgComposizione: centonove righe, e sono quelle decisive}

\file{composition/container.py} è un \cod{dataclass} con venti campi --- un
caso d'uso ciascuno --- e una funzione che lo costruisce.

\begin{codiceInterno}{python}{L'unico punto di scelta dei concreti}{composition/container.py}
    # --- unico punto di scelta dei CONCRETI ---
    repo = RegistryStrategyRepository(Path(strategies_dir) if strategies_dir else None)
    weighting = EnsembleWeightingModeProvider()
    runner = InProcessWalkForwardRunner()
    registry = FileModelRegistryAdapter(model_dir)
    sink = FileReportSink(report_dir)
    ...
    return Container(
        list_strategies=ListStrategies(repo),
        get_strategy=GetStrategy(repo),
        run_walkforward=RunWalkForward(runner, sink),
        ...
    )
\end{codiceInterno}

\noindent
Per cambiare sorgente dati, broker o registro dei modelli si tocca questo
file, e nessun altro. Per far girare i casi d'uso nei test con dei finti si
costruisce un \simbolo{Container} diverso, senza toccare né i casi d'uso né
le porte.

\begin{pratica}{Come si aggiunge una capacità, nell'ordine giusto}
\begin{enumerate}
  \item Un tipo in \pkgDominio, se serve un dato nuovo che attraversa il
        confine.
  \item Una porta in \pkgPorte: che cosa deve saper fare il mondo.
  \item Un caso d'uso in \pkgApplicazione: tre righe, nessun calcolo.
  \item Un adattatore in \pkgAdattatori{} che avvolge il codice del motore
        che fa davvero la cosa.
  \item Una riga in \file{composition/container.py} che li collega.
  \item Un router in \pkgApi{} che lo espone, se serve via HTTP.
\end{enumerate}
Saltare il passo~\num{2} è l'errore tipico: si collega il caso d'uso
direttamente all'adattatore, «tanto è uno solo», e da quel momento il caso
d'uso non è più verificabile senza l'infrastruttura.
\end{pratica}

% ===========================================================================
%  sviluppatore/03-architettura/02-strangler.tex
% ===========================================================================
\section{Due applicazioni web, e il motivo}\label{sez:strangler}

Dentro la metà locale ci sono due applicazioni web (\vediFig{strangler}):

\begin{center}
\begin{tabular}{@{}L{34mm}L{24mm}L{80mm}@{}}
\toprule
 & \textbf{Flask} & \textbf{FastAPI} \\
\midrule
\rigaTenue
File principale & \file{dashboard/app.py} & \file{api/app.py} \\
Porta           & \num{5000}              & \num{8000} \\
\rigaTenue
Dimensione      & \num{249} rotte         & \num{13} router \\
Restituisce     & pagine HTML e JSON      & solo JSON \\
\rigaTenue
Autenticazione  & guardia globale, quasi tutto aperto & nessuna \\
Parla con       & direttamente il motore  & i casi d'uso, via
                                            \simbolo{Container} \\
\bottomrule
\end{tabular}
\end{center}

\subsection{Che cos'è uno strangler fig}

La pianta da cui prende il nome cresce attorno a un albero ospite, lo
avvolge e lo sostituisce gradualmente: alla fine resta in piedi da sola, e
l'albero di dentro non c'è più. Nessun momento in cui il bosco resta senza
albero.

Applicato qui: la dashboard Flask ha dentro il codice che è stato scritto per
primo, con il motore chiamato direttamente. L'API FastAPI è il nuovo guscio e
parla solo con i casi d'uso. Le due convivono, e il confine si sposta una
funzionalità alla volta.

\begin{perche}{Perché non si riscrive la dashboard e basta}
Perché la dashboard funziona, e perché \num{28494} righe di Python più
\num{42072} di template non si riscrivono con la garanzia di non perdere
niente. Una riscrittura in blocco ha due esiti noti: o dura abbastanza da
essere abbandonata a metà, o arriva in fondo avendo perso silenziosamente i
casi limite che nessuno aveva scritto da nessuna parte --- e in un sistema che
muove denaro, un caso limite perso è un ordine sbagliato.
\end{perche}

\subsection{Come si avvia l'API}

\begin{codiceInterno}{python}{Il montaggio dell'applicazione}{api/app.py}
def create_app(container: Optional[Container] = None, jobs=None) -> FastAPI:
    container = container or build_container()
    app = FastAPI(title="QuantSys Engine API", version="0.2.0")
    app.state.container = container

    if jobs is None:
        # JobManager reale: max_workers=1 = serializzato (default sicuro del progetto)
        from execution.job_manager import JobManager
        job_manager = JobManager(max_workers=1)
        app.state.job_manager = job_manager
        jobs = JobManagerJobService(container.run_walkforward, job_manager)
    app.state.jobs = jobs

    app.include_router(health.router)
    app.include_router(strategies.router)
    ...
\end{codiceInterno}

\noindent
Due cose meritano attenzione. La prima è che \simbolo{create\_app} accetta un
\simbolo{Container} dall'esterno: nei test si passa un contenitore di finti e
l'API intera gira senza toccare il disco. La seconda è
\cod{max\_workers=1}: i job sono \emph{serializzati}, ed è il default sicuro
del progetto. Due backtest in parallelo si contendono la stessa cache dati e
gli stessi file di output.

\subsection{I tredici router}

\begin{center}
\begin{tabular}{@{}L{32mm}L{112mm}@{}}
\toprule
\textbf{Router} & \textbf{Che cosa espone} \\
\midrule
\rigaTenue
\file{health}             & Stato del servizio. \\
\file{strategies}         & Elenco e dettaglio delle strategie. \\
\rigaTenue
\file{strategy\_admin}    & Validazione e installazione di una strategia. \\
\file{runs}               & Avvio di un walk-forward, stato, log in SSE. \\
\rigaTenue
\file{streams}            & Catalogo degli stream prodotti. \\
\file{ensemble}           & Composizione: pesi, overlay, metriche. \\
\rigaTenue
\file{models}             & Registro dei modelli ML e promozione. \\
\file{factor\_attribution}& Scomposizione del rendimento per fattore. \\
\rigaTenue
\file{stress\_test}       & Prove sotto scenario. \\
\file{live}               & Controllo dell'operatività. \\
\rigaTenue
\file{generator}          & Generazione di codice da descrizione. \\
\file{nodes}              & Compilazione del grafo a nodi in Python. \\
\rigaTenue
\file{fidelity}           & Confronto fra trade disegnati a mano e strategia. \\
\bottomrule
\end{tabular}
\end{center}

\begin{attenzione}{Le rotte FastAPI non esistono per Flask}
Sono due applicazioni separate, in due processi. Nessuna rotta di \pkgApi{} è
registrata nell'app Flask, e viceversa. Cercare \percorsoHTTP{/api/strategies} sulla
porta \num{5000} e trovarci una rotta Flask con lo stesso nome ma
comportamento diverso è possibile, ed è successo: quando si documenta o si
collauda un endpoint, la porta fa parte dell'indirizzo.
\end{attenzione}

\subsection{Dove passa il confine, oggi}

La regola pratica: \textbf{ciò che è nuovo nasce nel guscio}. Una funzionalità
nuova si scrive come caso d'uso, si espone dall'API, e la dashboard la chiama
come chiamerebbe qualunque servizio. Ciò che esiste si sposta quando lo si
deve toccare comunque, non prima e non per ordine --- spostare codice che
funziona e che nessuno sta modificando è rischio senza ritorno.

% ===========================================================================
%  sviluppatore/03-architettura/03-rotte.tex
% ===========================================================================
\section{Le \num{246} rotte, e chi le registra}\label{sez:rotte}

La dashboard Flask ha \num{267} decoratori \cod{@route} che danno
\num{246} percorsi distinti --- la differenza sono le coppie
«\cod{GET} in una funzione, \cod{POST} in un'altra». Sono sparsi su undici
moduli, montati in un ordine che conta.

\begin{center}
\begin{tabular}{@{}L{36mm}R{14mm}L{86mm}@{}}
\toprule
\textbf{Modulo} & \textbf{Rotte} & \textbf{Come si registra} \\
\midrule
\rigaTenue
\file{app.py}             & \num{134} & \cod{@app.route} dentro
                                        \simbolo{create\_app()} \\
\file{community.py}       & \num{43}  & \cod{Blueprint("community")},
                                        senza \cod{url\_prefix} \\
\rigaTenue
\file{studio.py}          & \num{43}  & \cod{Blueprint("studio")} \\
\file{marketplace\_routes.py} & \num{12} & Agganciate \emph{al blueprint
                                        community}, non a uno proprio \\
\rigaTenue
\file{agent.py}           & \num{10}  & \cod{Blueprint("agent")} \\
\file{genetic.py}         & \num{8}   & \cod{Blueprint("genetic")} \\
\rigaTenue
\file{auth.py}            & \num{8}   & \cod{@app.route} dentro
                                        \simbolo{init\_auth(app)} \\
\file{ui\_ext\_a/b/e/f.py} & \num{9}  & Caricate automaticamente
                                        (\vediSez{ui-ext}) \\
\bottomrule
\end{tabular}
\end{center}

\begin{nota}{Due moduli senza rotte}
\file{live\_runners.py} e \file{autopubblica.py} non registrano niente: sono
moduli di servizio. Il primo è chiamato da \file{app.py} sotto
\percorsoHTTP{/api/runners/*} e \percorsoHTTP{/api/brokers/*}; il secondo parte come thread
all'avvio.
\end{nota}

\subsection{L'ordine di montaggio}\label{sez:montaggio}

\figuraTikz{montaggio-app}
  {Gli otto passi con cui \simbolo{create\_app()} assembla l'applicazione.}
  {montaggio-app}

\subsection{Il caricamento automatico delle estensioni v2}\label{sez:ui-ext}

L'ultimo passo di registrazione è diverso dagli altri: non c'è nessun elenco.

\begin{codiceInterno}{python}{Il meccanismo delle estensioni}{dashboard/app.py:4362}
    # Rotte aggiunte con il nuovo stile (v2): ogni modulo dashboard/ui_ext_*.py
    # espone register(app, data_layer, base_path). Tenerle in moduli propri
    # evita di far crescere ancora questo file.
    try:
        import importlib as _il, pkgutil as _pu, dashboard as _dpkg
        for _m in sorted(m.name for m in _pu.iter_modules(_dpkg.__path__)
                         if m.name.startswith("ui_ext_")):
            _il.import_module(f"dashboard.{_m}").register(app, data_layer, base_path)
    except Exception as _e:  # noqa: BLE001
        logging.getLogger("app").warning("v2 page extensions not loaded: %s", _e)
\end{codiceInterno}

\noindent
Cinque conseguenze, tutte rilevanti per chi estende:

\begin{enumerate}
  \item \textbf{Nessun registro da aggiornare.} Si crea
        \file{dashboard/ui\_ext\_<lettera>.py} con una funzione
        \simbolo{register(app, data\_layer, base\_path)} e al riavvio c'è.
  \item \textbf{Ordine alfabetico garantito} dal \cod{sorted}: \file{a},
        \file{b}, \file{e}, \file{f}. Le lettere \file{c} e \file{d} non
        esistono --- un buco nella serie, non un errore.
  \item \textbf{Rotte sull'app, non su un blueprint.} Gli endpoint sono
        globali (\cod{a\_overview}, \cod{b\_project\_detail}\ldots) e
        \cod{url\_for} li raggiunge senza prefisso.
  \item \textbf{Caricate dopo tutto il resto}, compreso
        \simbolo{init\_auth}: gli hook di queste estensioni girano dopo la
        guardia di autenticazione.
  \item \textbf{Fallimento silenzioso.} Un solo \cod{except} avvolge l'intero
        ciclo: se \file{ui\_ext\_a} solleva all'import, \emph{anche b, e, f non
        vengono caricati}, e tutto ciò che si vede è una riga di avviso nel
        registro.
\end{enumerate}

\begin{attenzione}{Il punto più fragile della dashboard}
Il punto~\num{5} è il modo tipico in cui metà dell'interfaccia sparisce dopo
una modifica apparentemente innocua. Il sintomo è una pagina che va in
\cod{404} o un pannello vuoto; la causa è un \cod{ImportError} in un altro
modulo. Quando qualcosa della v2 non c'è più, la prima cosa da cercare nel
registro è la riga \cod{v2 page extensions not loaded}.
\end{attenzione}

\subsection{Rotte statiche che coprono rotte dinamiche}

Werkzeug dà la precedenza ai segmenti statici, quindi
\percorsoHTTP{/api/marketplace/items} vince su \percorsoHTTP{/api/marketplace/<artifact\_id>}.
Funziona, ed è voluto. Ma significa che un artefatto il cui identificativo
fosse \cod{items}, \cod{mine}, \cod{acquisti}, \cod{diagnostica},
\cod{contestazioni}, \cod{video-check} o \cod{system-artifacts} sarebbe
\emph{irraggiungibile}. Lo stesso vale per \percorsoHTTP{/marketplace/sell} contro
\percorsoHTTP{/marketplace/<artifact\_id>}, e per \percorsoHTTP{/api/jobs/conclusi} contro
\percorsoHTTP{/api/jobs/<job\_id>}.

\begin{nota}{Perché oggi non è un problema}
Gli identificativi degli artefatti sono generati, non scelti dall'utente. Se
un giorno diventassero scegliibili --- uno \emph{slug} leggibile nell'URL, per
dire --- questa lista diventerebbe un elenco di parole riservate da vietare in
fase di pubblicazione.
\end{nota}

\subsection{Tre cose morte, da non seguire}

\begin{description}[leftmargin = 1.6em, style = nextline, itemsep = 5pt]
  \item[\rotta{POST}{/api/market/publish}]
    Ritirata. Il docstring dice \cod{"""RITIRATA."""} e la rotta risponde
    sempre \cod{403}, ma resta registrata e occupa il percorso. La
    pubblicazione vera è \rotta{POST}{/api/marketplace/publish}, in
    \file{community.py}.
  \item[L'endpoint \cod{auth\_google\_token}]
    Elencato fra quelli pubblici in \file{auth.py}, ma nessuna rotta ha
    quel nome. Riferimento pendente, innocuo.
  \item[Il prefisso \percorsoHTTP{/claim}]
    Elencato fra i percorsi che richiedono un account, ma nessuna rotta
    comincia con \percorsoHTTP{/claim}: la pagina delle contestazioni è
    \percorsoHTTP{/marketplace/rimborso/<claim\_id>}, già coperta dal prefisso
    \percorsoHTTP{/marketplace}. L'equivoco nasce dal template, che si chiama
    \file{claim.html}.
\end{description}

\subsection{Due template orfani}

In \file{dashboard/templates/} ci sono \num{35} file. Due non sono raggiunti
da nessuna rotta e da nessun \cod{include}:

\begin{itemize}
  \item \file{brokers.html} --- il suo JavaScript è stato estratto in
        \file{static/js/brokers.js} e l'interfaccia dei broker vive ora dentro
        \file{live.html}. Residuo di una ristrutturazione.
  \item \file{server\_runners.html} --- i runner non stanno più sul server, ma
        sul computer dell'utente. Funzionalità rimossa, template rimasto.
\end{itemize}

\noindent
Altri quattro non sono orfani ma non sono nemmeno pagine: \file{base.html} è
il genitore di tutte le \num{32} pagine, \file{\_robustness\_panel.html},
\file{\_v2\_profile\_security.html} e \file{\_v2\_studio\_bp.html} sono
frammenti inclusi. \file{serve\_account.html} è un caso a sé: non lo rende
nessuna rotta, lo rende la guardia di autenticazione
(\vediSez{permessi}).

\begin{attenzione}{\file{live.html} non è la pagina di Studio}
\file{studio.html} è servito da \percorsoHTTP{/studio} e da
\percorsoHTTP{/studio/dati-esterni}; \file{live.html} è servito da
\percorsoHTTP{/studio/runner-live}. Il nome del percorso dice Studio, la pagina è
Live. È facile sbagliare, e il blueprint che la registra si chiama
\cod{studio}.
\end{attenzione}

% ===========================================================================
%  sviluppatore/03-architettura/04-permessi.tex
% ===========================================================================
\section{Chi può fare che cosa}\label{sez:permessi}

Cercando un decoratore \cod{@login\_required} sulle rotte della dashboard non
se ne trova quasi nessuno. Non è una dimenticanza: il controllo è uno solo,
globale, e sta in \file{dashboard/auth.py}.

\begin{codiceInterno}{python}{La guardia, ridotta all'osso}{dashboard/auth.py:1474}
@app.before_request
def _require_login():
    if request.endpoint in _PUBLIC_ENDPOINTS or request.path.startswith("/static/"):
        return None
    if not session.get("uid") and (not _serve_account(request.path)
                                   or request.path in _ANONIMO_AMMESSO):
        g.user = None
        return None            # <-- si lavora ANONIMI
    if not session.get("uid"):
        if "/api/" in request.path or ...:
            return jsonify({"error": "not signed in"}), 401
        return render_template("serve_account.html", destinazione=request.path), 200
    g.user = {...}
\end{codiceInterno}

\figuraTikz{guardia-auth}
  {Il percorso di una richiesta. Il ramo di destra è quello che prende la
   grande maggioranza delle rotte.}
  {guardia-auth}

\subsection{Le tre liste}

\begin{description}[leftmargin = 0pt, style = nextline, itemsep = 6pt]

\item[\cod{\_SERVE\_ACCOUNT} --- i prefissi che richiedono un account]
  L'elenco è scritto a mano, e il confronto avviene per segmenti:
  \percorsoHTTP{/marketplace}, \percorsoHTTP{/profile}, \percorsoHTTP{/security}, \percorsoHTTP{/admin},
  \percorsoHTTP{/messages}, \percorsoHTTP{/api/market}, \percorsoHTTP{/api/marketplace},
  \percorsoHTTP{/api/account}, \percorsoHTTP{/api/admin}, \percorsoHTTP{/api/messages},
  \percorsoHTTP{/api/community}, \percorsoHTTP{/claim}.

\item[\cod{\_PUBLIC\_ENDPOINTS} --- gli esenti]
  \cod{auth\_login}, \cod{auth\_logout}, \cod{auth\_register},
  \cod{auth\_verify}, \cod{auth\_reset}, \cod{static}, più i quattro
  dell'accesso con Google. Sono nomi di \emph{endpoint}, non percorsi.

\item[\cod{\_ANONIMO\_AMMESSO} --- l'eccezione dentro l'eccezione]
  Una voce sola: \percorsoHTTP{/api/marketplace/claim/aperte}. Alimenta il contatore
  nella barra laterale, e senza account deve restituire un elenco vuoto
  invece di un \cod{401} --- altrimenti ogni pagina mostrerebbe un errore a
  chi non ha fatto l'accesso.

\end{description}

\subsection{I quattro livelli, sopra la guardia}

\begin{center}
\begin{tabular}{@{}L{26mm}L{44mm}L{72mm}@{}}
\toprule
\textbf{Livello} & \textbf{Come si verifica} & \textbf{Che cosa succede se manca} \\
\midrule
\rigaTenue
anonimo  & ---                           & Si lavora: tutto il motore è accessibile. \\
con account & \cod{session["uid"]}       & Pagina \file{serve\_account.html} con
                                           stato \cod{200}, oppure \cod{401} in JSON. \\
\rigaTenue
amministratore & \cod{auth.is\_admin()}  & \cod{403}, oppure \cod{404} per le
                                           sezioni che non devono rivelare di
                                           esistere. \\
super-amministratore & \cod{auth.is\_super\_admin()} & \cod{403}. È l'unico che
                                           può nominare o togliere un
                                           amministratore. \\
\bottomrule
\end{tabular}
\end{center}

\begin{perche}{Perché \cod{404} e non \cod{403} su alcune rotte di amministrazione}
Un \cod{403} dice «esiste, e non puoi». Un \cod{404} dice «non c'è». Per le
sezioni che elencano utenti, feedback e artefatti di tutti, la seconda
risposta non concede nemmeno l'informazione che quella sezione esista. La
distinzione è usata in modo consapevole: \percorsoHTTP{/admin}, \percorsoHTTP{/api/admin/users},
\percorsoHTTP{/api/admin/feedback}, \percorsoHTTP{/api/admin/artefatti} e \percorsoHTTP{/api/admin/beta}
rispondono \cod{404}; le altre rotte di amministrazione, che riguardano
impostazioni di cui l'esistenza non è un segreto, rispondono \cod{403}.
\end{perche}

\subsection{La scelta di fondo, e il suo prezzo}

L'applicazione gira sul computer di chi la usa. Obbligare all'accesso avrebbe
significato chiedere un account per entrare in casa propria --- per questo
\file{app.py} (tranne la parte di amministrazione), \file{studio.py},
\file{genetic.py} e \file{agent.py} sono aperti per intero.

Il prezzo è che \textbf{la porta \num{5000} va considerata una console di
amministrazione}: chi la raggiunge può lanciare backtest, leggere le
credenziali del broker e scrivere file sul disco. È una decisione, non una
falla, ma implica che l'applicazione non vada esposta in rete.

\subsection{Quattro punti dove la protezione non è dove ci si aspetta}

Conseguenze dirette del fatto che \cod{\_SERVE\_ACCOUNT} è un elenco scritto
a mano: ciò che non è elencato nasce aperto.

\begin{center}
\begin{tabular}{@{}L{42mm}L{102mm}@{}}
\toprule
\textbf{Percorso} & \textbf{Stato} \\
\midrule
\rigaTenue
\rotta{GET}{/subscription}
  & La pagina «Il mio abbonamento» non è fra i prefissi e non fa controlli:
    si apre senza account. I dati però arrivano da \percorsoHTTP{/api/abbonamento*} e
    \percorsoHTTP{/api/f/subscription/invoices}, che rispondono \cod{401}: si vede
    l'impalcatura vuota. \\
\rotta{GET}{/feedback},\newline\rotta{GET}{/account}
  & Chiamano \cod{\_uid()} senza verificarne l'esito. Si aprono, ed
    elencano niente. \\
\rigaTenue
\percorsoHTTP{/api/abbonamento*},\newline\percorsoHTTP{/api/e/*}, \percorsoHTTP{/api/f/*},\newline\percorsoHTTP{/api/feedback}
  & Fuori dai prefissi: \emph{ciascuno si difende da sé} con un controllo
    scritto nella rotta. Funziona, ma un endpoint nuovo aggiunto a queste
    famiglie nasce senza protezione. \\
\rotta{POST}{/api/feedback} e\newline\rotta{POST}{/api/community/feedback}
  & Due endpoint per la stessa cosa, con autenticazione diversa: il primo si
    difende da sé, il secondo è coperto dal prefisso. Il secondo prima si
    chiamava come il primo. \\
\bottomrule
\end{tabular}
\end{center}

\begin{pratica}{Aggiungere una rotta che tocca l'account}
\begin{enumerate}
  \item Metterla sotto uno dei prefissi già elencati in
        \cod{\_SERVE\_ACCOUNT}. È il modo che non si dimentica.
  \item Se il percorso non può stare lì, aggiungere il prefisso nuovo a
        quella lista --- non un controllo dentro la rotta.
  \item Decidere se la rotta debba rispondere \cod{403} o \cod{404} a chi non
        è amministratore, e usare \simbolo{\_require\_admin()} per il secondo
        caso.
  \item Se la risposta deve esserci anche per gli anonimi (un contatore, un
        badge), aggiungerla a \cod{\_ANONIMO\_AMMESSO} e restituire un
        risultato vuoto, non un errore.
\end{enumerate}
\end{pratica}

\subsection{Gli endpoint sull'account sono quattro}

\rotta{GET}{/api/account/security} in \file{app.py}, \rotta{GET}{/api/f/account} in
\file{ui\_ext\_f.py}, \rotta{GET}{/api/e/account-data} in \file{ui\_ext\_e.py} e
\rotta{POST}{/api/account} in \file{community.py}. Tre moduli, quattro percorsi,
tre stili di autenticazione diversi. Non è un difetto che rompa qualcosa, ma
è il posto dove cercare quando un dato dell'account è aggiornato in una
schermata e vecchio in un'altra.


% ===========================================================================
%  sviluppatore/04-motore/00-capitolo.tex
% ===========================================================================
\chapter{Il motore quantitativo}\label{cap:motore}

% ===========================================================================
%  sviluppatore/04-motore/01-barra.tex
% ===========================================================================
\section{Che cosa succede a ogni barra}\label{sez:barra}

Il ciclo principale sta in \file{backtest/engine.py}. Per ogni data del
calendario, dopo il periodo di riscaldamento, fa sei cose nello stesso ordine.

\figuraTikz{pipeline-motore}
  {I sei passi del ciclo. La divisione fra la riga superiore e quella
   inferiore non è grafica: è la divisione fra decidere ed eseguire.}
  {pipeline-motore}

\begin{enumerate}
  \item \textbf{Indicatori.} Non si calcolano dentro la strategia: si
        pre-calcolano una volta per simbolo e si passano già pronti. Che cosa
        serve lo dichiara ogni strategia nel proprio descrittore, e il
        registro ne fa l'unione.
  \item \textbf{Strategie.} Ciascuna riceve gli indicatori e il portafoglio e
        restituisce una lista di \simbolo{Signal}: che cosa, in che
        direzione, con quale convinzione. Non quanto.
  \item \textbf{Allocatore.} Dai segnali ai pesi di portafoglio. È qui che
        si decide quanto, e il come è configurabile
        (\vediSez{allocatore}).
  \item \textbf{Rischio.} Tetti su peso per posizione, esposizione lorda per
        classe di attività, leva complessiva, volatilità obiettivo.
  \item \textbf{Esecuzione.} Prezzo di apertura, slippage, costi.
  \item \textbf{Portafoglio.} Posizioni aggiornate, equity segnata a mercato,
        storia registrata.
\end{enumerate}

\subsection{Prima di tutto: i trigger intra-barra}

C'è un passo che precede la decisione, e non compare nella figura perché non
è parte del ciclo decisionale: il controllo degli \emph{stop} e dei
\emph{take profit} sulle posizioni già aperte.

\begin{codiceInterno}{python}{Il controllo che viene prima di tutto}{backtest/engine.py:745}
            # ---- INTRA-BAR TRIGGER CHECK ----
            # PRIMA della decision/execution phase, controlla se i SL/TP delle
            # posizioni aperte sono stati toccati nella bar t (high/low).
            trigger_hits = exit_tracker.check_bar(t, data)
\end{codiceInterno}

\noindent
Se un livello è stato toccato, la posizione si chiude subito e il simbolo
finisce in una lista nera fino al prossimo ribilanciamento vero. La filosofia
dichiarata nel codice è che «ribilanciare è quando rivaluti la tesi»: uno
stop toccato non aspetta il prossimo ribilanciamento, ma nemmeno autorizza a
riaprire la stessa posizione nella stessa barra.

\begin{attenzione}{La lista nera non blocca le chiusure}
Un simbolo in lista nera non può essere \emph{aperto}, ma un target a zero
passa sempre. Invertire questa asimmetria --- cioè bloccare anche le chiusure
--- significherebbe tenere aperta una posizione che il sistema ha deciso di
chiudere, ed è il tipo di errore che si scopre su un conto reale.
\end{attenzione}

% ===========================================================================
%  sviluppatore/04-motore/02-decisione.tex
% ===========================================================================
\section{Decisione ed esecuzione sono due istanti}\label{sez:decisione-esecuzione}

È la regola più importante del motore, e si riduce a due righe di codice.

\begin{codiceInterno}{python}{Le due righe che tengono in piedi tutto}{backtest/engine.py:740}
            # t_execution = t (open T)
            # t_decision = data point precedente in all_dates (close T-1)
            t_execution = t
            t_decision = all_dates[i - 1]
\end{codiceInterno}

\figuraTikz{decisione-esecuzione}
  {I due istanti. Tutto ciò che sta nella fascia a destra non esiste ancora
   quando si decide.}
  {decisione-esecuzione}

\subsection{Perché due istanti e non uno}

Una strategia che decide alla chiusura di $T$ e compra alla chiusura di $T$
non è implementabile: nel momento in cui si conosce il prezzo di chiusura, la
sessione è finita. Le implementazioni ingenue lo fanno lo stesso, e il
backtest che ne esce è eccellente.

Qui la decisione avviene con i dati disponibili fino alla chiusura di $T-1$ e
l'ordine viene eseguito al prezzo di apertura di $T$:

\begin{codiceInterno}{python}{La fase di esecuzione}{backtest/engine.py:945}
            # ---- EXECUTION PHASE: open T+1 = t_execution ----
            # Per ogni target ottengo l'open price a t_execution.
            ...
                df_sym = data.get(sym)
                if df_sym is None or t_execution not in df_sym.index:
                    # Symbol non disponibile, skip silenzioso
                    continue
                open_price = float(df_sym.loc[t_execution, "open"])
\end{codiceInterno}

\subsection{La firma che non si può cambiare}

Le strategie ricevono i due istanti come argomenti \emph{nominati}:

\begin{codiceInterno}{python}{Il contratto}{backtest/engine.py:834}
                        t_decision=t_decision,
                        t_execution=t_execution,
\end{codiceInterno}

\noindent
Il motore li passa per nome, quindi una strategia che chiami i suoi parametri
in modo diverso fa fermare il backtest alla prima barra con un
\cod{TypeError}. Il messaggio di errore del motore lo dice esplicitamente,
ricordando la firma richiesta:

\begin{codiceInterno}{python}{}{backtest/engine.py:857}
                            "Firma richiesta: generate_signals(self, t_decision, "
                            "t_execution, indicators_by_symbol, portfolio) -> "
\end{codiceInterno}

\begin{attenzione}{Il bug che migliora i risultati}
Dentro \simbolo{generate\_signals} si legge solo fino a \cod{t\_decision}:

\begin{center}
\begin{tabular}{@{}L{62mm}L{76mm}@{}}
\toprule
\textbf{Corretto} & \textbf{Lookahead} \\
\midrule
\cod{ind.close.loc[:t\_decision]} & \cod{ind.close.loc[:t\_execution]} \\
\cod{serie.shift(1)} prima di decidere & \cod{serie} senza shift \\
\bottomrule
\end{tabular}
\end{center}

\noindent
La seconda colonna non dà errore, non produce avvisi, e fa salire lo Sharpe.
È per questo che la separazione è imposta dalla firma invece che raccomandata
in un commento: una regola che si può violare senza accorgersene viene
violata.
\end{attenzione}

\subsection{La stessa regola altrove}

La causalità non riguarda solo le strategie. Si ripresenta identica in tre
altri punti, e in ciascuno è scritta nel codice come vincolo:

\begin{center}
\begin{tabular}{@{}L{34mm}L{104mm}@{}}
\toprule
\textbf{Dove} & \textbf{Come} \\
\midrule
\rigaTenue
Regime & Lo snapshot è calcolato con \cod{as\_of=True}, che usa solo dati
         anteriori a \cod{t\_decision}. \\
Overlay & Il moltiplicatore si calcola dal PnL \emph{base} e si applica
          shiftato: «la decisione di oggi non può usare il rendimento di
          oggi». \\
\rigaTenue
Regole di regime & Niente statistiche sull'intero campione: un quantile
          calcolato su tutta la storia è lookahead travestito da soglia.
          Gli helper condivisi (\cod{\_expanding\_median},
          \cod{\_hysteresis}) esistono per non doverlo riscrivere ogni volta. \\
Normalizzazione della volatilità & La vol trailing viene shiftata di un
          giorno nel backtest (\vediSez{ensemble}). \\
\bottomrule
\end{tabular}
\end{center}

% ===========================================================================
%  sviluppatore/04-motore/03-allocatore.tex
% ===========================================================================
\section{Dall'intenzione alla posizione}\label{sez:allocatore}

Una strategia dice \emph{che cosa}; l'allocatore decide \emph{quanto}. Sta in
\file{core/allocator.py}, e il suo comportamento si governa con una sola
\cod{dataclass} di configurazione.

\subsection{Le quattro modalità}

\begin{center}
\begin{tabular}{@{}L{28mm}L{110mm}@{}}
\toprule
\cod{strategy\_mix\_mode} & \textbf{Come distribuisce} \\
\midrule
\rigaTenue
\cod{equal\_weight}
  & Gerarchico classico: il portafoglio si divide fra le strategie, poi
    dentro ciascuna. Il fattore $\sqrt{N_\text{strat}}$ assume indipendenza
    fra strategie. È il valore predefinito, ed è anche il ripiego quando le
    altre modalità non hanno i dati che servono. \\
\cod{stream\_level}
  & Ogni strategia opera come stream indipendente con un proprio budget di
    volatilità (\cod{stream\_target\_vol\_annual}, \num{0.12} di
    default); i contributi si sommano con peso $1/N$. \\
\rigaTenue
\cod{explicit}
  & Si esegue il vettore di pesi dichiarato dalla strategia, punto. Serve a
    chi ha già calcolato le proprie dimensioni e non vuole che nessuno le
    tocchi. \\
\cod{covariance}
  & Le correlazioni si \emph{stimano} invece che assumerle: la volatilità del
    libro è $\sqrt{w^{\mathsf T}\Sigma w}$, con $\Sigma$ da
    \file{core/covariance.py}. \\
\rigaTenue
\cod{bayesian}
  & Dichiarata e non implementata: solleva \cod{NotImplementedError}. È
    elencata nei valori ammessi perché il messaggio d'errore è più utile di
    un «valore non valido». \\
\bottomrule
\end{tabular}
\end{center}

\begin{perche}{Perché \cod{explicit\_vol\_target} è un campo a sé}
In modalità \cod{explicit} si può chiedere che il libro venga scalato una
volta sola perché la sua volatilità prevista sia un certo valore --- le
proporzioni fra i titoli restano identiche, cambia solo quanto pesa il libro
nel suo insieme. Il parametro non è \cod{portfolio\_target\_vol\_annual}
perché quello ha sempre un valore (ha un default), e con un default non si
distingue «lo voglio» da «c'era». Qui l'assenza significa qualcosa:
\cod{None} vuol dire «esegui ciò che ti ho dato, senza scalare».
\end{perche}

\subsection{I tetti di rischio}

\begin{center}
\begin{tabular}{@{}L{44mm}R{16mm}L{76mm}@{}}
\toprule
\textbf{Parametro} & \textbf{Default} & \textbf{Che cosa limita} \\
\midrule
\rigaTenue
\cod{portfolio\_target\_vol\_annual} & \num{0.10} & Volatilità annua
  obiettivo del portafoglio. \\
\cod{max\_position\_weight} & \num{0.10} & $|w|$ del singolo strumento. \\
\rigaTenue
\cod{max\_asset\_class\_gross} & \num{0.20} & $\sum |w|$ dentro una classe
  di attività. \\
\cod{max\_gross\_leverage} & \num{1.5} & $\sum |w|$ totale. \\
\rigaTenue
\cod{min\_vol\_floor} & \num{0.05} & Pavimento sulla volatilità stimata. \\
\bottomrule
\end{tabular}
\end{center}

\begin{nota}{Il ragionamento dietro il tetto per classe}
Il commento nel codice lo dice senza giri: quattro obbligazioni tutte lunghe
fanno il \num{40}\,\% lordo senza alcun tetto, e in situazione di stress sono
una posizione sola. Il tetto al \num{20}\,\% limita a «un'esposizione
obbligazionaria con leva~2 al massimo», che è realistico. Non è un numero
scelto per prudenza generica: è la risposta a un modo preciso in cui una
diversificazione apparente sparisce.
\end{nota}

\subsection{La validazione è nel costruttore}

\simbolo{\_\_post\_init\_\_} controlla ogni campo e solleva \cod{ValueError}
con il valore ricevuto nel messaggio. Due vincoli meritano attenzione perché
non sono ovvi:

\begin{itemize}
  \item \cod{max\_gross\_leverage >= 1.0}. Un tetto inferiore a~\num{1}
        significherebbe non poter mai investire tutto il capitale, ed è quasi
        sempre un errore di unità di misura.
  \item \cod{max\_asset\_class\_gross} ammesso fino a~\num{4.0}. Il limite è
        stato ampliato per sostenere i portafogli lungo-corto, che arrivano
        a un lordo di~\num{2} e oltre.
\end{itemize}

\begin{attenzione}{La configurazione è congelata}
\simbolo{AllocatorConfig} è una \cod{dataclass} congelata: non si modifica un
campo dopo averla costruita. Per variarne uno in un test si usa
\cod{dataclasses.replace(config, max\_gross\_leverage=2.0)}, che ne
restituisce una copia nuova. Assegnare direttamente solleva
\cod{FrozenInstanceError}, ed è voluto: una configurazione di rischio
modificata a metà di un backtest renderebbe i risultati non riproducibili.
\end{attenzione}

% ===========================================================================
%  sviluppatore/04-motore/04-ensemble.tex
% ===========================================================================
\section{Comporre più stream}\label{sez:ensemble}

Una strategia da sola è uno stream. \simbolo{EnsembleEngine}, in
\file{reporting/ensemble.py}, ne prende diversi e ne fa uno solo. Il percorso
ha tre stadi, e due di essi sono punti d'innesto per i plugin.

\figuraTikz{composizione-ensemble}
  {La composizione. Pesatura e overlay fanno cose diverse e non vanno
   confuse: la prima distribuisce, la seconda scala.}
  {composizione-ensemble}

\subsection{Primo stadio: portare gli stream alla stessa scala}

Due strategie con lo stesso Sharpe ma volatilità diversa non vanno pesate
allo stesso modo: la più volatile dominerebbe il risultato
indipendentemente dal peso dichiarato. Ogni stream viene quindi riportato a
una volatilità obiettivo comune prima di essere pesato.

Il mattone è uno solo, e serve sia al backtest sia all'operatività dal vivo:

\begin{codiceInterno}{python}{Il mattone comune}{reporting/ensemble.py:335}
    def _vol_trailing(self, returns: pd.Series) -> pd.Series:
        """Vol annualizzata TRAILING e NON shiftata: il valore a t usa le
        osservazioni fino a t INCLUSO.

        E' il mattone comune fra backtest e live. Il backtest la shifta di un
        giorno (`_compute_scale_series`) perche' lo scale applicato al ritorno
        di t deve conoscere solo il passato; il live la usa non shiftata
        (`scala_corrente`) perche' li' l'ultima osservazione e' il ritorno
        appena realizzato e la posizione che si sta per aprire rendera' nel
        periodo SUCCESSIVO. Due usi, una sola formula: e' l'unico modo perche'
        il libro che si esegue sia quello che si e' misurato.
        """
\end{codiceInterno}

\begin{perche}{Perché uno shift di differenza non è un dettaglio}
Nel backtest il fattore applicato al rendimento di $t$ deve conoscere solo il
passato, quindi la serie si shifta di uno. Dal vivo l'ultima osservazione
\emph{è} il rendimento appena realizzato, e la posizione che si sta per
aprire renderà nel periodo successivo: shiftare anche lì significherebbe
dimensionare con una stima vecchia di un giorno. Stesso numero, due usi, e
se le due strade calcolassero la volatilità ciascuna per conto proprio il
libro eseguito non sarebbe quello misurato. Il motivo per cui la funzione è
una sola è esattamente questo.
\end{perche}

\noindent
Due costanti governano i casi limite:

\begin{center}
\begin{tabular}{@{}L{32mm}R{14mm}L{92mm}@{}}
\toprule
\textbf{Costante} & \textbf{Valore} & \textbf{A che serve} \\
\midrule
\rigaTenue
\cod{SCALE\_MIN\_OBS} & \num{20} & Sotto questa soglia non c'è stima, e il
  fattore vale \cod{1.0}: il libro resta quello dell'allocatore, mai
  amplificato da una volatilità stimata su quattro dati. \\
\cod{SCALE\_MAX\_LEVA} & \num{10.0} & Con volatilità che tende a zero il
  fattore esploderebbe. La volatilità ha quindi un pavimento a
  \cod{target\_vol / SCALE\_MAX\_LEVA}. \\
\bottomrule
\end{tabular}
\end{center}

\subsection{Secondo stadio: la pesatura}

Quattordici modalità installate, raggruppate per famiglia:

\begin{center}
\begin{tabular}{@{}L{34mm}L{18mm}L{86mm}@{}}
\toprule
\textbf{Famiglia} & \textbf{Tipo} & \textbf{Modalità} \\
\midrule
\rigaTenue
Ingenua & statica & \cod{equal} \\
Rischio / covarianza & statica & \cod{inverse\_cov}, \cod{min\_variance\_lw},
  \cod{min\_variance\_rmt}, \cod{erc} \\
\rigaTenue
Robusta & statica & \cod{hrp}, \cod{nco}, \cod{cluster\_risk\_budget},
  \cod{max\_div} \\
Coda & statica & \cod{cvar} \\
\rigaTenue
Con vista sui rendimenti & statica & \cod{black\_litterman}, \cod{kelly} \\
Adattiva al regime & \textbf{rolling} & \cod{regime\_ml}, \cod{regime\_style} \\
\bottomrule
\end{tabular}
\end{center}

\noindent
Il campo \cod{kind} del \cod{META} governa il dispatch, e non è
documentazione: è la condizione di un \cod{if}.

\begin{codiceInterno}{python}{Il dispatch, senza una catena di if}{reporting/ensemble.py:491}
        spec = next((m for m in self.WEIGHTING_MODE_REGISTRY
                     if m["key"] == self.weighting_mode), None)
        if spec is None:
            raise ValueError(f"Unknown weighting_mode: {self.weighting_mode}")
        weight_fn = _WEIGHTING_COMPUTERS[self.weighting_mode]
        if spec["kind"] == "static":
            weights = weight_fn(self, stream_pnls)
        else:  # rolling: pesi time-varying; per la UI esponi la MEDIA come scalare
            weights_over_time = weight_fn(self, stream_pnls)
            mean_w = weights_over_time.mean(axis=0)
            weights = {name: float(mean_w[name]) for name in stream_pnls}
\end{codiceInterno}

\begin{attenzione}{Il modello documenta un valore che non esiste}
\file{weights/WEIGHT\_MODE\_TEMPLATE.md} dice che \cod{kind} vale
\cod{"static"} o \cod{"dynamic"}. Nel codice il confronto è con
\cod{"static"}, e \emph{tutto il resto} finisce nel ramo rolling. I
quattordici plugin installati usano \cod{static} (dieci) e \cod{rolling}
(quattro); \cod{dynamic} non lo usa nessuno. Un plugin scritto seguendo il
modello funzionerebbe per caso. Il modello è da correggere, il codice no.
\end{attenzione}

\subsection{Terzo stadio: gli overlay}

Ogni overlay attivo produce un moltiplicatore in $(0,1]$ calcolato dal
PnL \emph{base}, cioè da quello prima di qualunque overlay. I moltiplicatori
si moltiplicano fra loro.

\begin{codiceInterno}{python}{L'applicazione degli overlay}{reporting/ensemble.py:555}
        # Overlay di esposizione CAUSALI (ortogonali al weighting_mode e componibili
        # fra loro). Ogni overlay attivo produce un moltiplicatore in (0,1] calcolato
        # dal PnL/stream BASE (=> ordine-indipendente); l'esposizione finale e' il
        # PRODOTTO dei moltiplicatori. Nessun overlay attivo => ensemble_pnl invariato.
        base_pnl = ensemble_pnl
        overlay_multipliers: dict[str, pd.Series] = {}
        for ov in self.OVERLAY_REGISTRY:
            if ov["key"] in self.enabled_overlays:
                m = _GATE_COMPUTERS[ov["key"]](self, base_pnl, stream_pnls)
                m = m.reindex(base_pnl.index).fillna(1.0)
                overlay_multipliers[ov["key"]] = m
        if overlay_multipliers:
            total_mult = pd.Series(1.0, index=base_pnl.index)
            for m in overlay_multipliers.values():
                total_mult = total_mult * m
            ensemble_pnl = base_pnl * total_mult
\end{codiceInterno}

\begin{nota}{L'indipendenza dall'ordine è una proprietà costruita}
Se ogni overlay calcolasse il proprio moltiplicatore dal PnL già modificato
dai precedenti, il risultato dipenderebbe dall'ordine in cui compaiono nel
registro --- cioè dal campo \cod{order}, che è un dato di presentazione.
Calcolandoli tutti dal PnL base, l'ordine smette di contare e attivare due
overlay dà lo stesso risultato comunque li si elenchi.
\end{nota}

\subsection{I sette overlay installati}

\begin{center}
\begin{tabular}{@{}L{26mm}L{32mm}L{80mm}@{}}
\toprule
\textbf{Chiave} & \textbf{Famiglia} & \textbf{Quando taglia} \\
\midrule
\rigaTenue
\cod{vol\_regime} & Volatilità & Quando la volatilità realizzata sale. Mira
  al drawdown massimo. \\
\cod{jump\_vol} & Volatilità & Reazione in un giorno a salti oltre $N$ sigma:
  copre il rischio di apertura che gli overlay lenti perdono. \\
\rigaTenue
\cod{corr\_gate} & Correlazione & Quando la correlazione media fra gli stream
  schizza, cioè quando la diversificazione sparisce. \\
\cod{regime\_gmm} & Stress & Miscela gaussiana su feature multivariate;
  riduzione graduale al crescere della probabilità di stress. \\
\rigaTenue
\cod{iso\_stress} & Stress & Isolation forest su dispersione, volatilità,
  ampiezza, correlazione media. \\
\cod{macro\_regime} & Macro & Fattore latente (PCA causale) su curva, spread
  di credito, VIX, disoccupazione. \\
\rigaTenue
\cod{news\_pressure} & Flusso informativo & Quando il volume di notizie
  supera la norma storica. È l'unico che dichiara dati esterni e una
  scadenza giornaliera. \\
\bottomrule
\end{tabular}
\end{center}

% ===========================================================================
%  sviluppatore/04-motore/05-live.tex
% ===========================================================================
\section{Lo stesso calcolo, su un conto vero}\label{sez:live}
% Sezione fitta di tabelle che non si spezzano: con il fondo pagina
% allineato, lo spazio avanzato finiva attorno a una tabella. Qui il fondo
% resta libero; alla fine della sezione si torna all'allineamento del libro.
\raggedbottom

\pkgLive{} non è un secondo motore: è lo stesso, con un broker al posto dei
dati storici. Il valore di tutto il progetto dipende da questa frase, e il
modo in cui viene mantenuta vera merita di essere letto per intero, perché è
anche l'esempio migliore di che cosa va storto quando non lo è.

\subsection{Il difetto che ha definito il progetto}

Fino a un certo punto, il backtest normalizzava la volatilità di ogni stream
e l'operatività dal vivo no. Il commento in testa a
\file{live/ensemble\_live.py} dice che cosa significava:

\begin{quote}\itshape
Senza il punto~2 un libro dichiarato al \num{4}\,\% di vol veniva mostrato
in backtest moltiplicato per \num{2.5} ed eseguito dal vivo a dimensione
piena: \num{0.4} volte il rischio e il rendimento misurati. Al contrario, un
libro al \num{25}\,\% veniva tagliato a \num{0.4} in backtest e girava dal
vivo a \num{2.5} volte.
\end{quote}

\noindent
Nessun errore, nessuna eccezione, nessun avviso. Il backtest mostrava un
numero e il conto ne produceva un altro, e la differenza era un fattore fra
\num{0.4} e \num{2.5} a seconda della volatilità dello stream --- cioè più
grande quanto più il libro si discostava dall'obiettivo, che è esattamente il
caso in cui si ha bisogno che il backtest sia attendibile.

\begin{perche}{Perché il rimedio è stato condividere la funzione}
Si sarebbe potuto riscrivere la normalizzazione in \pkgLive{} copiando la
formula. Due implementazioni della stessa formula restano identiche finché
qualcuno non ne tocca una, e questo difetto è nato esattamente così. La
soluzione è che \simbolo{\_vol\_trailing} è \emph{una}, e le due strade la
usano in modo diverso in un solo punto documentato: il backtest la shifta di
un giorno, il live no (\vediSez{ensemble}).
\end{perche}

\subsection{I due punti in cui la normalizzazione entra}

\begin{enumerate}
  \item \textbf{I pesi si calcolano sui rendimenti scalati}, non su quelli
        grezzi. Le modalità di pesatura guardano la matrice di covarianza: su
        serie non normalizzate uno stream al \num{25}\,\% di volatilità
        riceve un peso del tutto diverso da quello che riceveva in backtest.
  \item \textbf{\simbolo{scales()} restituisce il fattore da applicare adesso}
        al vettore di posizioni --- cioè la leva che il backtest stava
        mostrando e che il live non applicava.
\end{enumerate}

\subsection{Il riscaldamento}

La volatilità trailing non esiste prima di una ventina di osservazioni, e fra
venti e sessanta è ancora rumorosa. In quel tratto il sistema mescola la
stima con una volatilità \emph{a priori} --- per uno stream di approccio, la
sua volatilità di backtest --- con peso lineare crescente.

\begin{center}
\begin{tabular}{@{}L{40mm}L{98mm}@{}}
\toprule
\textbf{Osservazioni} & \textbf{Che cosa si usa} \\
\midrule
\rigaTenue
meno di \cod{SCALE\_MIN\_OBS} & Solo la volatilità a priori. Senza una, il
  fattore è \cod{1.0}. \\
fra le due soglie & Miscela lineare fra priori e stima. \\
\rigaTenue
oltre \cod{scale\_min\_obs} & Solo la stima: qui coincide con il backtest. \\
\bottomrule
\end{tabular}
\end{center}

\noindent
Il requisito dichiarato è: \emph{mai un salto di dimensione in un solo
giorno}. Una transizione netta fra priori e stima produrrebbe, il giorno in
cui la soglia viene superata, un ordine di ribilanciamento che non
corrisponde a nessun cambiamento del mercato.

\begin{nota}{Lo stato sopravvive ai riavvii}
I rendimenti per stream sono salvati su JSON, così la finestra non riparte da
zero a ogni riavvio del processo. Senza, ogni riavvio rimetterebbe il sistema
in riscaldamento, e il riscaldamento usa la volatilità a priori: il conto
cambierebbe dimensione per il solo fatto di essere stato riavviato.
\end{nota}

\subsection{I broker}

L'interfaccia è \simbolo{live.broker.LiveBroker}, la stessa classe che
\pkgPorte{} ri-esporta come porta \simbolo{Broker}. Le implementazioni
attuali coprono Interactive Brokers (\file{live/brokers/ibkr.py}),
MetaTrader~5 (\file{live/brokers/mt5.py}), l'esecuzione su carta e un broker
di prova con il tempo accelerato --- quest'ultimo serve alla pagina di
simulazione, che fa passare una giornata di mercato in pochi secondi.

Oltre ai metodi obbligatori, cinque \emph{ganci} con un default innocuo:

\begin{center}
\begin{tabular}{@{}L{46mm}L{92mm}@{}}
\toprule
\textbf{Metodo} & \textbf{A che serve} \\
\midrule
\rigaTenue
\simbolo{account\_label()} & Una riga che dice a quale conto ci si è
  collegati: numero, demo o reale, terminale, cambio applicato. Va nel
  registro all'avvio. \\
\simbolo{account\_identity()} & Chi è il conto, in un dizionario
  (\cod{broker}, \cod{account}, \cod{where}, \cod{kind}); \cod{None} se il
  broker non sa dirlo. Lega la conferma degli ordini reali a un conto
  (\vediSez{live-ordini-reali}). \\
\rigaTenue
\simbolo{real\_orders\_blockers()} & I motivi per cui il conto
  \emph{non} può ricevere ordini da Lyra. Se la lista non è vuota, un
  runner con ordini reali non parte; se compare mentre gira, i cicli si
  saltano. \\
\simbolo{open\_orders()} & Gli ordini ancora in lavorazione,
  \cod{\{simbolo: quanti\}}. Il motore non manda un ordine su un simbolo che
  ne ha uno in corso. \\
\rigaTenue
\simbolo{cancel\_open\_orders()} & Cancella gli ordini aperti \emph{di
  Lyra} (non quelli fatti a mano) e dice quanti. Il runner la chiama
  all'avvio e all'arresto. \\
\bottomrule
\end{tabular}
\end{center}

\noindent
L'esito di ogni ordine è un \simbolo{Fill} con \cod{status} fra
\cod{filled}, \cod{partial}, \cod{pending} (accettato, né eseguito né
cancellato), \cod{skipped} (sotto la quantità minima, non inviato),
\cod{rejected}, \cod{cancelled} ed \cod{error}, e una \cod{note} con il
perché nelle parole del broker. Un adattatore che non riesce a leggere il
conto solleva \simbolo{BrokerUnavailable} (una \simbolo{ConnectionError}):
il runner lo riprova invece di fermarsi. Anche un fill parziale
entra nel libro della strategia (\simbolo{LivePortfolioMirror.record\_fill}).

\begin{center}
\begin{tabular}{@{}L{30mm}L{108mm}@{}}
\toprule
\textbf{Adattatore} & \textbf{Che cosa garantisce} \\
\midrule
\rigaTenue
IBKR & L'ordine va al conto indicato (\cod{order.account}) e porta il
  riferimento del runner (\cod{orderRef} \cod{lyra:<id>}); equity e
  posizioni si leggono di quel conto e solo per i contratti \cod{STK}, così
  un'opzione su AAPL non viene letta come azioni AAPL. Un ordine non
  eseguito entro \cod{fill\_timeout\_s} viene cancellato; se la
  cancellazione non arriva entro \cod{cancel\_timeout\_s} l'esito è
  \cod{pending}, e un ordine \cod{Inactive} viene cancellato perché può
  riattivarsi. \simbolo{equity\_in\_valuta()} converte la
  \cod{NetLiquidation} della valuta base del conto con il tag
  \cod{ExchangeRate} (in riserva il cambio dei dati di Lyra): senza, un
  conto in euro risultava a equity zero. \simbolo{open\_orders()} legge
  \simbolo{reqAllOpenOrders()} (solo \cod{STK}, solo il conto);
  \simbolo{cancel\_open\_orders()} cancella solo i riferimenti
  \cod{lyra}. Il collegamento passa \cod{account} e un \cod{timeout} di
  \num{20}\,s. Posizioni, valori del conto e stati degli ordini di
  \pkg{ib\_insync} si aggiornano solo mentre il suo event loop gira, e fra
  un ciclo e l'altro il runner dorme: prima di leggere il conto
  \simbolo{\_aggiorna()} fa girare il loop e le posizioni si chiedono al
  server (\simbolo{reqPositions()}), altrimenti un ordine eseguito
  durante la notte risultava ancora da fare; se la richiesta fallisce
  l'errore sale, senza ripiegare sulla copia locale. Le richieste bloccanti
  hanno \cod{RequestTimeout} di \num{20}\,s (di default aspettano per
  sempre). Si importa \pkg{ib\_async}, il seguito mantenuto, oppure
  \pkg{ib\_insync}, nel thread che lo usa. \\
MT5 & Il riempimento segue \cod{symbol\_info.filling\_mode} (IOC, FOK o
  Return): imporre IOC fa rifiutare ogni ordine sui simboli che non lo
  accettano. Il retcode \num{10010} è un'esecuzione parziale, il
  \num{10008} un ordine accettato non ancora eseguito (\cod{pending}). Un
  conto \emph{hedging}, il trading disabilitato, \cod{trade\_expert}
  falso, \voce{Algo Trading} spento o \cod{tradeapi\_disabled} bloccano gli
  ordini reali. Un conto in un'altra valuta si converte con il tick della
  coppia nel terminale (\cod{EURUSD}, anche con il suffisso del broker).
  Gli ordini in attesa vengono da \simbolo{orders\_get()}; si cancellano
  con \cod{TRADE\_ACTION\_REMOVE} solo quelli con il \cod{magic} di Lyra.
  \cod{None} da \simbolo{positions\_get()}, \simbolo{orders\_get()} o
  \simbolo{account\_info()} è un errore dell'API, non un conto vuoto: letto
  come «nessuna posizione» faceva ricomprare il libro.
  \simbolo{is\_connected()} interroga il terminale, così dopo un suo
  riavvio ci si ricollega. Il pacchetto parla con un terminale per
  processo: le chiamate passano da un lucchetto, e \simbolo{shutdown()} lo
  chiama solo l'ultimo che si scollega. \\
\bottomrule
\end{tabular}
\end{center}

\begin{attenzione}{Le credenziali restano su file locale}
\percorsoHTTP{/api/brokers/credentials} legge e scrive un file sul computer
dell'utente. Non passano dal sito, non finiscono nel database, e non c'è
nessun percorso che le porti in rete. Questo è il motivo per cui la porta
della dashboard non va esposta: chi la raggiunge può leggerle. Per IBKR
porta e client ID sono controllati al salvataggio, e
\simbolo{live\_runners.cfg\_ibkr()} li passa al collegamento come
\cod{broker\_cfg}.
\end{attenzione}

\subsection{Simulazione o ordini reali}\label{sez:live-ordini-reali}

Un runner su un broker vero parte in simulazione: \cod{dry\_run} vero, il
motore calcola gli ordini e non li invia. Gli ordini reali richiedono tre
cose insieme, e nessuna basta da sola:

\begin{codiceInterno}{python}{Quando un runner invia davvero}{dashboard/live\_runners.py:216}
def ordini_reali(runner: dict) -> bool:
    r = runner or {}
    return (str(r.get("broker") or "") in BROKER_REALI
            and r.get("dry_run") is False
            and bool(r.get("real_orders_confirmed_at")))
\end{codiceInterno}

\noindent
\cod{dry\_run} deve essere \cod{False} in modo esplicito --- non basta che
manchi --- e la conferma la scrive solo una richiesta che la porta:
\cod{confirm\_real\_orders} quando si salva il runner, \cod{confirm} quando
si cambia modalità. L'avvio di un runner con ordini reali chiede a sua volta
\cod{confirm\_real}. Le rotte \percorsoHTTP{/api/live/*}, rimaste per le
prove da script, su un broker vero girano sempre in simulazione: gli ordini
reali passano solo dai runner, dove ci sono conferma, registro e kill-switch
salvato.

\subsubsection{La conferma è legata a un conto}

Al primo avvio dopo una conferma, \simbolo{\_prepara\_conto} scrive in
\file{data/runners/<id>/account.json} la chiave del conto
(\simbolo{chiave\_conto()}: broker, server per MT5, numero) insieme a
\cod{real\_orders\_confirmed\_at}. Agli avvii successivi con la stessa
conferma un conto diverso ferma l'avvio (\simbolo{\_NonParte}); a ogni ciclo
\simbolo{\_controlla\_conto} confronta di nuovo l'identità e, se è cambiata,
ferma il runner prima del ciclo (\simbolo{\_Fermo}) senza toccare gli
ordini del conto nuovo. Una conferma nuova lega il runner al conto del
prossimo avvio (\simbolo{legame\_valido()}). Senza identità
(\cod{account\_identity()} a \cod{None}), o con \file{account.json}
illeggibile (\simbolo{leggi\_legame()}), gli ordini reali non partono.

\subsection{Il registro del runner}\label{sez:live-registro}

Ogni riga scritta con \simbolo{REGISTRO.log(rid, testo, kind, data)} va in
memoria (le ultime \num{200}, per lo stato) e su disco, in
\file{data/runners/<id>/journal.jsonl}: un oggetto JSON per riga con
\cod{ts} (UTC), \cod{session} (una per avvio), \cod{mode}, \cod{broker},
\cod{kind}, \cod{msg} e, quando serve, \cod{data}. Oltre \num{5}\,MB il
file ruota su \file{journal.1.jsonl} \ldots{} \file{journal.9.jsonl}
(\simbolo{\_ruota\_registro}), e \simbolo{registro\_leggi} li ripercorre
all'indietro finché servono righe. I \cod{kind} sono \cod{mode},
\cod{start}, \cod{stop}, \cod{info}, \cod{cycle}, \cod{order} (simulati),
\cod{fill} (inviati, con l'esito), \cod{warning}, \cod{error}, \cod{risk} e
\cod{maintenance}.

\noindent
Nella stessa cartella stanno lo stato del motore (\file{live\_state.json}),
la storia della normalizzazione (\file{ensemble.json}), lo stato del
\simbolo{RiskManager} (\file{risk.json}) e il legame con il conto
(\file{account.json}). Prima stato e normalizzazione erano un file unico in
\file{data/}: due runner attivi insieme se li sovrascrivevano a vicenda.

\subsection{Dati vecchi}\label{sez:live-dati-vecchi}

\simbolo{LoaderLiveFeed} legge gli stessi file del backtest e non scarica
niente. Se nessuno li aggiorna, il motore continua a decidere sugli stessi
prezzi --- e con ordini veri dimensiona le quantità su prezzi di giorni
prima. \simbolo{live.freshness.motivo\_dati\_vecchi} conta le \emph{sedute
mancanti} fra l'ultima barra e l'ultimo giorno di borsa prima di oggi
(UTC), sul calendario del mercato del simbolo (\simbolo{calendario\_per},
dall'asset class). Prima di oggi, perché il feed legge i file con una data
di fine esclusiva: la barra di oggi la serve domani, e pretenderla dopo la
chiusura americana tratteneva ogni ordine dalle 22 a mezzanotte.

\begin{center}
\begin{tabular}{@{}L{30mm}L{78mm}L{28mm}@{}}
\toprule
\textbf{Calendario} & \textbf{Sedute} & \textbf{Tolleranza} \\
\midrule
\rigaTenue
\cod{nyse} & Lunedì--venerdì meno le festività NYSE (regole di
  \pkg{pandas}: Venerdì Santo, Juneteenth dal 2022, 4 luglio e Natale con
  il giorno osservato) & nessuna \\
\cod{fx} & Lunedì--venerdì & una seduta \\
\rigaTenue
\cod{crypto} & Tutti i giorni & una seduta \\
\bottomrule
\end{tabular}
\end{center}

\noindent
Mancano le chiusure straordinarie (lutti nazionali, uragani): quel giorno i
dati risultano vecchi di una seduta e gli ordini aspettano, cioè l'errore va
nella direzione prudente. Per barre infragiornaliere si contano le barre
mancanti in tempo di borsa (oltre tre, vecchi); per barre settimanali, oltre
due periodi e tre giorni.

\noindent
Il motore misura ogni simbolo a ogni ciclo, anche in simulazione
(\simbolo{\_dati\_vecchi}), e riporta il riassunto in
\cod{CycleReport.data\_stale}. Fuori dal dry-run trattiene gli ordini dei
simboli indietro e lo scrive in \cod{orders\_held}: un'ultima barra fresca
su SPY non rende freschi i prezzi di AAPL. Un simbolo del runner con una
posizione aperta e nessun prezzo nel ciclo (file mancante, simbolo saltato
dal caricamento) finisce nello stesso riassunto: il motore non può né
valutarlo né ridurlo. In entrambi i casi gli ordini che aumentano
l'esposizione aspettano (vedi \emph{Prima le riduzioni}, più avanti). Passano solo le chiusure
ordinate dal kill-switch del conto, che vendono la quantità che c'è e non
dipendono dal prezzo. Il controllo vale per ogni broker tranne quello su
carta (\cod{guard\_stale\_data}), che il simulatore TEST e la prova da
terminale guidano con il proprio orologio.

\subsection{Il kill-switch del conto}\label{sez:live-kill}

Il \simbolo{RiskManager} vive quanto il processo, e il livello del conto
protegge del denaro che sta nel conto. Per questo, con ordini reali, picco
di equity e chiusura sono salvati per \emph{conto}, in
\file{data/accounts/<chiave>.json} (\simbolo{carica\_stato\_conto},
\simbolo{salva\_stato\_conto}); mai nella cartella temporanea, che il
sistema svuota. Il resto dello stato --- strategie sospese con la data,
asset class ridotte con la scadenza, finestre di PnL --- va in
\file{risk.json} del runner con il suo contesto (modalità, broker, conto),
attraverso \simbolo{export\_state()} e \simbolo{restore\_state()}: uno
stato per contesto (\simbolo{contesti\_rischio}), così una corsa in
simulazione non cancella le strategie sospese della modalità reale.

\begin{itemize}
  \item Il drawdown si controlla all'inizio del ciclo, appena letta
        l'equity (\simbolo{check\_account}), non solo dopo gli ordini: un
        calo avvenuto fra due cicli faceva partire un altro giro di acquisti
        prima che il kill-switch scattasse. Con il conto chiuso il ciclo
        procede anche con l'equity a zero: chiudere non ne ha bisogno.
  \item Passare a simulazione e tornare agli ordini reali, o usare un altro
        runner sullo stesso conto, non riapre il conto: le rotte di
        creazione e di cambio modalità non toccano più lo stato. Lo riapre
        solo \rotta{POST}{/api/runners/<rid>/kill-switch/reset}, che azzera
        il file del conto a cui il runner è legato.
  \item Uno stato del conto illeggibile ferma gli ordini reali: trattarlo da
        vuoto riaprirebbe un conto chiuso.
  \item Se all'avvio l'equity è già sotto il picco del conto oltre la
        soglia, il kill-switch scatta e si scrive nel file del conto, ma il
        runner non manda ordini: non può sapere se è un prelievo o una
        perdita. Al prossimo avvio chiude (perdita); il reset lo evita
        (prelievo). Con l'equity non leggibile all'avvio non parte.
  \item Un conto chiuso non impedisce l'avvio: il runner parte in sola
        chiusura (il \simbolo{RiskManager} riparte chiuso, ogni peso è
        zero), chiude e a conto piatto (\simbolo{\_posizioni\_aperte}) si
        ferma; in sola chiusura nessun errore lo ferma. L'elenco dei runner
        segnala la sola chiusura (\cod{close\_only}) solo se il conto chiuso
        è quello della conferma attuale; altrimenti \cod{killed\_unbound}.
  \item Quando il kill scatta durante un ciclo, il ciclo di chiusura parte
        subito, senza aspettare la cadenza; se resta qualcosa di aperto si
        riprova dopo \cod{min(passo, 60\,s)} per tre volte, poi alla
        cadenza normale.
\end{itemize}

\noindent
Il drawdown del conto si controlla a ogni ciclo
(\simbolo{check\_account}); i livelli per strategia e per asset class, che
lavorano su finestre di barre, una volta per barra nuova
(\simbolo{\_nuova\_barra\_rischio}), come nel backtest: aggiornarli a ogni
ciclo --- quattro per barra --- accorciava la finestra di \num{30} giorni a
\num{35} cicli. Il PnL realizzato fra due barre si attribuisce tutto alla
barra nuova (\simbolo{LivePortfolioMirror.consolida\_realizzati}). Il
livello per asset class riceve dal ciclo il PnL di prezzo per classe una
volta per barra (\simbolo{\_registra\_pnl\_asset\_class}). Nel backtest
\simbolo{update} chiama \simbolo{check\_account} con la stessa regola di
prima: i risultati non cambiano.

\subsection{Il ciclo dal vivo}

Il runner fa il primo ciclo subito, poi uno ogni quarto della barra più
corta che guarda (\simbolo{cadenza\_reale\_secondi}, mai sotto i \num{15}
secondi): per strategie giornaliere, ogni sei ore. L'avvio dalla pagina
imponeva \num{60} secondi, e un runner giornaliero rifaceva lo stesso ciclo
ogni minuto.

\noindent
Ogni runner gira in un thread suo, con un proprio event loop
\cod{asyncio}: \pkg{ib\_insync} ne ha bisogno per collegarsi, e un thread
che non è il principale non ne ha uno. Fermato il runner, gli ordini di Lyra
ancora aperti si cancellano (con ordini reali, sul conto verificato), il
collegamento al broker si chiude --- altrimenti IBKR terrebbe occupato il
client ID --- e il loop con lui.

\begin{itemize}
  \item \textbf{Un avvio solo.} \simbolo{start} tiene un lucchetto dai
        controlli fino alla registrazione del runner: con due richieste
        ravvicinate potevano partire due motori sullo stesso conto. Un
        errore dopo la registrazione rimette il runner fermo.
  \item \textbf{Equity assente.} Un conto con equity zero o mancante salta
        il ciclo (\cod{skipped\_reason} \cod{no\_account\_equity}): prima
        ogni peso andava a zero e partiva una chiusura per ogni posizione.
  \item \textbf{Ordini in corso.} Un simbolo con un ordine ancora aperto al
        broker non riceve un ordine nuovo; se gli ordini aperti non si
        leggono, non parte nulla.
  \item \textbf{Prima le riduzioni.} Gli ordini che riducono una posizione
        (\simbolo{\_riduce}) partono prima di quelli che la aumentano, e
        un'inversione diventa chiusura più apertura
        (\simbolo{\_separa\_inversioni}; la gamba di apertura conta come
        aumento). Se qualcosa è trattenuto (dati vecchi, ordine in corso,
        posizione senza prezzo) o una riduzione non è \cod{filled} né
        \cod{skipped}, gli aumenti del ciclo non partono: in una rotazione
        con la vendita trattenuta, l'acquisto da solo portava l'esposizione
        oltre quella del libro vecchio e di quello nuovo.
  \item \textbf{Stop fra un ordine e l'altro.} Il runner assegna al motore
        \cod{stop\_requested}: premuto Stop, gli ordini non ancora inviati
        del ciclo in corso non partono.
  \item \textbf{Cicli falliti.} Un errore di collegamento
        (\simbolo{e\_transitorio}: \simbolo{ConnectionError} e quindi
        \simbolo{BrokerUnavailable}, \simbolo{TimeoutError},
        \simbolo{OSError}) si riprova senza limite, con attese che
        raddoppiano fino a \num{15} minuti. Un'eccezione di altro tipo si
        riprova dopo \cod{min(passo, 60\,s)}; tre di fila fermano il runner
        con la causa, tranne in sola chiusura. Un ciclo saltato non tocca né
        il kill-switch né la manutenzione.
\end{itemize}

\begin{center}
\begin{tabular}{@{}L{50mm}L{88mm}@{}}
\toprule
\textbf{Rotta} & \textbf{Che cosa fa} \\
\midrule
\rigaTenue
\rotta{GET}{/api/runners} & I runner salvati, con \cod{mode},
  \cod{real\_orders}, \cod{killed}, \cod{close\_only}, il conto legato,
  \cod{data\_stale}, \cod{account\_blocked} e le soglie del
  kill-switch. \\
\rotta{POST}{/api/runners} & Crea o aggiorna. Ordini reali solo con
  \cod{dry\_run: false} e \cod{confirm\_real\_orders: true}. \\
\rigaTenue
\rotta{POST}{/api/runners/<rid>/mode} & Simulazione o reale, a runner
  fermo; il reale chiede \cod{confirm}. \\
\rotta{POST}{/api/runners/<rid>/start} & Avvia. Con ordini reali chiede
  \cod{confirm\_real}; senza \cod{interval}, la cadenza viene dal
  timeframe. \\
\rigaTenue
\rotta{POST}{/api/runners/<rid>/stop} & Ferma il runner: gli ordini non
  ancora inviati del ciclo in corso non partono. \\
\rotta{GET}{/api/runners/<rid>/journal} & Le ultime righe del registro
  (\cod{limit}, di default \num{300}). \\
\rigaTenue
\rotta{POST}{/api/runners/<rid>/kill-switch/reset} & Riapre il conto
  chiuso dal kill-switch (il file del conto legato al runner), a runner
  fermo. \\
\rotta{POST}{/api/live/run\_cycle} & Un ciclo solo, per le prove da script.
  Su un broker vero sempre in simulazione. \\
\rigaTenue
\rotta{POST}{/api/live/start} & Ciclo continuo per le prove da script,
  con la stessa regola. \\
\rotta{GET}{/api/live/status} & Stato, numero di cicli, configurazione e le
  ultime \num{60} righe del ciclo di prova. \\
\bottomrule
\end{tabular}
\end{center}

\flushbottom



% --- PARTE III — come si cambia ---------------------------------------------
\part{Come si cambia}
% ===========================================================================
%  sviluppatore/05-estendere/00-capitolo.tex
% ===========================================================================
\chapter{Aggiungere un artefatto}\label{cap:estendere}

% ===========================================================================
%  sviluppatore/05-estendere/01-due-forme.tex
% ===========================================================================
\section{Non c'è un contratto solo: ce ne sono due}\label{sez:due-forme}

È la prima cosa da sapere prima di scrivere un plugin, e non è scritta da
nessuna parte nel codice: le cinque famiglie di artefatti si dividono in due
forme, con contratti, meccanismi di scoperta e modi di fallire diversi.

\figuraTikz{plugin-due-forme}
  {Le due forme. Scegliere la famiglia significa scegliere la forma, e con
   essa il modo in cui l'artefatto può non comparire.}
  {plugin-due-forme}

\begin{center}
\begin{tabular}{@{}L{22mm}L{34mm}L{36mm}L{40mm}@{}}
\toprule
 & \textbf{Famiglie} & \textbf{Oggetto esposto} & \textbf{Punto d'ingresso} \\
\midrule
\rigaTenue
\textbf{Forma A} & \pkgStrategie, \pkgApprocci
  & \cod{DESCRIPTOR},\newline\cod{APPROACH\_DESCRIPTOR}
  & \cod{factory(params)},\newline\cod{build(params)} \\
\textbf{Forma B} & \pkgPesi, \pkgOverlay, \pkgRegime
  & \cod{META} (dict di soli letterali)
  & \cod{compute(...)} \\
\bottomrule
\end{tabular}
\end{center}

\subsection{Come si differenziano i due registri}

\begin{center}
\begin{tabular}{@{}L{34mm}L{50mm}L{52mm}@{}}
\toprule
 & \textbf{Forma A} & \textbf{Forma B} \\
\midrule
\rigaTenue
Quando si costruisce & A richiesta: \cod{discover\_*()} restituisce un dict
  nuovo ogni volta & All'import del package, una volta sola \\
Come trova i file & \cod{importlib}, cartella per cartella & \cod{pkgutil},
  più un ripiego sulle cartelle senza \file{\_\_init\_\_.py} \\
\rigaTenue
File atteso & \file{strategy.py}, \file{approach.py} &
  \file{weight\_mode.py}, \file{overlay.py}, o un modulo \file{.py} \\
Se un plugin è rotto & Avviso nel registro, saltato & Avviso, saltato \\
\rigaTenue
Se un plugin è \emph{scritto male} & Spesso un errore all'uso &
  \textbf{Invisibile}, senza un rigo da nessuna parte \\
Per vederne uno nuovo & Basta richiamare \cod{discover\_*()} &
  Serve \cod{reload\_registry()}, o un riavvio \\
\bottomrule
\end{tabular}
\end{center}

\begin{attenzione}{Il modo tipico di fallire è l'invisibilità}
Vale per tutta la forma~B: un \cod{META} che non è fatto di soli letterali,
una \cod{key} che non coincide con il nome del file, un
\cod{APPROACH\_DESCRIPTOR} messo in \pkgOverlay, un \cod{compute} mancante.
Nessuno di questi solleva un'eccezione. Il plugin semplicemente non appare
nell'elenco, e non c'è niente da cercare nel registro perché non è successo
niente.

Chi scrive un artefatto e non lo vede comparire non ha un bug da trovare: ha
un contratto da rileggere. Il \vediSez{trappole} elenca i casi uno per uno.
\end{attenzione}

\subsection{Il nome canonico è il nome della cartella}

In entrambe le forme, e vale la pena insistere: il nome con cui il sistema
conosce un artefatto è il nome della sua cartella, non quello dichiarato
dentro il file. Il registro delle strategie arriva a sovrascriverlo:

\begin{codiceInterno}{python}{Il nome della cartella vince}{strategies/registry.py:106}
        if descriptor.name != name:
            # il nome CARTELLA vince sempre
            descriptor.name = name
\end{codiceInterno}

\noindent
Per la forma~B la regola è ancora più stretta: la \cod{key} del \cod{META}
\emph{deve} coincidere con il nome del file, e per le regole di regime il
controllo è esplicito. Negli altri casi il disallineamento non dà errore:
l'artefatto compare con un nome e si cerca con un altro.

% ===========================================================================
%  sviluppatore/05-estendere/02-strategia.tex
% ===========================================================================
\section{Una strategia}\label{sez:nuova-strategia}

\subsection{Che cosa si crea}

\begin{codiceInterno}{console}{La struttura minima}{strategies/}
strategies/mia_strategia/
|-- __init__.py       # ri-esporta DESCRIPTOR
|-- strategy.py       # espone DESCRIPTOR
`-- mia_impl.py       # l'implementazione
\end{codiceInterno}

\noindent
La cartella \textbf{non} deve cominciare con \cod{\_}: il prefisso è
l'interruttore di spegnimento, e il filtro che lo applica è questo.

\begin{codiceInterno}{python}{Il filtro delle cartelle}{strategies/registry.py:55}
def _is_valid_strategy_dir(p: Path) -> bool:
    if not p.is_dir():                                   return False
    if p.name.startswith("_") or p.name.startswith("."): return False
    if p.name in SKIP_DIR_NAMES:                         return False
    if not (p / "strategy.py").exists():                 return False
    return True
\end{codiceInterno}

\subsection{Il livello di esecuzione}

Si sottoclassa \simbolo{BaseAtomicStrategy} da \file{strategies/base.py}, che
offre già \simbolo{applicable\_symbols()}, \simbolo{get\_value()} e
\simbolo{get\_last\_value\_until()}. Tre cose sono obbligatorie:

\begin{codiceInterno}{python}{Il contratto di esecuzione}{strategies/base.py:55}
class AtomicStrategy(ABC):
    @property
    @abstractmethod
    def style(self) -> StrategyStyle: ...

    @property
    @abstractmethod
    def allowed_asset_classes(self) -> set[AssetClass]: ...

    @abstractmethod
    def generate_signals(
        self,
        t_decision: datetime,
        t_execution: datetime,
        indicators_by_symbol: dict[str, PrecomputedIndicators],
        portfolio: Portfolio,
    ) -> list[Signal]:
\end{codiceInterno}

\noindent
I valori ammessi per \cod{style} vengono da \simbolo{core.types.StrategyStyle}:
\cod{TREND}, \cod{MEAN\_REVERSION}, \cod{MOMENTUM}, \cod{BREAKOUT},
\cod{STAT\_ARB}, \cod{STAT\_ARB\_INTRADAY}, \cod{FACTOR}, \cod{VALUE},
\cod{MACRO}, \cod{CARRY}, \cod{VOL}.

\begin{attenzione}{I nomi dei parametri fanno parte del contratto}
Il motore chiama \cod{generate\_signals} per \emph{keyword}
(\vediSez{decisione-esecuzione}). Rinominare \cod{t\_decision} in \cod{t}, o
\cod{indicators\_by\_symbol} in \cod{ind}, fa fermare il backtest alla prima
barra con un \cod{TypeError}.
\end{attenzione}

\subsection{Il livello dei metadati}

\simbolo{StrategyDescriptor} ha sei campi obbligatori e una ventina di
opzionali.

\begin{codiceInterno}{python}{I sei obbligatori}{strategies/descriptor.py:107}
    # Required
    name: str
    description: str
    category: StrategyCategory
    timeframe: StrategyTimeframe
    factory: Callable[[dict], Any]
    default_params: dict
\end{codiceInterno}

\noindent
Un esempio reale, che mostra anche il modo idiomatico di scrivere la
\cod{factory}:

\begin{codiceInterno}{python}{Il descrittore di una strategia installata}{strategies/momentum/strategy.py:16}
DESCRIPTOR = StrategyDescriptor(
    name="momentum",
    description="Cross-Sectional Momentum 12-1 (Jegadeesh-Titman). "
                "Top decile, monthly rebalance, long-only.",
    category=StrategyCategory.D1_LONG_ONLY,
    timeframe=StrategyTimeframe.D1,

    excluded_instruments=["EURUSD", "GBPUSD", "USDJPY", "USDCHF", "AUDUSD"],
    required_indicators=["momentum_252_21"],

    factory=lambda params: CrossSectionalMomentumStrategy(**{
        k: v for k, v in params.items() if not k.startswith("exit_")
    }),
    default_params={
        "formation_window": 252,
        "skip_window": 21,
        "top_decile_pct": 0.10,
        "max_positions": 25,
        "rebalance_every_n_days": 21,
        "exit_catastrophic_sl_pct": None,
    },

    min_lookback_days=300,
    backtest_script="strategies/momentum/run_backtest.py",
    output_path_pattern="output/strategies/wf_momentum_oos_equity*.parquet",
    stream_type="equity_curve",
    stream_category="d1",
    exit_policy_type="signal_only",
    regime_weights={"bull": 1.0, "neutral": 1.0, "bear": 0.7},
)
\end{codiceInterno}

\begin{nota}{Il filtro delle chiavi \cod{exit\_}}
La \cod{factory} scarta le chiavi che cominciano con \cod{exit\_} prima di
passarle al costruttore. Quei parametri non appartengono alla strategia: li
consuma la politica di uscita, costruita a parte da
\cod{exit\_policy\_type}. Scrivere la \cod{factory} senza quel filtro
produce un \cod{TypeError} su un argomento inatteso.
\end{nota}

\subsection{I campi opzionali che contano davvero}

\begin{center}
\begin{tabular}{@{}L{40mm}L{98mm}@{}}
\toprule
\textbf{Campo} & \textbf{A che serve} \\
\midrule
\rigaTenue
\cod{output\_path\_pattern},\newline\cod{stream\_type},
\cod{stream\_category}
  & \textbf{Integrazione con la dashboard, senza toccarla.} Con questi tre
    campi lo stream compare da solo in \voce{Streams}: non c'è nessun
    elenco da aggiornare in \file{dashboard/config.py}. Senza
    \cod{output\_path\_pattern} lo stream non compare. \\
\cod{required\_indicators}
  & Nomi presi dal pool comune. Il registro ne calcola l'unione su tutte le
    strategie e li pre-calcola una volta sola. \\
\rigaTenue
\cod{regime\_weights}
  & Modulatori \cod{bull}/\cod{neutral}/\cod{bear}. Le chiavi devono essere
    esattamente quelle tre e i valori numerici non negativi: è l'unico
    controllo del descrittore che solleva \cod{ValueError} invece di un
    avviso. \\
\cod{allocator\_overrides}
  & Scavalca per questa sola strategia \cod{max\_position\_weight},
    \cod{max\_gross\_leverage}, \cod{max\_asset\_class\_gross},
    \cod{portfolio\_target\_vol\_annual}. \\
\rigaTenue
\cod{extra\_timeframes}
  & Per le strategie multi-timeframe. Il caricamento parte dal più fine. \\
\cod{required\_external\_data}
  & \cod{news}, \cod{fundamentals}, \cod{policy\_rates}. Il sistema li
    \emph{scarica} se mancano. \\
\rigaTenue
\cod{required\_mms\_features},\newline\cod{required\_volume\_profile}
  & Feature specialistiche calcolate su richiesta. \\
\cod{maintenance}
  & \cod{retrain\_every\_days}, \cod{retrain\_every\_bars},
    \cod{retrain\_on\_drawdown\_pct}, \cod{data\_refresh\_every\_days}. È
    da qui che il sistema sa quando un modello va riaddestrato. \\
\rigaTenue
\cod{\_module\_path}, \cod{\_folder\_path}
  & \textbf{Li scrive il registro.} Non vanno impostati a mano. \\
\bottomrule
\end{tabular}
\end{center}

\begin{attenzione}{\cod{regime\_weights} si dichiara a priori}
Il commento nel descrittore è esplicito: modificarli guardando il backtest
«sarebbe curve-fitting mascherato da regime adaptation». Si dichiarano sul
ragionamento strutturale --- questa strategia compra il momentum, quindi in
regime ribassista deve pesare meno --- e poi si verifica che il modulato non
peggiori rispetto al neutro. Non si cercano i valori che massimizzano il
risultato.
\end{attenzione}

\subsection{La procedura}

\begin{pratica}{Aggiungere una strategia}
\begin{enumerate}
  \item Creare la cartella \textbf{senza} prefisso \cod{\_}.
  \item \file{mia\_impl.py}: sottoclassare \simbolo{BaseAtomicStrategy},
        implementare \cod{style}, \cod{allowed\_asset\_classes},
        \cod{reset()} e \cod{generate\_signals()} \emph{con quei nomi esatti
        di parametro}.
  \item Leggere solo fino a \cod{t\_decision}. Mettere
        \cod{Signal.timestamp = t\_execution}.
  \item \file{strategy.py}: \cod{DESCRIPTOR = StrategyDescriptor(...)} con i
        sei campi obbligatori; \cod{factory} che filtra le chiavi
        \cod{exit\_}.
  \item Per vedere lo stream in dashboard: \cod{output\_path\_pattern} +
        \cod{stream\_type} + \cod{stream\_category}.
  \item \file{\_\_init\_\_.py}: ri-esportare \cod{DESCRIPTOR}.
  \item Verificare con \file{tests/diagnostica\_strategie.py}, che esiste
        proprio per rispondere alla domanda «perché
        \cod{discover\_strategies()} torna vuota».
\end{enumerate}
\end{pratica}

\noindent
Il riferimento completo è \file{strategies/GUIDE\_STRATEGIE\_CUSTOM.md}
(\num{2281} righe), affiancato da sei guide tematiche: il grafo a nodi come
alternativa al Python, l'uso dei regimi dentro una strategia, le posizioni
multiple sullo stesso simbolo, le feature MMS, il volume profile e un
documento di proposta sulle strategie in valuta.

% ===========================================================================
%  sviluppatore/05-estendere/03-approccio.tex
% ===========================================================================
\section{Un approccio}\label{sez:nuovo-approccio}

Un approccio \emph{scavalca il motore}. Non implementa
\simbolo{AtomicStrategy}, non emette \simbolo{Signal}, non passa per
allocatore, rischio ed esecuzione. Costruisce il proprio calcolo e consegna
direttamente la serie dei rendimenti.

\begin{center}
\begin{tabular}{@{}L{28mm}L{52mm}L{54mm}@{}}
\toprule
 & \textbf{Strategia} & \textbf{Approccio} \\
\midrule
\rigaTenue
File & \file{strategy.py} & \file{approach.py} \\
Oggetto & \cod{DESCRIPTOR} & \cod{APPROACH\_DESCRIPTOR} \\
\rigaTenue
Entrata & \cod{factory(params)} \freccia{} istanza & \cod{build(params)}
  \freccia{} \cod{pd.Series} \\
Chi esegue & il motore di backtest & l'approccio stesso \\
\rigaTenue
Produce & segnali, poi posizioni & rendimenti giornalieri \\
Runner & \file{execution/run\_walkforward.py} &
  \file{execution/run\_approach.py} \\
\bottomrule
\end{tabular}
\end{center}

\subsection{Il contratto}

\begin{codiceInterno}{python}{Dichiarato nel docstring del modulo}{approaches/descriptor.py:17}
CONTRATTO della funzione build:
    build(params: dict) -> pandas.Series
        - index = date (tz-naive); valori = rendimenti DAILY (frazione)
        - usa SOLO dati fino a t per ogni t (niente lookahead)
\end{codiceInterno}

\noindent
«Tz-naive» non è pignoleria: una serie con fuso orario e una senza non si
allineano, e il risultato è un ensemble che perde silenziosamente le
osservazioni di uno stream. Per questo gli approcci installati chiudono tutti
allo stesso modo:

\begin{codiceInterno}{python}{La chiusura obbligata di build}{approaches/momentum\_ema\_approach/approach.py:74}
    # Assicuriamo che l'indice sia tz-naive come richiesto
    pnl.index = pd.to_datetime(pnl.index).tz_localize(None)
    return pnl
\end{codiceInterno}

\subsection{Il descrittore}

Più corto di quello delle strategie, e con tre campi che si derivano da soli
se li si omette:

\begin{codiceInterno}{python}{I default derivati}{approaches/descriptor.py:79}
        if not self.label:
            self.label = self.name
        if not self.output_path_pattern:
            self.output_path_pattern = f"output/strategies/{self.name}_pnl_daily.parquet"
        if not self.backtest_script:
            self.backtest_script = f"approaches/{self.name}/run_backtest.py"
\end{codiceInterno}

\noindent
Due campi hanno invece una lista chiusa di valori ammessi, e un valore fuori
lista solleva \cod{ValueError} al caricamento:

\begin{codiceInterno}{python}{Le due allowlist}{approaches/descriptor.py:31}
APPROACH_STREAM_TYPES = ("pnl_daily", "pnl_sparse", "equity_curve")
APPROACH_STREAM_CATEGORIES = ("ml", "d1", "intraday", "ml_overlay")
\end{codiceInterno}

\noindent
Un esempio completo:

\begin{codiceInterno}{python}{Un descrittore di approccio}{approaches/momentum\_ema\_approach/approach.py:79}
APPROACH_DESCRIPTOR = ApproachDescriptor(
    name=NAME,
    label="EMA Momentum Approach",
    description=("Generates weights from the distance between the closing "
                 "price and its exponential moving average (EMA)."),
    build=build,
    default_params={"lookback_ema": 12, "short_lookback": 20},
    stream_type="pnl_daily",
    stream_category="d1",
    tags=["momentum", "ema", "d1"],
)
\end{codiceInterno}

\subsection{Approcci e overlay: lo stesso registro, due elenchi}

Il registro degli approcci serve \emph{due} cartelle, \pkgApprocci{} e
\pkgOverlay, e il package si deduce da dove viene la cartella:

\begin{codiceInterno}{python}{Il package dipende dalla provenienza}{approaches/registry.py:56}
    # il package dipende dalla cartella di provenienza: approaches/ oppure overlays/
    pkg = approach_dir.parent.name if approach_dir.parent.name in ("approaches", "overlays") else "approaches"
    module_dotted = f"{pkg}.{name}.approach"
\end{codiceInterno}

\noindent
Di conseguenza ci sono tre funzioni pubbliche con semantiche distinte:

\begin{center}
\begin{tabular}{@{}L{32mm}L{106mm}@{}}
\toprule
\textbf{Funzione} & \textbf{Che cosa restituisce, e perché} \\
\midrule
\rigaTenue
\cod{discover\_approaches()} & Solo \pkgApprocci. Sono quelli eseguibili dal
  pulsante \pulsante{Run backtest}. \\
\cod{discover\_overlays()} & Solo le cartelle \file{overlays/*/approach.py}.
  Non sono eseguibili da sole. \\
\rigaTenue
\cod{discover\_all()} & L'unione, \textbf{solo per risolvere un nome}, mai
  per costruire un elenco da mostrare. \\
\bottomrule
\end{tabular}
\end{center}

\begin{perche}{Perché gli elenchi sono separati}
Il razionale è scritto nel registro: mettere gli overlay fra i backtest
lanciabili «li faceva comparire fra i backtest lanciabili, dove sarebbero
falliti o --- peggio --- avrebbero prodotto uno stream vuoto senza dirlo». Uno
stream vuoto che non si lamenta è peggio di un errore, perché entra
nell'ensemble e lo cambia.
\end{perche}

\begin{nota}{Il campo \cod{family} non esiste nella dataclass}
\cod{\_scansiona} attacca all'oggetto un attributo \cod{family} che non è
dichiarato in \simbolo{ApproachDescriptor} e non compare in \cod{to\_dict()}:

\begin{codiceInterno}{python}{}{approaches/registry.py:113}
        # Da dove viene: e' l'unico modo affidabile per distinguere un
        # approccio da un overlay, perche' il contratto del file e' identico.
        try:
            desc.family = etichetta
        except Exception:      # dataclass congelata: non e' un problema
            pass
\end{codiceInterno}

\noindent
Funziona perché la dataclass non è congelata. Chi legge
\file{descriptor.py} non sa che esiste.
\end{nota}

\begin{pratica}{Aggiungere un approccio}
\begin{enumerate}
  \item Cartella senza \cod{\_}, con \file{approach.py} e
        \file{\_\_init\_\_.py}.
  \item \cod{def build(params: dict) -> pd.Series}: indice di date tz-naive,
        valori di rendimento giornaliero in frazione, causale.
  \item \cod{APPROACH\_DESCRIPTOR} con almeno \cod{name}, \cod{description},
        \cod{build}.
  \item Controllare che \cod{stream\_type} e \cod{stream\_category} stiano
        nelle due allowlist.
  \item Provarlo: \file{python -m execution.run\_approach <nome>}.
\end{enumerate}
\end{pratica}

% ===========================================================================
%  sviluppatore/05-estendere/04-forma-b.tex
% ===========================================================================
\section{Una modalità di pesatura, o un overlay}\label{sez:forma-b}

Le due famiglie condividono forma, meccanismo di scoperta e modi di fallire.
Cambiano il file atteso e la firma di \cod{compute}.

\begin{center}
\begin{tabular}{@{}L{22mm}L{34mm}L{80mm}@{}}
\toprule
 & \textbf{File} & \textbf{Firma e significato del ritorno} \\
\midrule
\rigaTenue
\pkgPesi & \file{weight\_mode.py}
  & \cod{compute(eng, stream\_pnls)}\newline
    \cod{kind="static"} \freccia{} \cod{dict[str, float]}, somma \num{1}\newline
    \cod{kind="rolling"} \freccia{} \cod{DataFrame} di pesi nel tempo \\
\pkgOverlay & \file{overlay.py}
  & \cod{compute(eng, ensemble\_pnl, stream\_pnls)}\newline
    \freccia{} \cod{Series} di \emph{moltiplicatori} in $[0,1]$, allineata
    a \cod{ensemble\_pnl.index} \\
\bottomrule
\end{tabular}
\end{center}

\subsection{Due forme di cartella, e perché}

Entrambe le famiglie accettano sia un modulo singolo
(\file{weights/<chiave>.py}, forma storica) sia una cartella
(\file{weights/<chiave>/weight\_mode.py}, da preferire per i nuovi). La
scoperta ha un ripiego che vale la pena conoscere:

\begin{codiceInterno}{python}{Il ripiego sulle cartelle senza \_\_init\_\_.py}{weights/\_\_init\_\_.py:26}
def _nomi_disponibili():
    """I nomi da provare: moduli, package e CARTELLE SENZA __init__.py.

    `pkgutil.iter_modules` non vede una cartella priva di `__init__.py`. Il
    marketplace pero' esclude dai pacchetti i file che iniziano con "__",
    quindi un artefatto installato arriva proprio cosi': con il solo
    weight_mode.py. Senza questa aggiunta un overlay comprato non si sarebbe
    caricato, e l'unico sintomo sarebbe stato la sua assenza dall'elenco.
    """
\end{codiceInterno}

\noindent
È il tipo di dettaglio che non si deduce leggendo il meccanismo: il
marketplace e il caricatore dei plugin hanno un'interazione, e il caricatore
si adatta a come il marketplace impacchetta.

\subsection{Il registro si costruisce all'import}

\begin{codiceInterno}{python}{La costruzione del registro}{weights/\_\_init\_\_.py:73}
for _nome in _nomi_disponibili():
    try:
        _m = importlib.import_module(f"{__name__}.{_nome}")
    except Exception as exc:  # un plugin rotto non deve impedire l'avvio
        logger.debug("%s: import diretto fallito (%s)", _nome, exc)
        _m = None
    _m = _risolvi(_m, _nome)
    if _m is not None and hasattr(_m, "META") and hasattr(_m, "compute"):
        _mods.append(_m)

_mods.sort(key=lambda m: (m.META.get("order", 999), m.META["key"]))

REGISTRY  = tuple(m.META for m in _mods)
COMPUTERS = {m.META["key"]: m.compute for m in _mods}
\end{codiceInterno}

\noindent
La riga che conta è l'ultima condizione: senza \cod{META} \emph{e}
\cod{compute} il modulo viene scartato in silenzio. Non c'è nemmeno un
\cod{logger.warning}: l'oggetto semplicemente non entra nel registro.

\subsection{Il campo \cod{META}}

\begin{center}
\begin{tabular}{@{}L{34mm}L{16mm}L{86mm}@{}}
\toprule
\textbf{Campo} & \textbf{Obbl.} & \textbf{A che serve} \\
\midrule
\rigaTenue
\cod{key} & sì & Identificativo \cod{[a-z0-9\_]}, uguale al nome del file o
  della cartella. È la chiave del dizionario dei \cod{compute}. \\
\cod{label} & sì & Il nome mostrato nell'interfaccia. \\
\rigaTenue
\cod{kind} & solo \pkgPesi & \cod{"static"} o \cod{"rolling"}. Governa il
  dispatch. \\
\cod{group} & no & Raggruppamento nell'interfaccia. \\
\rigaTenue
\cod{order} & no & Ordinamento, default \num{999}. \\
\cod{hint} & no & Una riga di descrizione. \\
\rigaTenue
\cod{maintenance} & no & \cod{retrain\_every\_days},
  \cod{data\_refresh\_every\_days}, \cod{notes}. \\
\cod{required\_external\_data} & no & \cod{news}, \cod{macro},
  \cod{fundamentals}, \cod{rates}. \\
\bottomrule
\end{tabular}
\end{center}

\noindent
Un esempio minimo e completo:

\begin{codiceInterno}{python}{La modalità più semplice che esista}{weights/equal/weight\_mode.py:16}
META = {
    "key": 'equal',
    "label": 'Equal',
    "kind": 'static',
    "group": 'Naive',     # nel sorgente: Naive con la dieresi
    "order": 10,
    "hint": "Classic 1/N",
}

def compute(eng, stream_pnls: dict[str, pd.Series]) -> dict[str, float]:
    """Pesi equal-weight 1/N. Output: dict {stream_name: weight}, somma 1.0."""
    n = len(stream_pnls)
    if n == 0:
        return {}
    return {name: 1.0 / n for name in stream_pnls}
\end{codiceInterno}

\noindent
E uno che dichiara dati esterni e una scadenza:

\begin{codiceInterno}{python}{Quando l'artefatto ha bisogno di dati freschi}{overlays/news\_pressure/overlay.py:30}
    # QUALI dati servono e OGNI QUANTO vanno rinfrescati. Il sistema legge di
    # qui per sapere quando chiamare la porta: senza, la scadenza non
    # esisterebbe e in live i dati resterebbero fermi al giorno
    # dell'installazione.
    "required_external_data": ["news"],
    "maintenance": {
        "data_refresh_every_days": 1,
        "notes": "news goes stale in hours: it is refreshed every day",
    },
\end{codiceInterno}

\subsection{Che cosa si può importare}

Il codice di un artefatto di forma~B passa da un validatore \simbolo{AST}
quando arriva dall'editor. Le radici ammesse e quelle vietate:

\begin{center}
\begin{tabular}{@{}L{30mm}L{108mm}@{}}
\toprule
 & \textbf{Moduli} \\
\midrule
\rigaTenue
Ammessi ovunque & \pkg{numpy}, \pkg{pandas}, \pkg{math}, \pkg{statistics},
  \pkg{logging}, \pkg{typing} \\
Ammessi in \pkgPesi & più \pkg{scipy}, \pkg{sklearn}, \pkg{lightgbm},
  \cod{from ..metrics import ...} \\
\rigaTenue
Ammessi in \pkgOverlay & più \pkg{approaches}, \pkg{core}, \pkg{reporting} \\
\textbf{Vietati} & \pkg{os}, \pkg{sys}, \pkg{subprocess}, \pkg{socket},
  \pkg{requests}, \cod{open}, \cod{eval}, \cod{exec}, \cod{\_\_import\_\_},
  \pkg{importlib} \\
\bottomrule
\end{tabular}
\end{center}

\begin{perche}{Perché un validatore e non la fiducia}
Un artefatto del marketplace è codice di qualcun altro che gira sulla
macchina di chi lo installa, con i suoi dati e le sue credenziali del broker
a portata di \cod{open()}. Il validatore non rende sicuro l'ineseguibile: un
calcolo malizioso scritto con solo \pkg{numpy} resta possibile. Toglie però
la via diretta --- leggere un file, aprire una connessione, eseguire una
stringa --- e rende il resto visibile a chi guarda il sorgente, che il
marketplace mostra sempre.
\end{perche}

\subsection{Le regole dei pesi}

Oltre al contratto di firma, una modalità di pesatura deve rispettare quattro
condizioni, che il motore \emph{non} verifica:

\begin{itemize}
  \item pesi \textbf{solo lunghi};
  \item limitati all'intervallo \cod{[eng.min\_weight, eng.max\_weight]};
  \item normalizzati a somma \num{1};
  \item \cod{\{\}} se non calcolabili --- l'ensemble ha un ripiego, ma solo
        se glielo si dice restituendo il vuoto. Restituire pesi sbagliati non
        attiva nessun ripiego.
\end{itemize}

\subsection{Le regole degli overlay}

\begin{itemize}
  \item il ritorno è una serie di \emph{moltiplicatori}, non di rendimenti;
  \item causale: \cod{shift(1)} sulla decisione --- «la decisione di oggi non
        può usare il rendimento di oggi»;
  \item \textbf{non} esporre \cod{APPROACH\_DESCRIPTOR}. Il registro cerca
        \cod{META} e \cod{compute}; non trovandoli, l'artefatto resta
        invisibile senza alcun errore.
\end{itemize}

\begin{pratica}{Aggiungere una modalità di pesatura o un overlay}
\begin{enumerate}
  \item \file{weights/mia\_chiave/weight\_mode.py} (o
        \file{overlays/mia\_chiave/overlay.py}), con
        \file{\_\_init\_\_.py} che può restare vuoto.
  \item \cod{META} con \cod{key} uguale al nome della cartella.
  \item \cod{compute} con la firma giusta per la famiglia.
  \item Se il processo è già avviato:
        \cod{from weights import reload\_registry; reload\_registry()}.
\end{enumerate}
I modelli di riferimento sono \file{weights/WEIGHT\_MODE\_TEMPLATE.md} e
\file{overlays/OVERLAY\_TEMPLATE.md}; quest'ultimo si apre con una tabella
che separa overlay e approccio su sei assi, e la frase è «\emph{An overlay and
an approach are two different constructs. This is not a nuance.}»
\end{pratica}

% ===========================================================================
%  sviluppatore/05-estendere/05-regime.tex
% ===========================================================================
\section{Un regime: due file, in quest'ordine}\label{sez:nuovo-regime}

\pkgRegime{} è l'unica famiglia divisa in due artefatti distinti, ed è la
divisione giusta: \emph{una regola serve più mercati}. Azioni e cripto usano
entrambe \cod{trend\_vol}, con parametri diversi.

\figuraTikz{catena-regime}
  {Dal simbolo allo stato. I due riquadri con il file sotto sono le due cose
   che si scrivono.}
  {catena-regime}

\subsection{La regola}

\begin{codiceInterno}{python}{Il contratto della regola}{regime/logics/\_\_init\_\_.py:7}
    META = {
        "key": "mia_regola",
        "label": "La mia regola",
        "description": "a cosa serve",
        "defaults": {"finestra": 60},
    }

    def compute(close: pd.Series, p: dict) -> pd.Series:
        ...   # Series di 'bull' / 'neutral' / 'bear'

Due vincoli non negoziabili:

1. **Causalita'**: il valore al tempo t usa solo dati fino a t. Niente
   statistiche sull'intero campione (quantili full-sample = look-ahead).
2. **Etichette**: sempre ``bull``/``neutral``/``bear``, cosi' i
   ``regime_weights`` gia' dichiarati dalle strategie continuano a valere.
\end{codiceInterno}

\noindent
Gli helper causali condivisi stanno in \file{regime/logics/\_base.py}:
\simbolo{\_realized\_vol}, \simbolo{\_expanding\_median},
\simbolo{\_hysteresis}. Vanno usati: riscrivere un quantile a mano è il modo
più comune di introdurre una statistica full-sample senza volerlo.

Le tre regole installate:

\begin{center}
\begin{tabular}{@{}L{26mm}L{54mm}L{58mm}@{}}
\toprule
\textbf{Chiave} & \textbf{Che cosa misura} & \textbf{Che cosa significano le
  etichette} \\
\midrule
\rigaTenue
\cod{trend\_vol} & Posizione rispetto alla media mobile lunga, corretta per
  volatilità realizzata & bull = sopra il trend e volatilità accettabile;
  bear = sotto il trend, o volatilità alta \\
\cod{vol\_state} & Volatilità realizzata contro i quantili espansivi della
  propria storia & bull = calmo; bear = sotto stress. Per attività senza
  deriva \\
\rigaTenue
\cod{trend\_range} & Rapporto di varianza sui rendimenti: la persistenza, non
  la direzione & bull = in tendenza; bear = laterale. Indipendente dalla
  deriva \\
\bottomrule
\end{tabular}
\end{center}

\subsection{Il mercato}

\begin{codiceInterno}{python}{Il contratto del mercato}{regime/markets/\_\_init\_\_.py:12}
    META = {
        "key": "equity",
        "label": "Azioni",
        "logic": "trend_vol",          # regola ereditata da regime/logics.py
        "driver": ["SPY"],             # serie che GUIDA il regime (in ordine)
        "match": [r"^[A-Z]{1,5}$"],    # simboli assegnati a questo mercato
        "priority": 50,                # match piu' bassa = valutata prima
        "enabled": True,
        "params": {"trend_window": 200},   # override sui default della logica
    }
\end{codiceInterno}

\noindent
Un esempio reale, che mostra come si usa \cod{priority} e che cosa
significa \cod{"self"} fra i driver:

\begin{codiceInterno}{python}{Il mercato valutato per primo}{regime/markets/forex.py:17}
META = {
    "key": "forex",
    "label": "Forex",
    "logic": "vol_state",
    "driver": ["DXY", "self"],
    "match": [r"^(EUR|USD|GBP|JPY|CHF|AUD|NZD|CAD|SEK|NOK)"
              r"(EUR|USD|GBP|JPY|CHF|AUD|NZD|CAD|SEK|NOK)$",
              r"^[A-Z]{3}/[A-Z]{3}$",
              r"^DXY$"],
    "priority": 10,          # valutato per primo: pattern molto specifico
    "enabled": True,
    "params": {
        "vol_window": 21,
        "min_obs": 252,
        # banda ampia + isteresi lunga: un regime deve essere uno STATO, non
        # una reazione. Valori strutturali, non tarati sui dati.
        "calm_quantile": 0.30,      # sotto questo quantile: risk-on
        "stress_quantile": 0.80,    # sopra questo quantile: risk-off
        "hysteresis_days": 10,
    },
}
\end{codiceInterno}

\begin{nota}{\cod{driver} è una lista di ripieghi}
Si valutano in ordine: il primo disponibile vince. \cod{"self"} significa «il
simbolo stesso». Se nessuno dei driver è disponibile il mercato non modula
--- cioè si comporta come neutro --- e lo dichiara in un avviso nel registro,
non in un errore. Un mercato che non modula non si distingue da un mercato
che modula con pesi \num{1}, se non leggendo il registro.
\end{nota}

\subsection{\cod{META} di soli letterali: qui è un vincolo vero}

Per le regole di regime il sistema legge il \cod{META} con
\simbolo{ast.literal\_eval}, cioè \textbf{senza eseguire il file}. La
conseguenza è netta: un valore come \cod{"x": np.pi} non è un letterale, la
lettura fallisce, e la regola diventa invisibile.

\begin{attenzione}{La \cod{key} deve coincidere con il nome del file}
Per le regole di regime il controllo è esplicito ed è l'unico caso in cui il
disallineamento viene rifiutato. Negli altri casi non lo è, e il plugin
compare con un nome e si cerca con un altro.
\end{attenzione}

\subsection{Il merge dei parametri}

I \cod{params} dichiarati dal mercato scavalcano i \cod{defaults} dichiarati
dalla regola, e il punto in cui accade è uno solo:

\begin{codiceInterno}{python}{Dove i due artefatti si incontrano}{regime/logics/\_\_init\_\_.py:89}
def apply(logic_key: str, close, params: dict | None = None):
    if logic_key not in REGISTRY:
        raise KeyError(f"unknown regime rule: '{logic_key}' "
                       f"(available: {sorted(REGISTRY)})")
    entry = REGISTRY[logic_key]
    p = dict(entry["defaults"])
    p.update(params or {})
    return entry["fn"](close, p)
\end{codiceInterno}

\begin{pratica}{Aggiungere un regime}
\begin{enumerate}
  \item \file{regime/logics/mia\_regola.py}: \cod{META} di soli letterali con
        \cod{key} uguale al nome del file, e \cod{compute(close, p)} che
        restituisce \cod{'bull'}/\cod{'neutral'}/\cod{'bear'} allineati a
        \cod{close}, causali, con isteresi.
  \item \file{regime/markets/mio\_mercato.py}: \cod{META} con almeno
        \cod{key} e \cod{logic}; la regola nominata deve \emph{già} esistere.
  \item Scaricare le serie dei driver:
        \file{python scripts/download\_regime\_drivers.py}.
  \item Scorciatoia: \cod{regime.logics.scaffold(key)} genera lo scheletro
        corretto, ed è quello che usa l'editor dell'interfaccia.
\end{enumerate}
Le guide sono \file{regime/GUIDE\_REGIMI.md} («come costruirne uno») e
\file{strategies/GUIDE\_REGIME.md} («come usarlo dentro una strategia»): sono
due documenti distinti e si rimandano a vicenda.
\end{pratica}

% ===========================================================================
%  sviluppatore/05-estendere/06-trappole.tex
% ===========================================================================
\section{«L'ho scritto e non compare»}\label{sez:trappole}

Questa sezione è un elenco di cause, ordinate per frequenza. Serve quando un
artefatto è stato scritto, il file è al suo posto, e nell'interfaccia non
c'è.

\begin{center}
\begin{tabular}{@{}L{54mm}L{84mm}@{}}
\toprule
\textbf{Causa} & \textbf{Come si riconosce} \\
\midrule
\rigaTenue
Il registro è vecchio
  & Il processo era già avviato. Il registro di forma~B si costruisce
    all'import: chiamare \cod{reload\_registry()}, o riavviare. \\
Il nome della cartella comincia con \cod{\_}
  & È l'interruttore di spegnimento. Si vede guardando, ma è anche la cosa
    che non si guarda. \\
\rigaTenue
Manca il file atteso
  & \file{strategy.py}, \file{approach.py}, \file{weight\_mode.py},
    \file{overlay.py}. Il filtro delle cartelle lo richiede per nome. \\
Manca \cod{META} o \cod{compute}
  & La condizione di ammissione al registro li richiede entrambi. Senza, il
    modulo viene scartato senza una riga di registro. \\
\rigaTenue
\cod{META} non è fatto di soli letterali
  & Vale per i regimi, letti con \cod{ast.literal\_eval}: un \cod{np.pi}
    dentro i \cod{defaults} rende la regola invisibile. \\
La \cod{key} non coincide col nome del file
  & Rifiutata per i regimi; altrove l'artefatto compare con un nome diverso
    da quello che si sta cercando. \\
\rigaTenue
In \pkgOverlay{} si è esposto \cod{APPROACH\_DESCRIPTOR}
  & Il registro degli overlay cerca \cod{META}+\cod{compute}. Quel file è un
    \emph{approccio}, e come overlay non esiste. \\
\cod{stream\_type} fuori allowlist
  & Per gli approcci è l'unico caso che dà un \cod{ValueError} chiaro al
    caricamento. \\
\rigaTenue
\cod{output\_path\_pattern} vuoto
  & L'artefatto c'è, ma lo \emph{stream} non compare in \voce{Streams}: la
    dashboard salta i descrittori senza pattern. \\
Un altro modulo \file{ui\_ext\_*} è rotto
  & Se l'artefatto si vede ma manca il pannello che lo mostra, è questo
    (\vediSez{ui-ext}). \\
\bottomrule
\end{tabular}
\end{center}

\subsection{Il nodo della ricarica}

Il registro di forma~B si costruisce una volta, all'import del package. Un
artefatto scritto a runtime --- dallo Studio, dall'agente, dal marketplace ---
esiste su disco e non compare. \cod{reload\_registry()} risolve, ma deve fare
due cose che non sono ovvie:

\begin{codiceInterno}{python}{Le due accortezze}{overlays/\_\_init\_\_.py:144}
    Due accortezze necessarie perche' la ricarica sia davvero visibile:
    - COMPUTERS viene aggiornato SUL POSTO: `reporting.ensemble` ne tiene un
      riferimento preso all'import, e rilegarlo qui non lo aggiornerebbe;
    - EnsembleEngine.OVERLAY_REGISTRY e' un attributo di classe legato alla
      creazione della classe: va riassegnato esplicitamente.
\end{codiceInterno}

\noindent
Per le modalità di pesatura c'è un terzo punto, ed è quello che produce il
sintomo più sconcertante:

\begin{codiceInterno}{python}{Il terzo oggetto da aggiornare}{weights/\_\_init\_\_.py:117}
    try:                      # l'attributo di classe e' legato alla creazione
        from reporting.ensemble import EnsembleEngine
        EnsembleEngine.WEIGHTING_MODE_REGISTRY = REGISTRY
        # WEIGHTING_MODES e' una tupla DERIVATA dal registro alla creazione
        # della classe ed e' quella usata per validare le richieste: senza
        # aggiornarla, un weight mode nuovo compare nell'elenco ma il compute
        # risponde 400 "Invalid weighting_mode".
        EnsembleEngine.WEIGHTING_MODES = tuple(m["key"] for m in REGISTRY)
\end{codiceInterno}

\begin{attenzione}{Compare nell'elenco e risponde \cod{400}}
È il sintomo di una ricarica parziale: l'elenco che l'interfaccia mostra
viene da \cod{WEIGHTING\_MODE\_REGISTRY}, la validazione della richiesta da
\cod{WEIGHTING\_MODES}. Aggiornare il primo e non il secondo produce una voce
cliccabile che risponde «modalità non valida». Chi aggiunge un registro
nuovo di forma~B deve chiedersi quali altri oggetti ne derivino una copia al
momento della creazione della classe.
\end{attenzione}

\subsection{Due package che sono solo rinvii}

\file{reporting/weighting/} e \file{reporting/gates/} non contengono logica:
ri-esportano \pkgPesi{} e \pkgOverlay. Sono il ricordo di dove quelle due
famiglie stavano prima. Non vanno estesi, e il codice nuovo importa dalle
cartelle vere.

\subsection{Un modulo morto da non seguire}

Accanto al package \file{regime/logics/} esiste un file
\file{regime/logics.py}, \num{237} righe, con una bella tabella della
semantica delle etichette nel docstring. \textbf{Non è raggiungibile}: a
parità di nome il package vince sempre, e \cod{regime.logics} risolve a
\file{regime/logics/\_\_init\_\_.py}. Chi lo modifica non cambia niente e non
se ne accorge.

\subsection{Due incoerenze nelle mappe dell'interfaccia}

\begin{itemize}
  \item \file{dashboard/studio\_browser.py} dichiara \cod{overlay} e
        \cod{weightmode} come artefatti «a file», mentre tutti e ventuno
        quelli installati sono a cartella.
  \item \file{dashboard/agent.py} ha la stessa incoerenza al suo interno, e i
        suoi commenti la annotano già: «\emph{Qui indicava un file singolo, e
        altrove \cod{approach.py}: tre posti, tre risposte diverse alla stessa
        domanda}».
\end{itemize}

\noindent
Non rompono niente oggi, perché la scoperta accetta entrambe le forme. Vanno
sapute quando si aggiunge una famiglia: la mappa canonica da seguire è
\cod{\_FILE\_ZONA} in \file{dashboard/agent.py}.

\begin{codiceInterno}{python}{Il file principale atteso per famiglia}{dashboard/agent.py:66}
_FILE_ZONA = {"strategies": "strategy.py", "approaches": "approach.py",
              "overlays": "overlay.py", "weights": "weight_mode.py",
              "metrics": "metric.py", "indicators": "indicator.py"}
\end{codiceInterno}

\subsection{Un elenco obsoleto}

\file{strategies/\_\_init\_\_.py} elenca nel docstring «momentum, quality\_lv,
value, custom\_example». Delle quattro, due sono oggi spente (\cod{\_value},
\cod{\_custom\_example}) e mancano le altre sette attive. È documentazione
nel codice che ha smesso di essere vera, e il modo in cui ci si accorge è
cercarvi una strategia che esiste e non trovarla.

% ===========================================================================
%  sviluppatore/05-estendere/07-metrica.tex
% ===========================================================================
\section{Una metrica: i dati del sistema e la pagina dello stream}\label{sez:nuova-metrica}

Una metrica è il più piccolo degli artefatti: un file \file{metrics/<nome>.py}
con un \cod{META} di soli letterali e, sotto, codice vettoriale sull'equity
della pagina. Produce un numero (\cod{result(...)}) o uno o più grafici
(\cod{plot}, \cod{bars}, \cod{heatmap}, \cod{chart}). La si scrive in Studio,
che il \cod{META} lo compone da sé con i campi \voce{Output} e \voce{Chart
type}; il contratto completo, con gli esempi, è in
\file{metrics/METRIC\_TEMPLATE.md}.

Gira in \file{dashboard/metric\_runtime.py}, in una sandbox: builtins ridotti,
niente \cod{import}, niente attributi \cod{\_\_dunder\_\_} né privati. Fino a
questa versione vedeva soltanto l'equity della pagina e, quando la pagina lo
passava, il benchmark. Ora legge anche i dati del sistema.

\subsection{L'oggetto \cod{data}}

Nel codice di una metrica c'è \cod{data}, istanza di \cod{Data}
(\file{dashboard/metric\_data.py}). Ogni cartella di \pkgDati{} è una
categoria e ogni categoria è una classe, tutte con gli stessi metodi:

\begin{center}
\begin{tabular}{@{}L{36mm}L{102mm}@{}}
\toprule
\textbf{Metodo} & \textbf{Che cosa restituisce} \\
\midrule
\rigaTenue
\cod{names()} & i file della categoria: simboli per le categorie 1:1, file
  condivisi per le 1:N (\cod{settore\_tech}, \cod{USD}) \\
\cod{timeframes(nome)} & le sottocartelle di timeframe presenti, dalla più
  fine \\
\rigaTenue
\cod{columns(nome)} & i campi veri di un file \\
\cod{get(nome)} & il file com'è sul disco (con \cod{start} ed \cod{end} se
  ne ritaglia un periodo) \\
\rigaTenue
\cod{series(nome, campo)} & un campo, allineato alle date della metrica
  (\cod{align=False} lo dà grezzo) \\
\bottomrule
\end{tabular}
\end{center}

\noindent
Tranne \cod{timeframes}, accettano tutti \cod{timeframe=}: senza, vale quello
predefinito (sotto).
Le categorie del sistema hanno una classe propria: \cod{Assets} (con
\cod{close}, \cod{returns}, \cod{ohlcv}), \cod{Earnings} (con \cod{events}),
\cod{Fundamentals}, \cod{Rates}, \cod{Macro} (anche i campi del driver
macro), \cod{News} (titoli al giorno), \cod{MMS} (file \cod{<SIMBOLO>\_<TF>}).
Quelle create dalla pagina Dati sono \cod{Category}.

\begin{codiceInterno}{python}{Una metrica che legge un asset}{metrics/beta\_spy.py}
META = {"label": "Beta vs SPY", "type": "number", "chart_kind": None, "active": True}

r_m = data.assets.returns("SPY", timeframe="D1")   # aligned to the equity dates
both = returns.notna() & r_m.notna()
result(returns[both].cov(r_m[both]) / r_m[both].var(), "Beta vs SPY", fmt="ratio")
\end{codiceInterno}

Tre regole fanno il comportamento.

\begin{description}[leftmargin = 1.4em, style = nextline, itemsep = 4pt]
  \item[Il timeframe predefinito] È quello della strategia dietro la pagina,
    letto dal \cod{DESCRIPTOR} per albero sintattico, senza importare la
    strategia (\cod{strategy\_timeframes}). Per una strategia su più
    timeframe vale il più basso, cioè il più fine; per approcci e pagine
    senza strategia D1; in un ensemble il più basso fra le strategie dei suoi
    stream. Il codice ne sceglie un altro con \cod{timeframe="H1"}.
  \item[Il timeframe che manca] Se su disco non c'è, lo si costruisce da uno
    più fine che lo divide esattamente (H4 da H1 prima che da M5, D1 da M5),
    con le regole OHLCV del motore. Se non si può, \cod{DataError} elenca i
    timeframe che ci sono e propone \cod{timeframe=...}.
  \item[L'allineamento] \cod{series}, \cod{close} e \cod{returns} sono
    allineate alle date della metrica con la regola del motore: alla data
    $t$ l'ultimo valore noto a $t$, mai uno successivo. I dati intraday
    prima si chiudono al giorno, così la data $t$ riceve la chiusura di $t$.
\end{description}

\begin{nota}{Un dato che manca non è un errore del codice}
\cod{DataError} (sottoclasse di \cod{LookupError}) arriva alla pagina come
\enquote{Data not available}, non come errore. Studio salva comunque la
metrica, con un avviso: funzionerà quando il dato sarà caricato. La prova che
precede il salvataggio gira sullo stream con più dati, non su un'equity
finta: una metrica che legge \cod{data} si prova davvero.
\end{nota}

\subsection{Grafici a più serie}

Con \cod{chart\_kind: "figure"} un grafico ha più serie, punti con
etichetta, dispersioni, barre, aree e linee di riferimento, su un asse x di
date oppure di numeri --- quello che serve per la Security Market Line del
CAPM, dove sull'asse x c'è il beta. Il runtime produce una specifica JSON
(\cod{kind: "figure"}, \cod{traces}); \cod{plotCustomFigure} di
\file{static/js/charts.js} la disegna con Plotly.

\begin{codiceInterno}{python}{Un grafico con l'asse x numerico}{}
fig = chart("CAPM - Security Market Line", x="Beta", y="Annual return", y_format="pct")
fig.line([0.0, 2.0], [rf, rf + 2.0 * premium], name="SML", dash="dash")
fig.point(1.0, market, name="S&P 500")
fig.point(beta, holding, name="Holding", label="beta %.2f" % beta)
\end{codiceInterno}

\subsection{La pagina di uno stream}

Ogni metrica della pagina, predefinita o personale, numerica o grafica, si
nasconde e si rimette. Le predefinite sono gli elementi con \cod{data-sv} in
\file{templates/stream.html}; le personali le dispone
\file{static/js/stream\_metrics.js}: i numeri dopo quelli predefiniti, i
grafici in fondo alla pagina, dopo grafici e tabelle predefiniti, con
\voce{New chart metric} per ultima.

\begin{center}
\begin{tabular}{@{}L{44mm}L{94mm}@{}}
\toprule
\textbf{Pezzo} & \textbf{Che cosa fa} \\
\midrule
\rigaTenue
\cod{/api/stream-view/prefs} & la scelta, una per tutti gli stream:
  \cod{shown} (le chiavi accese) e \cod{known} (quelle esistenti quando la
  scelta è stata fatta). \cod{shown: null} = predefinite più personali
  attive \\
\cod{sv\_off(chiave)} & nel template: le predefinite nascoste partono
  nascoste, senza comparire per un istante \\
\rigaTenue
\cod{/api/metrics/compute} & con \cod{stream}, \cod{ids} e
  \cod{include\_inactive}: calcola solo le personali accese, sui risultati
  dello stream e con il suo benchmark \\
\cod{/studio?tab=metric} & con \cod{output=number} o \cod{output=chart}
  (\voce{New number metric}, \voce{New chart metric}): l'uscita scelta, un
  nome libero e, con \cod{from=}, il ritorno allo stream; con
  \cod{metric=<id>} riapre una metrica. \voce{Save} valida e salva in un
  colpo solo \\
\bottomrule
\end{tabular}
\end{center}

\begin{perche}{Perché \cod{known} oltre a \cod{shown}}
Una metrica creata dopo la scelta --- per esempio con \voce{New chart
metric} dalla pagina stessa --- non è stata scartata da nessuno. Senza
\cod{known} resterebbe nascosta finché qualcuno non va a cercarla
nell'elenco; con \cod{known} compare da sola, se è attiva.
\end{perche}

\subsection{La guardia: numpy e pandas filtrati}

Il codice di una metrica o di un indicatore può arrivare dalla Community.
numpy e pandas gli arrivavano così com'erano, e da un loro sottomodulo si
raggiungeva il sistema operativo senza scrivere un import
(\cod{pd.core.config\_init.os}); \cod{pd.read\_pickle} di un indirizzo
eseguiva codice arrivato dalla rete. \file{dashboard/sandbox\_guard.py} li
passa filtrati: funzioni e classi restano quelle vere, i sottomoduli solo
quelli di calcolo (\cod{np.linalg}, \cod{np.random}, \cod{pd.offsets},
\ldots); lettura e scrittura di file, \cod{eval}, \cod{query} e \cod{plot}
sono rifiutati già nel sorgente, anche chiamati per nome
(\cod{agg("to\_parquet")}). I dati del sistema si leggono con \cod{data}.

\begin{attenzione}{Uno strato, non un isolamento}
Come il validatore AST delle strategie (\file{adapters/ast\_validator.py}),
la guardia lavora nel processo dell'applicazione: chiude gli abusi diretti,
non un attaccante che costruisce i nomi a pezzi. La barriera vera resta
l'isolamento di processo. Il validatore delle strategie ha la stessa
apertura su numpy e pandas: una strategia può arrivare a \cod{os} per la
stessa strada.
\end{attenzione}


% ===========================================================================
%  sviluppatore/06-dashboard/00-capitolo.tex
% ===========================================================================
\chapter{L'interfaccia}\label{cap:dashboard}

% ===========================================================================
%  sviluppatore/06-dashboard/01-anatomia.tex
% ===========================================================================
\section{Anatomia di \pkgDashboard}\label{sez:anatomia-dashboard}

Trentasei file Python, \num{35} template Jinja, un centinaio di file
JavaScript e CSS. È l'oggetto più grande del progetto, e il modo più rapido
per orientarsi è guardare le dimensioni: dicono dove sta il peso.

\begin{center}
\begin{tabular}{@{}L{42mm}R{16mm}L{80mm}@{}}
\toprule
\textbf{File} & \textbf{KB} & \textbf{Di che cosa si occupa} \\
\midrule
\rigaTenue
\file{app.py}            & \num{226} & Il montaggio e \num{137} rotte. \\
\file{agent.py}          & \num{132} & L'assistente che scrive artefatti. \\
\rigaTenue
\file{community.py}      & \num{132} & Marketplace, contestazioni, account,
                                       messaggi, abbonamento. \\
\file{live\_runners.py}  & \num{90}  & I runner live: avvio, conto legato,
                                       kill-switch del conto, registro. \\
\rigaTenue
\file{web\_db.py}        & \num{90}  & L'accesso al Postgres del sito. \\
\file{agent\_tools\_sistema.py} & \num{80} & Gli strumenti con cui l'agente
                                       agisce sul sistema vivo. \\
\rigaTenue
\file{auth.py}           & \num{71}  & Accesso, registrazione, guardia
                                       globale, pagine HTML in linea. \\
\file{autopubblica.py}   & \num{49}  & Pubblicazione automatica degli
                                       artefatti, in un thread. \\
\rigaTenue
\file{genetic.py}        & \num{49}  & Strategie per programmazione genetica. \\
\file{visual\_debug.py}  & \num{45}  & Il debug visivo sui grafici. \\
\rigaTenue
\file{studio.py}         & \num{42}  & Editor, grafo a nodi, generazione. \\
\bottomrule
\end{tabular}
\end{center}

\subsection{Le quattro generazioni di codice, tutte vive}

Guardando come una rotta è registrata si capisce quando è stata scritta, e
quale stile seguire per estenderla.

\begin{center}
\begin{tabular}{@{}L{32mm}L{40mm}L{66mm}@{}}
\toprule
\textbf{Generazione} & \textbf{Dove} & \textbf{Come si riconosce} \\
\midrule
\rigaTenue
1 --- tutto in \file{app.py}
  & \cod{@app.route} dentro \simbolo{create\_app} & \num{137} rotte in un
    file di \num{226}~KB. \\
2 --- blueprint
  & \file{studio}, \file{community}, \file{agent}, \file{genetic}
  & \cod{Blueprint(...)}, nessun \cod{url\_prefix}. \\
\rigaTenue
3 --- estensioni v2
  & \file{ui\_ext\_*.py} & \simbolo{register(app, data\_layer, base\_path)},
    caricate da sole. \\
4 --- il guscio
  & \pkgApi{} + \pkgApplicazione & Non è nella dashboard: è l'altra
    applicazione (\vediSez{strangler}). \\
\bottomrule
\end{tabular}
\end{center}

\begin{pratica}{Quale stile usare oggi}
Per una rotta nuova che serve a una pagina: un modulo \file{ui\_ext\_*.py}.
Per una capacità nuova del sistema: un caso d'uso nel guscio, esposto
dall'API, chiamato dalla dashboard. In nessun caso si aggiunge una rotta a
\file{app.py}: è il file che tutte e tre le generazioni successive sono
servite a non far crescere.
\end{pratica}

\subsection{I template}

Trentadue pagine, tutte figlie di \file{base.html}, più tre frammenti
inclusi e \file{serve\_account.html}, che non è servito da nessuna rotta ma
dalla guardia di autenticazione.

Le pagine di accesso sono l'eccezione: \file{auth.py} non usa i template, ma
HTML in linea renderizzato con \simbolo{render\_template\_string}
(\cod{\_LOGIN\_HTML}, \cod{\_REGISTER\_HTML}, \cod{\_VERIFY\_HTML},
\cod{\_RESET\_HTML}, \cod{\_ATTESA\_HTML}). Chi cerca la pagina di login in
\file{templates/} non la trova.

\begin{nota}{Dove sta l'interfaccia del Live}
\file{live.html} non è servito da \percorsoHTTP{/live}: è servito da
\percorsoHTTP{/studio/runner-live}, cioè dal blueprint \cod{studio}. È il nome di
percorso più fuorviante del progetto (\vediSez{rotte}).
\end{nota}

% ===========================================================================
%  sviluppatore/06-dashboard/02-dati.tex
% ===========================================================================
\section{Come i dati arrivano alle pagine}\label{sez:data-layer}

La dashboard non legge i \file{parquet} a ogni richiesta. All'avvio li carica
tutti in memoria, in un oggetto che viene poi passato a tutto il resto:

\begin{codiceInterno}{python}{L'interfaccia pubblica}{dashboard/data\_layer.py}
    DataLayer()
        .startup()                  - carica streams + computa ensembles
        .get_ensemble(name)         - 'core' | 'full'
        .get_stream(name)           - stream singolo
        .get_stream_names()         - lista
        .get_spy_benchmark()        - SPY returns
\end{codiceInterno}

\noindent
È il secondo argomento che ogni modulo \file{ui\_ext\_*} riceve in
\simbolo{register(app, data\_layer, base\_path)}, e il motivo per cui una
pagina nuova non deve aprire file.

\subsection{Che cosa è pigro e che cosa no}

Gli stream si caricano all'avvio; due serie costose no.

\begin{codiceInterno}{python}{Due eccezioni dichiarate}{dashboard/data\_layer.py:86}
        # Stress di mercato (iso_stress) e risk-off macro (macro_regime): caricati o
        # CALCOLATI ON-DEMAND e cacheati (lazy), solo quando l'overlay corrispondente e'
        # attivo -> vedi _get_market_stress / _get_macro_riskoff. Nessun calcolo a startup.
\end{codiceInterno}

\noindent
Il criterio: una serie che serve solo a un overlay si calcola quando
quell'overlay è attivo. Caricarla all'avvio allungherebbe la partenza per
tutti, a beneficio di chi usa una funzione fra le sette.

\subsection{Gli ensemble preconfigurati non esistono più}

\begin{codiceInterno}{python}{La decisione, scritta nel codice}{dashboard/data\_layer.py:100}
        # Ensemble pre-configurati DISABILITATI: il sistema usa solo l'ensemble
        # custom (pagina /ensemble_custom), costruito a runtime dagli stream
        # selezionati. CORE/FULL rimossi su richiesta. Set vuoto => nessuna
        # scheda ENSEMBLE fissa in nav/home (restano stream singoli + custom).
        self.configured_ensemble_names: set[str] = set()
\end{codiceInterno}

\noindent
Il metodo \cod{get\_ensemble(name)} esiste ancora e il suo docstring parla
ancora di \cod{'core'} e \cod{'full'}: sono nomi che non si possono più
ottenere. Chi cerca di capire perché la pagina \percorsoHTTP{/ensemble/core} dà
\cod{404} lo trova qui.

\subsection{Come uno stream entra nell'elenco}

\file{dashboard/config.py} costruisce l'elenco degli stream in tre passi, e
solo il primo è un elenco scritto a mano:

\begin{enumerate}
  \item Gli stream dichiarati esplicitamente nella mappa interna.
  \item \textbf{Le strategie scoperte}, tramite
        \cod{discover\_strategies()}: entra chi ha un
        \cod{output\_path\_pattern}; chi non ce l'ha viene saltato con un
        \cod{continue} e un commento che dice «strategia non self-contained
        per dashboard».
  \item \textbf{Gli approcci scoperti}, con lo stesso meccanismo.
\end{enumerate}

\begin{perche}{Perché il passo 3 usa \cod{discover\_all()} e non \cod{discover\_approaches()}}
Il commento lo spiega: «anche gli overlay producono uno stream, e uno stream
va elencato indipendentemente da chi lo ha generato. Escluderli qui li
farebbe sparire dalle composizioni, che è l'unico posto in cui servono
davvero». È il complemento della scelta fatta nel registro
(\vediSez{nuovo-approccio}): gli overlay restano \emph{fuori} dall'elenco dei
backtest lanciabili e \emph{dentro} l'elenco degli stream componibili. Due
elenchi, due domande diverse.
\end{perche}

\begin{attenzione}{Il fallimento della scoperta è un avviso}
Entrambi i passi di scoperta sono dentro un \cod{try/except} che registra un
avviso e prosegue. Un errore di import in una strategia non impedisce alla
dashboard di partire: fa sparire quella strategia dall'elenco, e lascia una
riga \cod{Self-contained stream discovery failed} nel registro. È lo stesso
schema delle estensioni v2, con la stessa conseguenza: si nota guardando
l'interfaccia, non guardando un errore.
\end{attenzione}

% ===========================================================================
%  sviluppatore/06-dashboard/03-job.tex
% ===========================================================================
\section{Il lavoro lungo: job e registro in diretta}\label{sez:job}

Un backtest dura minuti o ore. Una richiesta HTTP non può aspettarlo, quindi
tutto ciò che è lungo passa dal \simbolo{JobManager} di
\file{execution/job\_manager.py}: la rotta restituisce subito un
identificativo, il lavoro continua in un thread, e il registro arriva al
browser man mano che viene prodotto.

\begin{center}
\begin{tabular}{@{}L{40mm}L{98mm}@{}}
\toprule
\textbf{Rotta} & \textbf{Che cosa fa} \\
\midrule
\rigaTenue
\rotta{GET}{/api/jobs} & Elenco. \\
\rotta{GET}{/api/job/<id>} & Stato di uno. \\
\rigaTenue
\rotta{GET}{/api/job/<id>/stream} & \textbf{SSE}: le righe di registro, in
  diretta. \\
\rotta{POST}{/api/job/<id>/cancel} & Annulla. \\
\rigaTenue
\rotta{DELETE}{/api/job/<id>} & Toglie dall'elenco un job concluso. \\
\rotta{DELETE}{/api/jobs/conclusi} & Svuota tutti i conclusi. \\
\bottomrule
\end{tabular}
\end{center}

\subsection{Gli stati}

\cod{PENDING} \freccia{} \cod{RUNNING} \freccia{} \cod{SUCCESS},
\cod{FAILED} o \cod{CANCELLED}. Sono un \cod{str, Enum}, quindi il valore
serializzato in JSON è direttamente la stringa minuscola.

\subsection{Il registro in diretta}

\begin{codiceInterno}{python}{Lo stream SSE}{dashboard/app.py:3720}
    def api_job_stream(job_id):
        """Server-Sent Events stream of log lines."""
        job = job_manager.get(job_id)
        if job is None:
            return jsonify({"error": "Job not found"}), 404

        def event_stream():
            # Manda subito le log lines gia' esistenti
            existing = list(job.log_lines)
            for line in existing:
                yield f"data: {line}\n\n"

            # Subscribe via queue per nuove lines
            q = Queue(maxsize=10000)
            def listener(line):
                try:
                    q.put_nowait(line)
                except Exception:
                    pass
            unsub = job.subscribe(listener)
\end{codiceInterno}

\noindent
Due dettagli meritano di essere notati, perché sono la ragione per cui
questo codice regge:

\begin{description}[leftmargin = 1.6em, style = nextline, itemsep = 4pt]
  \item[Prima lo storico, poi la sottoscrizione]
    Chi apre la pagina a metà di un backtest vede subito tutto ciò che è già
    successo, e poi il resto. Sottoscrivere senza mandare lo storico
    mostrerebbe una finestra vuota fino alla riga successiva.
  \item[La coda è limitata e il \cod{put} non blocca]
    \cod{Queue(maxsize=10000)} con \cod{put\_nowait} dentro un
    \cod{try/except} che ignora il fallimento. Se un browser smette di
    leggere, la sua coda si riempie e le righe nuove vengono scartate --- ma
    il backtest \emph{non} rallenta. L'alternativa (bloccare sulla coda
    piena) farebbe sì che una scheda lasciata aperta fermi un calcolo.
\end{description}

\begin{nota}{Il registro è limitato a \num{5000} righe}
\cod{deque(maxlen=5000)}: un backtest lungo ne produce molte di più, e le
prime scorrono via. Per conservarle tutte si guarda il file di registro, non
l'interfaccia.
\end{nota}

\subsection{Un lavoro per volta}

Il \simbolo{JobManager} nasce con \cod{max\_workers=1}. È il default
dichiarato del progetto, e vale sia per l'API FastAPI sia per la dashboard.
Due backtest in parallelo si contenderebbero la stessa cache dei dati e gli
stessi file di uscita.

\begin{attenzione}{Alzare \cod{max\_workers} non è una ottimizzazione}
La serializzazione non è un limite di prestazioni da togliere: è ciò che
rende i risultati riproducibili. Due job concorrenti che scrivono sotto
\file{output/strategies/} possono produrre uno stream mezzo di un
calcolo e mezzo dell'altro, e il risultato non dà errore --- è un
\file{parquet} valido con dentro numeri di due esecuzioni diverse.
\end{attenzione}

\subsection{Le notifiche}

\rotta{GET}{/api/notifications} alimenta la campanella, e mette insieme due
sorgenti: i job reali del \simbolo{JobManager} e i messaggi che arrivano dal
sito. È il motivo per cui la campanella mostra sia «backtest finito» sia
«hai una risposta su una contestazione»: una rotta, due origini.

% ===========================================================================
%  sviluppatore/06-dashboard/04-pagina-nuova.tex
% ===========================================================================
\section{Aggiungere una pagina o una rotta}\label{sez:pagina-nuova}

\begin{pratica}{La procedura}
\begin{enumerate}
  \item Creare \file{dashboard/ui\_ext\_<lettera>.py}. La lettera decide
        l'ordine di caricamento; \file{c} e \file{d} sono libere.
  \item Esporre \simbolo{register(app, data\_layer, base\_path)} ---
        \emph{quella firma esatta}: è il contratto che il caricatore si
        aspetta.
  \item Dentro, registrare le rotte con \cod{@app.route}.
  \item Se serve una pagina, aggiungere il template in
        \file{dashboard/templates/} con \cod{\{\% extends "base.html" \%\}}.
  \item Se la rotta tocca l'account, metterla sotto uno dei prefissi di
        \cod{\_SERVE\_ACCOUNT} (\vediSez{permessi}).
  \item Riavviare. Non c'è nessun elenco da aggiornare.
\end{enumerate}
\end{pratica}

\subsection{Il modello da copiare}

\file{ui\_ext\_e.py} è l'esempio migliore, per due ragioni che valgono per
qualunque rotta nuova.

\subsubsection{Il docstring dichiara il contratto}

\begin{commento}{dashboard/ui\_ext\_e.py}
Rotte del gruppo E (Marketplace: vetrina, Sell yours, messaggi, feedback,
account) per la UI v2.

    GET  /api/e/sell-status   chi puo' pubblicare: piano (livello e nome),
                              amministratore, lucchetto del marketplace e, se
                              non si puo', il motivo nelle parole del sito.
    GET  /api/e/account-data  per la pagina My account, in una richiesta sola:
                              prodotti pubblicati, acquisti, richieste di
                              rimborso proprie e messaggi non letti.

Solo letture, e mai un errore verso la pagina: se il sito non risponde la
rotta dice «non so» (\cod{unknown: true}) e la pagina lo mostra come tale, invece
di far credere che il piano manchi.
\end{commento}

\subsubsection{«Non so» non è «no»}

È la regola di progetto più utile di tutto il capitolo, e vale ben oltre
questo file. Quando la dashboard interroga il sito e il sito non risponde,
la risposta sbagliata è \cod{can\_publish: false}: l'utente vede «non puoi
pubblicare», che è un'affermazione, e va a cercare un problema nel proprio
abbonamento. La risposta giusta è \cod{unknown: true}, che la pagina mostra
come incertezza.

Lo stesso schema ritorna in \percorsoHTTP{/api/e/account-data}, che riporta un campo
\cod{errors}: «se una lettura fallisce lo dice, così la pagina non scambia
\emph{non riuscito} per \emph{nessun acquisto}».

\begin{attenzione}{Un elenco vuoto è un'affermazione}
Restituire \cod{[]} quando non si è riusciti a leggere significa dire
all'utente «non hai acquisti». Per chi ne ha, è un errore che preoccupa, e
che porta ad aprire una contestazione. La distinzione fra «zero» e «non so»
va mantenuta fino alla pagina.
\end{attenzione}

\subsection{Una funzione di servizio che vale la pena copiare}

Le risposte JSON della dashboard contengono spesso date e \cod{Decimal} che
arrivano dal Postgres del sito, e che \cod{jsonify} non sa serializzare:

\begin{codiceInterno}{python}{Serializzazione difensiva}{dashboard/ui\_ext\_e.py:30}
def _plain(v):
    """Valori serializzabili in JSON (date in ISO, Decimal in float)."""
    if isinstance(v, (str, int, float, bool, type(None))):
        return v
    if isinstance(v, dict):
        return {str(k): _plain(x) for k, x in v.items()}
    if isinstance(v, (list, tuple, set)):
        return [_plain(x) for x in v]
    iso = getattr(v, "isoformat", None)
    if callable(iso):
        try:
            return iso()
        except Exception:  # noqa: BLE001
            pass
    try:
        return float(v)
    except Exception:  # noqa: BLE001
        return str(v)
\end{codiceInterno}

\noindent
L'ultimo ripiego è \cod{str(v)}: qualunque cosa arrivi, la pagina riceve
qualcosa invece di un \cod{500}. In una rotta di sola lettura è la scelta
giusta.

\subsection{Le due pagine speciali}

\begin{description}[leftmargin = 1.6em, style = nextline, itemsep = 5pt]
  \item[\file{ui\_ext\_b.py} fa un redirect in un \cod{before\_request}]
    Una richiesta a \percorsoHTTP{/ensemble\_custom?project=<id>} viene rediretta a
    \percorsoHTTP{/projects/<id>}. È il modo in cui i collegamenti vecchi continuano a
    funzionare dopo che una pagina ha cambiato indirizzo. Chi cerca perché un
    URL si trasforma in un altro lo trova lì, non nelle rotte.
  \item[\file{ui\_ext\_a.py} intercetta le installazioni in un \cod{after\_request}]
    Annota le installazioni fatte dallo Studio per i tutorial. Un
    \cod{after\_request} che osserva le risposte altrui è potente e
    invisibile: va dichiarato nel docstring del modulo, come lì è fatto.
\end{description}


% ===========================================================================
%  sviluppatore/07-marketplace/00-capitolo.tex
% ===========================================================================
\chapter{Account, abbonamenti e marketplace}\label{cap:marketplace}

% ===========================================================================
%  sviluppatore/07-marketplace/01-identita.tex
% ===========================================================================
\section{Il ponte verso il sito}\label{sez:ponte}

La dashboard non sa chi è l'utente: lo chiede al sito. Non conserva password
né segreti: li tiene il sito. Tutto ciò che riguarda identità, denaro e
catalogo passa da \file{dashboard/app\_api.py}, e il meccanismo è un token
personale.

\subsection{Dove sta il token}

\begin{codiceInterno}{python}{La ricerca del token, in ordine}{dashboard/app\_api.py:67}
def token() -> str | None:
    global _TOKEN, _TOKEN_LOADED
    if _TOKEN_LOADED:
        return _TOKEN
    _TOKEN_LOADED = True
    env = (os.environ.get("LYRA_APP_TOKEN") or "").strip()
    if env:
        _TOKEN = env
        return _TOKEN
    try:
        p = _token_path()
        if p.exists():
            data = json.loads(p.read_text(encoding="utf-8"))
            _TOKEN = (data.get("token") or "").strip() or None
    except Exception:
        _TOKEN = None
    return _TOKEN
\end{codiceInterno}

\noindent
Variabile d'ambiente prima, file dopo, e il risultato è messo in cache per
tutta la vita del processo. La cache è la ragione per cui, dopo un accesso
nuovo, può servire un riavvio: \cod{\_TOKEN\_LOADED} resta vero.

\subsection{La chiamata che non solleva}

\begin{commento}{dashboard/app\_api.py:137 $\cdot$ \cod{chiama\_sito()}}
Chiamata al sito che NON solleva: ritorna (ok, dati).

Serve alle rotte che fanno solo da ponte — sicurezza dell'account, per
esempio. L'app non conserva password ne' segreti: li tiene il sito, e qui
si inoltra soltanto aggiungendo il token di chi ha fatto l'accesso.

Gli errori del sito arrivano con il proprio messaggio e il proprio codice
di stato: tradurli qui significherebbe perdere spiegazioni che il sito ha
scritto apposta.
\end{commento}

\begin{perche}{Perché gli errori del sito si inoltrano senza tradurli}
Il sito scrive messaggi pensati per essere letti dall'utente: «il tuo
abbonamento è scaduto il \ldots», «questa email è già registrata». Sostituirli
con un generico «errore del server» cancella l'informazione. La dashboard fa
da ponte, e un ponte non riscrive ciò che trasporta.
\end{perche}

\subsection{La regola del token giusto}

C'è un secondo segreto, oltre al token dell'utente: quello della dashboard.
Servono a cose diverse, e scambiarli produce il più insidioso dei difetti
di questo capitolo.

\begin{commento}{dashboard/community.py:76 $\cdot$ \cod{SI()}}
SI CHIAMA CON IL TOKEN DELL'UTENTE, non col segreto della dashboard.
Prima no, e per questo non ha mai funzionato: \percorsoHTTP{/api/app/abbonamento}
riconosce la persona dal token, quindi rispondeva 401, l'eccezione veniva
inghiottita e il risultato era «non abbonato» per tutti. Il sintomo era
invisibile perche' in beta gli unici a pubblicare erano gli
amministratori, che passano da un'altra strada.
\end{commento}

\begin{attenzione}{Tre cose insieme rendono un difetto invisibile}
Una risposta \cod{401}, un'eccezione inghiottita, e un ripiego prudente che
risponde «non abbonato». Presi uno per uno sono ragionevoli; insieme
producono un sistema che funziona male e non lo dice. Chi aggiunge una
chiamata al sito controlli quale dei due segreti serve, e verifichi il
percorso d'errore con l'utente giusto --- non con un amministratore, che
spesso passa da un'altra strada e non lo attraversa mai.
\end{attenzione}

\subsection{«Non sapere vale come no»}

Per la pubblicazione il ripiego prudente è invece la scelta corretta, ed è
dichiarata:

\begin{commento}{dashboard/community.py:72 $\cdot$ \cod{In()}}
In caso di errore risponde «livello 0, non puo' pubblicare»: non sapere
vale come no. Pubblicare per errore mette in vetrina roba di cui nessuno
risponde e si toglie dopo; non pubblicare si risolve riprovando.
\end{commento}

\noindent
Il criterio è l'asimmetria del danno: una pubblicazione sbagliata si corregge
con fatica e intanto è pubblica; una pubblicazione mancata si corregge
riprovando. Si noti che è il contrario della regola vista per le \emph{letture}
(\vediSez{pagina-nuova}): lì «non so» non va mostrato come «no», perché
nessuna azione irreversibile ne dipende. La differenza fra i due casi è se
c'è un'azione, non se c'è un'incertezza.

% ===========================================================================
%  sviluppatore/07-marketplace/02-lucchetto.tex
% ===========================================================================
\section{I due strati del permesso di pubblicare}\label{sez:lucchetto}

Chi può caricare un artefatto sul marketplace è deciso da due condizioni in
cascata, e confonderle è il modo tipico di sbagliare una modifica in
quest'area.

\figuraTikz{lucchetto-marketplace}
  {I due strati, nell'ordine in cui si valutano.}
  {lucchetto-marketplace}

\subsection{Primo strato: il lucchetto}

\begin{commento}{dashboard/community.py:30 $\cdot$ \cod{marketplace\_bloccato()}}
Il marketplace e' chiuso?

E' il PRIMO dei due strati che regolano la pubblicazione, e durante la
beta e' quello che conta: chiuso significa che nessuno carica niente — ne'
gratis ne' a pagamento — e nessuno compra. Prima al suo posto c'era un
interruttore che, spento, lasciava pubblicare gratis: non bastava, perche'
cio' che entra in vetrina ci resta e il filtro che decide cosa merita
starci (l'abbonamento) non era ancora in piedi.

Lo decide il sito, non l'applicazione: e' l'unico punto che chi usa il
programma non puo' modificare. In caso di dubbio risponde CHIUSO.
\end{commento}

\noindent
Tre proprietà, tutte deliberate:

\begin{itemize}
  \item \textbf{Lo decide il sito.} È l'unico permesso che chi usa il
        programma non può modificare agendo sulla propria macchina --- tutto
        il resto della dashboard è aperto (\vediSez{permessi}), questo no.
  \item \textbf{In caso di dubbio: chiuso.} Il \cod{except} finale restituisce
        \cod{True}.
  \item \textbf{Gli acquisti ne derivano.} \simbolo{vendite\_attive()} non è
        più un interruttore a sé: è \cod{not marketplace\_bloccato()}. Il
        commento dice perché resta come funzione: «due nomi per la stessa
        domanda sono due occasioni di rispondere diversamente».
\end{itemize}

\subsection{Secondo strato: l'abbonamento}

A marketplace aperto, serve uno dei tre livelli. I livelli sono quattro
contando chi non è abbonato:

\begin{codiceInterno}{python}{I quattro livelli}{dashboard/community.py:98}
# I QUATTRO LIVELLI, con i nomi che l'interfaccia mostra. Il livello 0 non e'
# un ripiego: e' un utente pieno, che sfoglia, scarica, compra e usa. Non
# mette soltanto roba in vetrina.
NOMI_LIVELLO = {0: "Free", 1: "Author", 2: "Studio", 3: "Desk"}
\end{codiceInterno}

\begin{center}
\begin{tabular}{@{}R{10mm}L{20mm}L{108mm}@{}}
\toprule
\textbf{Liv.} & \textbf{Nome} & \textbf{Che cosa può fare} \\
\midrule
\rigaTenue
\num{0} & Free & Tutto il sistema: sfoglia, scarica, compra, usa. \emph{Non}
  pubblica. \\
\num{1} & Author & In più: pubblicare sul marketplace. \\
\rigaTenue
\num{2} & Studio & Idem, prezzo diverso. \\
\num{3} & Desk & Idem, prezzo diverso. \\
\bottomrule
\end{tabular}
\end{center}

\begin{nota}{Perché oggi i tre piani fanno la stessa cosa}
È uno stato intermedio dichiarato, non una dimenticanza. Durante la beta i
tre livelli concedono solo la pubblicazione e differiscono per prezzo; la
decisione su che cosa aggiungere a Studio e Desk --- o se ridurli a uno solo
--- è rimandata all'uscita dalla beta. I prezzi attuali sono segnaposto e si
cambiano dal pannello di amministrazione senza toccare il codice.
\end{nota}

\subsection{Gli amministratori}

\begin{commento}{dashboard/community.py:111 $\cdot$ \cod{livello\_utente()}}
La categoria dell'utente: 0 se non abbonato, 1..3 secondo il piano.

Gli amministratori risultano al livello massimo, perche' possono fare
tutto cio' che fa il livello massimo e anche di piu'. Non sono un livello
vero — un livello scade, loro no — ma all'interfaccia serve un numero.
\end{commento}

\noindent
Un amministratore può inoltre installare un artefatto a pagamento senza
comprarlo (\rotta{POST}{/api/marketplace/<id>/installa-admin}) e,
simmetricamente, può scegliere di pagarlo comunque --- il che per un
amministratore equivale a una donazione all'autore, visto che non ne avrebbe
bisogno.

\begin{attenzione}{Gli amministratori non attraversano i percorsi d'errore}
È il motivo per cui il difetto del token (\vediSez{ponte}) è rimasto
invisibile per mesi: in beta gli unici a pubblicare erano gli
amministratori, che passano da un'altra strada. Chi collauda una modifica in
quest'area deve farlo con un utente \emph{non} amministratore, altrimenti
collauda il ramo sbagliato.
\end{attenzione}

% ===========================================================================
%  sviluppatore/07-marketplace/03-pagamenti.tex
% ===========================================================================
\section{I pagamenti}\label{sez:pagamenti}

Questa è l'area in cui un errore costa denaro vero, e la regola che la governa
è una sola: \textbf{l'accesso lo concede il webhook, e nessun altro}.

\figuraTikz{pagamento}
  {Il percorso di un acquisto. La freccia sbarrata è quella che esisteva e non
   esiste più.}
  {pagamento}

\subsection{La regola, e perché c'è}

\begin{commento}{dashboard/community.py:1982 $\cdot$ \cod{marketplace\_purchase\_return()}}
Ritorno dal browser dopo il pagamento.

NON concede nulla. Prima questa rotta creava l'acquisto come «paid»
fidandosi della sessione indicata nell'URL: chiunque poteva aprirlo e
ottenere gratis qualunque artefatto. Ora l'accesso lo apre soltanto il
webhook del sito, quando Stripe conferma di aver incassato; qui si
verifica solo se la licenza risulta gia' attiva, per dare all'utente
un messaggio sensato.

Il pagamento e' appena avvenuto e il webhook puo' arrivare un istante
dopo: se la licenza non c'e' ancora non e' un errore, si dice di
attendere.
\end{commento}

\begin{attenzione}{Un identificativo nell'URL non è una prova di pagamento}
È il difetto classico dell'integrazione con un fornitore di pagamenti:
l'utente torna su un indirizzo che contiene l'identificativo della sessione,
e il codice lo tratta come conferma. Ma quell'indirizzo è visibile,
indovinabile e condivisibile. L'unica conferma è quella che arriva
\emph{dal} fornitore, firmata, sul canale del webhook. Qualunque modifica in
quest'area che aggiunga una concessione fuori dal webhook va considerata un
difetto di sicurezza, non una comodità.
\end{attenzione}

\subsection{Quando il webhook non arriva}

I webhook si perdono. Esiste quindi una rotta di riconciliazione, costruita
con lo stesso principio:

\begin{codiceInterno}{python}{Riconciliare senza concedere}{dashboard/community.py:2072}
    def marketplace_riconcilia(artifact_id):
        """Completa un pagamento che il webhook non ha portato a termine.

        Non concede nulla di suo: chiede a Stripe se il pagamento risulta
        incassato, e solo in quel caso fa cio' che avrebbe fatto il webhook.
        """
\end{codiceInterno}

\noindent
«Non concede nulla di suo» è la frase che la rende legittima: non è una
scorciatoia, è una seconda interrogazione alla stessa autorità.

\subsection{Distinguere «non comprato» da «non ho potuto chiedere»}

\begin{commento}{dashboard/community.py:2094 $\cdot$ \cod{marketplace\_perche\_non\_mio()}}
Perche' l'artefatto non risulta acquistato.

Il possesso dipende da due letture: la licenza e l'acquisto. Se
entrambe tornano vuote non si distingue «non l'hai comprato» da «non
sono riuscito a chiedere», e sono cose opposte. Qui si guardano una
per una, dicendo anche se la lettura e' stata rifiutata.
\end{commento}

\noindent
È la stessa distinzione fra zero e «non so» del \vediSez{pagina-nuova},
applicata al caso in cui la differenza è fra «devi pagare» e «c'è un
problema nostro». Senza questa rotta, un utente che ha pagato e non vede
l'artefatto non ha modo di sapere di chi sia il problema, e l'unica mossa
che gli resta è aprire una contestazione.

\subsection{Gli abbonamenti: una rotta sola per due cose}

\begin{commento}{dashboard/community.py:2680 $\cdot$ \cod{abbonamento\_checkout()}}
Apre il pagamento di un piano, o cambia quello in corso.

Una rotta sola per entrambe le cose perche' la domanda dell'utente e'
una sola — «voglio il livello X con cadenza Y» — e il sito sa se
quella richiesta e' un abbonamento nuovo o un cambio con conguaglio.
Due rotte avrebbero costretto l'interfaccia a indovinarlo.
\end{commento}

\noindent
Il corpo della richiesta porta \cod{tier} e \cod{periodo}; le cadenze sono
mensile e annuale, con l'annuale a dieci mensilità (Desk solo annuale). La
disdetta, il cambio di carta e le fatture passano invece dal portale di
fatturazione (\rotta{POST}{/api/abbonamento/portale}), che è una pagina
ospitata dal fornitore: non c'è codice da mantenere per quelle funzioni.

Tre regole stanno sul sito, nella rotta che crea la sessione di pagamento,
e non nell'app:

\begin{itemize}
  \item \textbf{La prova di due mesi.} Se il piano la prevede e l'account non
        ne ha mai avuta una né ha mai avuto un abbonamento su Stripe, la
        sessione porta la fine della prova: mesi di calendario, calcolati
        dal minuto della chiave di idempotenza, così due clic nello stesso
        minuto mandano a Stripe la stessa richiesta. La carta si chiede
        comunque, e il primo addebito parte da solo alla fine della prova. Il
        webhook scrive \cod{subscriptionTrialEnd} finché la prova è in corso
        e \cod{trialUsedAt} la prima volta, che non si svuota più.
  \item \textbf{Desk non si compra dal checkout.} Il piano porta
        \cod{sales\_only} e la rotta risponde \num{409} con il codice
        \cod{sales\_only}, anche a un amministratore.
  \item \textbf{Il piano si riconosce dal codice.} Nei metadati
        dell'abbonamento il checkout scrive il codice del piano accanto al
        livello; il webhook e la rotta delle fatture cercano il piano prima
        dal codice, poi dal prodotto Stripe, poi dal numero. Desk era il
        livello~\num{3} prima che arrivasse Quant, e un abbonamento nato
        allora porta ancora quel numero.
\end{itemize}

\subsection{Le contestazioni}

Un acquirente che ritiene che un artefatto non renda come dichiarato può
aprire una pratica. Due rotte la rendono una verifica e non una discussione:

\begin{center}
\begin{tabular}{@{}L{54mm}L{84mm}@{}}
\toprule
\textbf{Rotta} & \textbf{Che cosa confronta} \\
\midrule
\rigaTenue
\cod{.../verifica-integrita} & I file installati contro i file venduti,
  uno per uno. \\
\cod{.../verifica-performance} & Le metriche dichiarate dal venditore
  contro quelle ottenute rieseguendo qui. \\
\bottomrule
\end{tabular}
\end{center}

\begin{perche}{Perché la verifica è automatica}
Una contestazione fra due persone su che cosa «renda» una strategia non si
risolve discutendo. Le due rotte producono un confronto riproducibile: o i
file sono quelli, o non lo sono; o le metriche si riproducono, o non si
riproducono. L'amministratore che chiude la pratica decide su un risultato,
non su due racconti.
\end{perche}

\subsection{Le variabili d'ambiente}

Le chiavi Stripe, SMTP e simili \emph{non} stanno in un \file{.env} locale:
si gestiscono dal sito e arrivano alla metà locale attraverso il canale~\num{2}
(\vediSez{due-meta}). È il motivo per cui un ambiente di sviluppo può avere
il circuito di pagamento funzionante senza che nessuna chiave stia sul disco
dello sviluppatore.

% ===========================================================================
%  sviluppatore/07-marketplace/04-artefatti.tex
% ===========================================================================
\section{Pubblicare e installare}\label{sez:pubblicare}

\subsection{Nessuna moderazione, un filtro all'ingresso}

\begin{codiceInterno}{python}{La scelta, scritta dove si applica}{dashboard/community.py:1537}
        # CHI PUO' PUBBLICARE
        # Solo amministratori e abbonati, anche per gli artefatti
        # gratuiti. Il motivo non e' il denaro: e' che chi pubblica
        # mette il proprio nome accanto a una strategia che altri
        # faranno girare sui propri soldi. Un filtro all'ingresso vale
        # piu' di qualunque moderazione a valle.
        #
        # Chi non e' abbonato resta un utente pieno: sfoglia, scarica,
        # compra e usa. Non puo' soltanto mettere roba in vetrina.
\end{codiceInterno}

\noindent
Un artefatto pubblicato va online subito: nessuna coda di approvazione. La
difesa è a monte --- serve un abbonamento --- e a valle le contestazioni
(\vediSez{pagamenti}), non in mezzo.

\subsection{I conflitti si dicono prima}

\rotta{GET}{/api/marketplace/<id>/conflicts} risponde a «che cosa esiste già con
lo stesso nome». Si chiama \emph{prima} di installare, e serve a non scoprire
dopo che un'installazione ha scavalcato una strategia su cui si stava
lavorando. Il nome canonico di un artefatto è il nome della sua cartella
(\vediSez{due-forme}): due artefatti con lo stesso nome sono la stessa
cartella.

\subsection{La pubblicazione automatica}

\file{dashboard/autopubblica.py} gira in un thread dall'avvio e pubblica da
sé gli artefatti che superano una soglia.

\begin{codiceInterno}{python}{Che cosa fa}{dashboard/autopubblica.py}
COSA FA
Ogni volta che una strategy viene scritta, e ogni volta che un project viene
creato, si guarda lo Sharpe out-of-sample del walk-forward. Sopra 2.0 e non
ancora in vetrina, l'artefatto va online GRATIS da solo. Per un project la
regola e' piu' articolata (vedi `valuta_progetto`): dipende da quante strategie
lo compongono e da cosa e' gia' pubblicato.
\end{codiceInterno}

\subsubsection{Il file di tracciamento}

Senza, ogni avvio rifarebbe gli stessi controlli e --- peggio ---
riproverebbe le stesse pubblicazioni. Il file registra, per ogni artefatto,
se è stato visionato e con quale conclusione. È un JSON indentato, leggibile
a mano: non un formato interno.

\subsubsection{Quando una voce viene riesaminata}

Tre casi, e solo questi:

\begin{enumerate}
  \item non c'è (artefatto nuovo, o file di tracciamento appena nato);
  \item \textbf{la conclusione non era definitiva} --- mancava il backtest,
        mancava la rete, nessuno era connesso;
  \item l'impronta del codice o quella dei risultati è cambiata.
\end{enumerate}

\begin{perche}{La distinzione fra un ostacolo e un giudizio}
Il commento del modulo la enuncia meglio di qualunque parafrasi: gli
ostacoli «non sono giudizi sull'artefatto: sono ostacoli, e un ostacolo che
passa deve poter cambiare l'esito. Una soglia non superata, invece, è un
giudizio: resta tale finché il codice o i risultati non cambiano».

È il criterio per decidere se un esito negativo vada riprovato o meno, e si
applica ben oltre questo modulo: ogni volta che un sistema registra un
«no», la domanda è se quel no dipendeva dall'oggetto o dalle circostanze.
\end{perche}

\begin{attenzione}{Un artefatto può finire in vetrina da solo}
È un comportamento attivo che agisce sul mondo esterno senza che nessuno lo
chieda in quel momento. Chi lavora su \file{autopubblica.py} tenga presente
che un errore qui non produce un messaggio: produce una pubblicazione. Il
file di tracciamento, essendo leggibile a mano, è lo strumento per
verificare che cosa il modulo abbia deciso e perché.
\end{attenzione}

\subsection{L'identità di chi pubblica}

Il marketplace ammette due identità diverse per lo stesso artefatto: il conto
dell'utente e la macchina. «Già caricata» ha senso solo rispetto a una delle
due, ed è la ragione per cui il pannello di amministrazione conta gli
artefatti \emph{per conto e per macchina}
(\rotta{GET}{/api/admin/artefatti}), non solo in totale.



% --- PARTE IV — tenerlo in piedi --------------------------------------------
\part{Tenerlo in piedi}
% ===========================================================================
%  sviluppatore/08-qualita/00-capitolo.tex
% ===========================================================================
\chapter{Prove, rilascio e diagnosi}\label{cap:qualita}

% ===========================================================================
%  sviluppatore/08-qualita/01-prove.tex
% ===========================================================================
\section{Come sono scritte le prove}\label{sez:prove}

Non c'è un \pkg{pytest} con mille asserzioni minute. Ci sono cinque suite,
ciascuna dedicata a un \emph{contratto} del sistema, che si lanciano come
script e stampano un resoconto leggibile. La scelta ha una conseguenza utile:
ogni suite si apre con un docstring che dice quale problema sorveglia, e
quello è il posto giusto dove cercare prima di modificare l'area che copre.

\begin{center}
\begin{tabular}{@{}L{48mm}R{14mm}L{76mm}@{}}
\toprule
\textbf{Suite} & \textbf{Righe} & \textbf{Che contratto sorveglia} \\
\midrule
\rigaTenue
\file{prova\_normalizzazione\_stream.py} & \num{691}
  & Backtest e live devono produrre lo \emph{stesso libro}
    (\vediSez{live}). \\
\file{prova\_beta\_accesso.py} & \num{614}
  & L'accesso riservato in beta, \textbf{contro il sito vero}. \\
\rigaTenue
\file{prova\_covarianza.py} & \num{489}
  & Lo shrinkage di Ledoit--Wolf e la modalità \cod{covariance}. \\
\file{prova\_pesi\_dichiarati.py} & \num{484}
  & Il contratto di \cod{Signal.target\_weight} e la modalità
    \cod{explicit}. \\
\rigaTenue
\file{prova\_permessi\_marketplace.py} & \num{429}
  & I due strati del permesso di pubblicare, \textbf{contro il sito vero}. \\
\bottomrule
\end{tabular}
\end{center}

\subsection{Due suite girano contro il sito vero}

Non usano finti. Avviano un Next.js e un Postgres reali e interrogano gli
endpoint veri. È più lento e più fastidioso da far partire, e la ragione è
dichiarata nel docstring: i permessi e l'accesso non vivono nel codice
Python, vivono nelle regole del sito. Una prova con dei finti verificherebbe
che la dashboard sa leggere una risposta --- cosa che non è mai stata il
problema --- e non che la risposta è quella giusta.

\subsection{Le proprietà che una prova deve avere qui}

Tre, e tutte e tre sono state imparate rompendole.

\begin{description}[leftmargin = 1.6em, style = nextline, itemsep = 6pt]

\item[Ripetibile]
  Una suite che modifica lo stato (prezzi dei piani, elenchi di ammessi) deve
  azzerarlo all'inizio, non alla fine: se una esecuzione si interrompe a metà,
  la pulizia finale non viene eseguita e la successiva parte da uno stato
  sporco. \file{prova\_permessi\_marketplace.py} comincia con una
  reimpostazione in SQL; \file{prova\_beta\_accesso.py} ha una funzione
  \simbolo{azzera\_prova()} e un elenco fisso di indirizzi.

\item[Senza soglie indovinate]
  Un'asserzione come \cod{assert scala < 1.3} è una soglia inventata: passa o
  fallisce per ragioni che nessuno ha verificato. La forma corretta è
  relativa --- \cod{assert scala\_degenere < 0.75 * scala\_corretta} --- perché
  confronta due cose che il sistema produce, invece di confrontarne una con
  un numero scritto a mano.

\item[Che isola ciò che misura]
  Misurare la proporzionalità confrontando due chiamate complete non
  funziona se ciascuna chiamata registra un'osservazione in più e cambia lo
  stato. Il confronto va fatto sulla funzione che si vuole misurare, da sola.

\end{description}

\begin{attenzione}{Una prova che non è ripetibile è peggio di nessuna prova}
Passa la prima volta, fallisce la seconda, e chi la vede fallire conclude
che la prova è inaffidabile --- non che il codice lo sia. Da quel momento il
risultato di quella suite viene ignorato, cioè la suite ha smesso di
esistere pur continuando a girare.
\end{attenzione}

\subsection{Il resto di \file{tests/}}

Oltre alle cinque suite ci sono: prove end-to-end (\file{e2e\_hedging.py},
\file{e2e\_regime\_modulation.py}), confronti A/B fra strategie, diagnosi
specifiche per singole strategie (\file{diagnose\_ising.py},
\file{diagnose\_regime\_signals.py}), un server SMTP finto
(\file{posta\_finta.py}) usato dalla prova dell'accesso, e verifiche
numeriche di migrazione (\file{verify\_migration\_numbers.py}).

% ===========================================================================
%  sviluppatore/08-qualita/02-diagnosi.tex
% ===========================================================================
\section{Quando qualcosa non va}\label{sez:diagnosi}

Un elenco dei sintomi che capitano davvero, con la causa più probabile
accanto. Ciascuno rimanda al punto della guida che lo spiega.

\begin{center}
\begin{longtable}{@{}L{50mm}L{88mm}@{}}
\toprule
\textbf{Sintomo} & \textbf{Dove guardare} \\
\midrule
\endfirsthead
\toprule
\textbf{Sintomo} & \textbf{Dove guardare} \\
\midrule
\endhead

\rigaTenue
Metà dell'interfaccia è sparita dopo una modifica
  & Un modulo \file{ui\_ext\_*} solleva all'import, e il \cod{try} unico
    impedisce anche agli altri di caricarsi. Nel registro:
    \cod{v2 page extensions not loaded}. \vediSez{ui-ext} \\

Una strategia non compare nell'elenco
  & \file{python tests/diagnostica\_strategie.py}, che stampa l'errore di
    import \emph{visibile} --- nella dashboard finisce in un avviso che
    nessuno legge. \\

\rigaTenue
Un artefatto nuovo non compare
  & Dieci cause possibili, in ordine di frequenza:
    \vediSez{trappole}. \\

Una modalità di pesatura compare ma risponde \cod{400}
  & Ricarica parziale del registro:
    \cod{EnsembleEngine.WEIGHTING\_MODES} non è stato aggiornato.
    \vediSez{trappole} \\

\rigaTenue
Il backtest si ferma alla prima barra con \cod{TypeError}
  & La firma di \cod{generate\_signals} usa nomi di parametro diversi da
    quelli attesi. \vediSez{decisione-esecuzione} \\

Lo Sharpe è straordinario
  & Prima di festeggiare: lookahead. Si legge oltre \cod{t\_decision},
    oppure una soglia è un quantile calcolato sull'intero campione.
    \vediSez{decisione-esecuzione} \\

\rigaTenue
Dal vivo il conto si muove molto più (o molto meno) del backtest
  & La normalizzazione per stream. È il difetto storico del progetto, e la
    suite che lo sorveglia è \file{prova\_normalizzazione\_stream.py}.
    \vediSez{live} \\

Il marketplace è vuoto in locale e pieno sul sito
  & Canale~\num{1}: token mancante o scaduto. Il token è in cache per tutta
    la vita del processo. \vediSez{ponte} \\

\rigaTenue
«Non sei abbonato» per tutti
  & Si sta chiamando il sito con il segreto della dashboard invece che con
    il token dell'utente. \vediSez{ponte} \\

Un download dati fallisce solo su una macchina
  & Canale~\num{2}: la chiave non è arrivata. Le variabili d'ambiente le
    gestisce il sito. \vediSez{due-meta} \\

\rigaTenue
Una pagina mostra «nessun acquisto» a chi ne ha
  & Una lettura fallita trattata come elenco vuoto. \vediSez{pagina-nuova} \\

Un mercato di regime non modula
  & Nessuno dei suoi \cod{driver} è disponibile: lo dice un avviso nel
    registro, non un errore. Serve
    \file{scripts/download\_regime\_drivers.py}. \vediSez{nuovo-regime} \\

\rigaTenue
Le pagine sono tutte vuote all'avvio
  & Non ci sono dati. Non è un errore: non c'è niente da mostrare.
    \vediSez{dati} \\

\percorsoHTTP{/ensemble/core} risponde \cod{404}
  & Gli ensemble preconfigurati sono stati disattivati: esiste solo
    l'ensemble costruito a runtime. \vediSez{data-layer} \\

\bottomrule
\end{longtable}
\end{center}

\subsection{Gli strumenti di diagnosi}

\begin{center}
\begin{tabular}{@{}L{48mm}L{90mm}@{}}
\toprule
\textbf{Strumento} & \textbf{A che serve} \\
\midrule
\rigaTenue
\file{tests/diagnostica\_strategie.py}
  & Per ogni cartella di strategia: importa o dà errore, con il traceback
    intero. Si lancia dalla radice, con lo \emph{stesso} interprete con cui
    si lancia la dashboard --- la differenza fra due ambienti virtuali è una
    causa frequente. \\
\rotta{GET}{/api/marketplace/diagnostica}
  & Verifica la catena della vetrina e dice in quale punto si rompe. \\
\rigaTenue
\rotta{GET}{/api/marketplace/<id>/perche-non-mio}
  & Perché un artefatto non risulta acquistato, distinguendo «non
    comprato» da «non sono riuscito a chiedere». \vediSez{pagamenti} \\
\rotta{GET}{/api/data/cache\_status}
  & Che cosa c'è su disco: numero di file, dimensione, intervallo di date. \\
\rigaTenue
\rotta{GET}{/api/universo}
  & Che cosa il sistema \emph{sa} di avere, contro che cosa ha davvero. La
    differenza fra le due è il registro \file{data/\_universe.json}. \\
\rotta{GET}{/api/streams/orfani}
  & Risultati senza una strategia che li abbia prodotti: restano dopo una
    cancellazione, e sporcano gli elenchi. \\
\bottomrule
\end{tabular}
\end{center}

\begin{pratica}{Il primo posto dove guardare, sempre}
Il registro. Questo sistema è pieno di \cod{try/except} che registrano un
avviso e proseguono, ed è una scelta giusta --- un plugin rotto non deve
impedire l'avvio --- che però sposta l'informazione dal comportamento al
registro. Un sintomo che non ha un errore corrispondente nell'interfaccia ha
quasi sempre una riga nel registro, e quasi sempre quella riga contiene il
nome del modulo colpevole.
\end{pratica}

% ===========================================================================
%  sviluppatore/08-qualita/03-rilascio.tex
% ===========================================================================
\section{Impacchettare e rilasciare}\label{sez:rilascio}

\subsection{Un comando solo}

\begin{codiceInterno}{console}{Dalla radice del progetto}{build/build\_app.py}
python build/build_app.py                # backend + app per il sistema corrente
python build/build_app.py --backend-only # solo l'eseguibile del backend
python build/build_app.py --app-only     # solo Electron (backend gia' pronto)
\end{codiceInterno}

\noindent
Tre passi, in ordine: PyInstaller produce un backend autonomo che non
richiede Python sulla macchina di chi riceve l'app; il risultato viene
copiato dove electron-builder lo cerca; electron-builder produce
l'installatore per il sistema corrente (\file{.exe}, \file{.dmg},
\file{AppImage}).

\subsection{Dove finiscono i dati dell'utente}

Non dentro l'applicazione. Il bundle è di sola lettura e ogni aggiornamento
lo sovrascrive: al primo avvio il backend crea uno spazio di lavoro nella
cartella utente e vi copia le librerie iniziali --- strategie, approcci,
modalità di pesatura. È lo stesso meccanismo di \simbolo{PROJECT\_ROOT}
descritto nel \vediSez{dati}, visto dall'altro capo.

\begin{attenzione}{Le librerie iniziali si copiano una volta}
Al primo avvio, non a ogni aggiornamento. Un artefatto nuovo aggiunto a una
versione nuova dell'applicazione \emph{non} compare automaticamente nello
spazio di lavoro di chi ha già installato: quello spazio è suo, e
sovrascriverlo cancellerebbe le sue modifiche. È una conseguenza voluta del
punto precedente, e va tenuta presente quando si aggiunge un artefatto di
serie.
\end{attenzione}

\subsection{Gli aggiornamenti automatici}

Chi ha già installato riceve le versioni nuove da solo: all'avvio
\pkg{electron-updater} interroga le release di GitHub, scarica in silenzio e
mostra la notifica «aggiornamento pronto» nella campanella. L'utente riavvia
quando gli fa comodo. Il repository si ricava dal remote git, o si dichiara
con \cod{GH\_REPO}.

\section{Il manifesto delle query}\label{sez:manifesto}

È il meccanismo di sicurezza più importante fra la metà locale e il sito, e
merita di essere capito prima di toccare qualunque cosa interroghi il
database.

\figuraTikz{manifest-query}
  {Il rapporto di fiducia, prima e dopo.}
  {manifest-query}

\begin{codiceInterno}{python}{Il motivo, dal generatore}{build/gen\_query\_manifest.py}
Perche' esiste
--------------
Finora l'app inviava SQL libero a ``/api/dashboard/query``, protetto da un
segreto uguale per tutte le installazioni. Chi estraeva quel segreto poteva
eseguire qualsiasi comando sul database di produzione.

Con questo file il rapporto si rovescia: il sito conosce in anticipo l'elenco
esatto delle query ammesse e l'app puo' solo chiederne una per NOME (un hash
del testo SQL). Qualunque altra cosa viene rifiutata, segreto o no.
\end{codiceInterno}

\subsection{L'identificativo di una query}

\cod{op\_id} è l'impronta SHA-256 del testo SQL normalizzato negli spazi,
troncata a \num{32} cifre esadecimali. Conseguenza diretta: \textbf{cambiare
una query cambia il suo identificativo}. Una query modificata e non
rigenerata nel manifesto viene rifiutata, segreto o no.

\subsection{La legatura all'utente}

Limitare le query non basta: chi è autenticato potrebbe chiedere i dati di un
altro passando un identificativo diverso. Per ogni query il generatore
individua il parametro che rappresenta l'utente, in due forme:

\begin{codiceInterno}{sql}{Le due forme riconosciute}{build/gen\_query\_manifest.py}
    WHERE "userId" = %s              -> conta i %s che lo precedono
    INSERT INTO t(id, "userId", ...) -> posizione nella lista delle colonne
\end{codiceInterno}

\noindent
Il sito \emph{sovrascrive} quel parametro con l'identificativo ricavato dal
token. Anche mentendo, il client ottiene solo i propri dati.

\begin{pratica}{Quando si tocca una query}
\begin{enumerate}
  \item Modificare l'SQL.
  \item \file{python build/gen\_query\_manifest.py} --- riscrive il JSON nel
        sito.
  \item \file{python build/gen\_query\_manifest.py --check} --- verifica che
        sia aggiornato. È il comando da mettere in una verifica automatica.
  \item Rilasciare \emph{prima} il sito, poi l'app: un'app che chiede un
        \cod{op\_id} che il sito non conosce ancora viene rifiutata.
\end{enumerate}
\end{pratica}

\begin{attenzione}{L'ordine di rilascio non è simmetrico}
Il sito accetta sia gli identificativi vecchi sia quelli nuovi finché il
manifesto li contiene entrambi; l'app no. Aggiornare l'app prima del sito
rompe le installazioni appena aggiornate, e il sintomo è un rifiuto senza
spiegazione --- perché dal punto di vista del sito è esattamente ciò che
deve succedere a una query sconosciuta.
\end{attenzione}


% ===========================================================================
%  sviluppatore/09-riferimento/00-capitolo.tex
% ===========================================================================
\chapter{Riferimento}\label{cap:riferimento}

Questo capitolo non si legge: si consulta. Le due tabelle lunghe --- le porte
e le rotte --- sono \emph{generate dal sorgente}, non trascritte, e si
rigenerano con gli script in \file{strumenti/}. Una tabella di riferimento
sbagliata è peggio di una tabella assente, perché la si consulta invece di
leggere il codice.

% ===========================================================================
%  sviluppatore/09-riferimento/01-porte.tex
% ===========================================================================
\section{Le porte}\label{sez:porte-tabella}

Trentuno \cod{Protocol} in \file{ports/\_\_init\_\_.py}, più la porta
\simbolo{Broker} ri-esportata da \file{live/broker.py}: trentadue contratti
in tutto. Sono ciò che i casi d'uso conoscono del mondo
(\vediSez{esagonale}).

\begin{center}
\begin{longtable}{@{}L{44mm}L{94mm}@{}}
\toprule
\textbf{Porta} & \textbf{Che cosa promette} \\
\midrule\endfirsthead
\toprule
\textbf{Porta} & \textbf{Che cosa promette} \\
\midrule\endhead

\multicolumn{2}{@{}l}{\textit{\small Strategie e risultati}}\\[2pt]
\rigaTenue
\simbolo{StrategyRepository} & Le strategie disponibili, come
  \simbolo{StrategyInfo}. \\
\simbolo{WalkForwardRunner} & Esegue un walk-forward. Restituisce il codice
  di uscita. \\
\rigaTenue
\simbolo{StreamCatalog} & Gli stream fuori campione disponibili. \\
\simbolo{ReportSink} & Dove finiscono i report di una esecuzione. \\
\rigaTenue
\simbolo{StrategyStore} & Persistenza su filesystem, sotto
  \file{strategies/<nome>/}. \\
\simbolo{StrategyValidator} & Validazione di sicurezza statica del sorgente
  (\simbolo{AST}). \\

\addlinespace
\multicolumn{2}{@{}l}{\textit{\small Composizione e analisi}}\\[2pt]
\rigaTenue
\simbolo{WeightingModeProvider} & L'elenco delle modalità di pesatura. \\
\simbolo{EnsembleService} & Calcola l'ensemble di $N$ stream con una
  pesatura scelta. \\
\rigaTenue
\simbolo{FactorAttributionService} & Scompone i rendimenti in alfa più beta
  ai fattori standard. \\
\simbolo{StressTestService} & Prova uno stream sulle crisi storiche
  (\num{2018}, \num{2020}, \num{2022}\ldots). \\

\addlinespace
\multicolumn{2}{@{}l}{\textit{\small Dati, tempo, mercato}}\\[2pt]
\rigaTenue
\simbolo{DataSource} & Accesso agli OHLCV; avvolge i loader esistenti. \\
\simbolo{Clock} & Orologio iniettabile. Esiste perché
  \cod{datetime.now()} sparso nel codice rende i test non deterministici. \\
\rigaTenue
\simbolo{BrokerGateway} & Verso un broker, parlando in tipi di dominio:
  l'adattatore traduce, e i casi d'uso non toccano i tipi
  dell'infrastruttura. \\
\simbolo{Broker} & \emph{Ri-esportata} da \simbolo{live.broker.LiveBroker}:
  era già una classe base astratta, cioè già una porta. \\
\rigaTenue
\simbolo{ModelRegistryPort} & Registro versionato dei modelli ML, con
  promozione. \\

\addlinespace
\multicolumn{2}{@{}l}{\textit{\small Generazione e compilazione degli artefatti}}\\[2pt]
\rigaTenue
\simbolo{CodeGenerator} & Generazione di codice. Il backend è Ollama in
  locale o un modello proprietario; \cod{list\_models()} alimenta il menu. \\
\simbolo{NodeCompiler} & Compila un grafo a nodi in sorgente Python. \\
\rigaTenue
\simbolo{FidelityScorer} & Confronta i trade disegnati a mano con quelli
  della strategia. \\
\simbolo{WeightModeValidator} & Sicurezza più struttura (\cod{META} e
  \cod{compute}). \\
\rigaTenue
\simbolo{WeightModeStore} & \cod{exists(key)}, \cod{save(key, sorgente,
  overwrite)}. \\
\simbolo{WeightModeNodeCompiler} & Grafo a nodi \freccia{} modalità di
  pesatura. \\
\rigaTenue
\simbolo{IndicatorValidator} & Come sopra, per un indicatore. \\
\simbolo{IndicatorNodeCompiler} & Grafo a nodi \freccia{} indicatore. \\
\rigaTenue
\simbolo{OverlayValidatorPort} & Come sopra, per un overlay. \\
\simbolo{OverlayStore} & \cod{exists(nome)}, \cod{save(nome, sorgente,
  overwrite)}. \\
\rigaTenue
\simbolo{OverlayNodeCompiler} & Grafo a nodi \freccia{} overlay. \\

\addlinespace
\multicolumn{2}{@{}l}{\textit{\small Operatività dal vivo}}\\[2pt]
\rigaTenue
\simbolo{LiveRunnerValidatorPort} & \cod{validate(cfg)} \freccia{}
  \simbolo{LiveRunnerValidation}. \\
\simbolo{LiveRunnerStore} & \cod{exists}, \cod{save}, \cod{load},
  \cod{list}. \\
\rigaTenue
\simbolo{LiveRunnerUploader} & Carica una configurazione di runner. \\

\addlinespace
\multicolumn{2}{@{}l}{\textit{\small Marketplace}}\\[2pt]
\rigaTenue
\simbolo{ArtifactPublisherPort} & Carica un artefatto sul proprio conto sul
  sito. \\
\simbolo{AccountArtifactsPort} & Legge gli artefatti già presenti sul conto
  --- serve a non duplicarli. \\
\rigaTenue
\simbolo{ArtifactCodeSourcePort} & Legge il codice locale di un artefatto,
  per tipo e nome. \\

\bottomrule
\end{longtable}
\end{center}

\begin{nota}{Dieci porte senza docstring}
Le porte della famiglia \cod{WeightMode*}, \cod{Indicator*}, \cod{Overlay*} e
\cod{LiveRunner*} non hanno un docstring: si capiscono dalla firma, che è di
due o tre metodi. Non è grave --- la firma \emph{è} il contratto --- ma le
altre ventuno ce l'hanno, e l'asimmetria si nota aprendo il file.
\end{nota}

% ===========================================================================
%  sviluppatore/09-riferimento/02-costanti.tex
% ===========================================================================
\section{Le costanti e la configurazione}\label{sez:costanti}

\subsection{\file{config/quantsys.yaml}}

È il file centrale. Dopo una modifica serve riavviare la dashboard: non
viene riletto a caldo.

\begin{center}
\begin{tabular}{@{}L{46mm}R{18mm}L{72mm}@{}}
\toprule
\textbf{Chiave} & \textbf{Valore} & \textbf{Che cosa governa} \\
\midrule
\multicolumn{3}{@{}l}{\textit{\small \cod{data:}}}\\[2pt]
\rigaTenue
\cod{start\_date}        & \num{2010}-01-01 & Inizio del campione. \\
\cod{train\_years}       & \num{5}  & Ampiezza della finestra di
                                      addestramento. \\
\rigaTenue
\cod{test\_years}        & \num{1}  & Ampiezza della finestra di prova. \\
\cod{slide\_years}       & \num{1}  & Di quanto avanza la finestra. \\
\rigaTenue
\cod{warmup\_days}       & \num{300} & Barre saltate all'inizio. \\
\cod{min\_folds}         & \num{3}  & Minimo di finestre perché il
                                      walk-forward abbia senso. \\
\rigaTenue
\cod{benchmark\_symbol}  & SPY      & Il riferimento dei confronti. \\

\addlinespace
\multicolumn{3}{@{}l}{\textit{\small \cod{portfolio:}}}\\[2pt]
\rigaTenue
\cod{initial\_capital}   & \num{1000000} & Capitale iniziale. \\
\cod{portfolio\_target\_vol\_annual} & \num{0.12} & Volatilità annua
                                      obiettivo. \\
\rigaTenue
\cod{max\_position\_weight} & \num{0.05} & Tetto per strumento. \\
\cod{max\_asset\_class\_gross} & \num{1.00} & Tetto lordo per classe. \\
\rigaTenue
\cod{max\_gross\_leverage} & \num{1.0} & Tetto lordo totale. \\

\addlinespace
\multicolumn{3}{@{}l}{\textit{\small \cod{ensemble:}}}\\[2pt]
\rigaTenue
\cod{target\_vol}        & \num{0.10} & Volatilità a cui ogni stream viene
                                      portato \emph{prima} di essere pesato. \\
\cod{scale\_lookback}    & \num{252} & Barre della volatilità trailing. \\
\rigaTenue
\cod{regime\_window}     & \num{90}  & Finestra del regime. \\
\cod{regime\_sh\_low} / \cod{\_high} & \num{0.0} / \num{1.5} & Soglie di
                                      Sharpe per la modulazione. \\
\bottomrule
\end{tabular}
\end{center}

\begin{attenzione}{\cod{ensemble.target\_vol} non è un parametro di reporting}
Il commento nel file lo dice: «non è un dettaglio di reporting: lo stesso
fattore moltiplica il libro eseguito dal vivo». Uno stream a volatilità molto
bassa chiede molta leva per arrivare all'obiettivo, e il tetto resta
\cod{max\_gross\_leverage}. Cambiare questo valore cambia le dimensioni delle
posizioni su un conto reale (\vediSez{live}).
\end{attenzione}

\noindent
I valori della tabella sono quelli del file come si trova nel progetto; i
default del codice, usati quando una chiave manca, sono diversi e stanno in
\simbolo{AllocatorConfig} (\vediSez{allocatore}).

\subsection{Le costanti del motore}

\begin{center}
\begin{tabular}{@{}L{40mm}R{14mm}L{82mm}@{}}
\toprule
\textbf{Costante} & \textbf{Valore} & \textbf{Dove, e a che serve} \\
\midrule
\rigaTenue
\cod{SCALE\_MIN\_OBS} & \num{20} & \file{reporting/ensemble.py}. Sotto
  questa soglia non c'è stima di volatilità e il fattore vale \cod{1.0}. \\
\cod{SCALE\_MAX\_LEVA} & \num{10.0} & Idem. Rapporto massimo fra obiettivo e
  volatilità stimata: è il tetto alla leva quando la volatilità tende a
  zero. \\
\rigaTenue
\cod{max\_workers} & \num{1} & \file{execution/job\_manager.py}. I job sono
  serializzati (\vediSez{job}). \\
\cod{max\_log\_lines} & \num{5000} & Idem. Righe di registro conservate per
  ogni job. \\
\rigaTenue
\cod{min\_lookback\_days} & \num{252} & Default di
  \simbolo{StrategyDescriptor}. \\
\bottomrule
\end{tabular}
\end{center}

\subsection{Le variabili d'ambiente}

\begin{center}
\begin{tabular}{@{}L{38mm}L{100mm}@{}}
\toprule
\textbf{Variabile} & \textbf{Effetto} \\
\midrule
\rigaTenue
\cod{LYRA\_HOME} & Sposta la cartella di lavoro (\vediSez{dati}). \\
\cod{LYRA\_APP\_TOKEN} & Scavalca il token personale letto da file
  (\vediSez{ponte}). \\
\rigaTenue
\cod{FLASK\_HOST} & Default \cod{127.0.0.1}. Cambiarlo espone la dashboard
  (\vediSez{avvio}). \\
\cod{FLASK\_PORT} & Default \num{5000}. \\
\rigaTenue
\cod{GH\_REPO} & Repository per gli aggiornamenti automatici
  (\vediSez{rilascio}). \\
\bottomrule
\end{tabular}
\end{center}

\noindent
Le chiavi dei servizi esterni --- Polygon, FRED, Stripe, SMTP, Google,
Postgres --- \emph{non} stanno qui: le gestisce il sito e arrivano alla metà
locale attraverso il canale~\num{2} (\vediSez{due-meta}).

% ===========================================================================
%  sviluppatore/09-riferimento/03-guide.tex
% ===========================================================================
\section{Le guide dentro il progetto}\label{sez:guide}

Il progetto porta con sé una trentina di file Markdown. Non sono ridondanti
rispetto a questa guida: sono \emph{riferimenti di dettaglio} per chi sta già
scrivendo un artefatto, mentre questo documento spiega come il sistema sta in
piedi. Dove i due divergono, vince il codice --- e quando è stato possibile,
questa guida lo dice (\vediSez{forma-b}).

\begin{center}
\begin{longtable}{@{}L{54mm}R{12mm}L{72mm}@{}}
\toprule
\textbf{File} & \textbf{Righe} & \textbf{A che serve} \\
\midrule\endfirsthead
\toprule
\textbf{File} & \textbf{Righe} & \textbf{A che serve} \\
\midrule\endhead

\multicolumn{3}{@{}l}{\textit{\small Architettura}}\\[2pt]
\rigaTenue
\file{ARCHITETTURA.md} & \num{1220} & Il documento di architettura originale. \\
\file{README.md} & \num{122} & Il ragionamento dietro il guscio a strati. \\
\rigaTenue
\file{MODIFICHE\_QUANTSYS.md} & \num{358} & Storia delle modifiche al motore. \\
\file{dashboard/MODIFICHE.md} & \num{778} & Storia delle modifiche
  all'interfaccia. \\
\rigaTenue
\file{dashboard/MODIFICHE-UI.md} & \num{494} & Idem, lato grafico. \\
\file{MODIFICHE\_DASHBOARD\_V2.md} & \num{118} & Che cosa ha portato la v2:
  pagine nuove, finestra dei pagamenti, marchio. \\

\addlinespace
\multicolumn{3}{@{}l}{\textit{\small Strategie}}\\[2pt]
\rigaTenue
\file{GUIDE\_STRATEGIE\_CUSTOM.md} & \num{2281} & \textbf{Il riferimento
  principale.} Contratto completo, avvio rapido, descrittore campo per
  campo. \\
\file{GUIDE\_STRATEGIE\_BLUEPRINT.md} & \num{288} & La strategia come grafo
  di nodi invece che in Python. Alternativo, non complementare. \\
\rigaTenue
\file{GUIDE\_REGIME.md} & \num{487} & Come \emph{usare} i regimi dentro una
  strategia: \cod{regime\_weights}, anti-pattern, validazione. \\
\file{GUIDE\_HEDGING.md} & \num{248} & Posizioni multiple sullo stesso
  simbolo (il default è il netting). \\
\rigaTenue
\file{GUIDE\_STRATEGIE\_MMS.md} & \num{625} & Le feature della heatmap di
  liquidità, via \cod{required\_mms\_features}. \\
\file{GUIDE\_VOLUME\_PROFILE.md} & \num{526} & Le zone di supporto e
  resistenza, via \cod{required\_volume\_profile}. \\
\rigaTenue
\file{STRATEGIE\_FX\_CANDIDATE.md} & \num{231} & Un documento di
  \emph{proposta}: strategie in valuta scrivibili nel framework attuale. Non
  è implementazione. \\

\addlinespace
\multicolumn{3}{@{}l}{\textit{\small Approcci}}\\[2pt]
\rigaTenue
\file{approaches/README.md} & \num{303} & Che cos'è un approccio e perché
  scavalca il motore. \\
\file{GUIDE\_APPROCCI\_CUSTOM.md} & \num{359} & Il riferimento. Si apre con la
  tabella di confronto strategia/approccio. \\
\rigaTenue
\file{GUIDE\_APPROCCI\_BLUEPRINT.md} & \num{76} & L'approccio a nodi; copre
  solo le differenze dalla guida delle strategie. \\

\addlinespace
\multicolumn{3}{@{}l}{\textit{\small Artefatti di forma B}}\\[2pt]
\rigaTenue
\file{weights/WEIGHT\_MODE\_TEMPLATE.md} & \num{197} & Contratto, import
  ammessi e vietati, esempio funzionante. \textbf{Attenzione al valore di
  \cod{kind}} (\vediSez{ensemble}). \\
\file{overlays/OVERLAY\_TEMPLATE.md} & \num{212} & Si apre con una tabella
  che separa overlay e approccio su sei assi. \\
\rigaTenue
\file{indicators/INDICATOR\_TEMPLATE.md} & \num{127} & Modello di un
  indicatore. \\
\file{metrics/METRIC\_TEMPLATE.md} & \num{320} & Modello di una metrica: i dati
  del sistema (\cod{data}), i grafici a più serie, il CAPM come esempio. \\

\addlinespace
\multicolumn{3}{@{}l}{\textit{\small Regimi}}\\[2pt]
\rigaTenue
\file{regime/GUIDE\_REGIMI.md} & \num{190} & Come \emph{costruire} un regime:
  i due artefatti, la catena, l'isteresi. \\
\file{logics/LOGIC\_TEMPLATE.md} & \num{84} & Modello della regola. \\
\rigaTenue
\file{markets/MARKET\_TEMPLATE.md} & \num{80} & Modello del mercato. \\

\addlinespace
\multicolumn{3}{@{}l}{\textit{\small Rilascio}}\\[2pt]
\rigaTenue
\file{build/README.md} & \num{289} & L'impacchettamento. \\
\file{build/RILASCIO.md} & \num{186} & Pubblicare un aggiornamento. \\
\rigaTenue
\file{DOCKER.md} & \num{153} & Avviare con Docker. \\

\bottomrule
\end{longtable}
\end{center}

\begin{nota}{Due guide che sembrano la stessa e non lo sono}
\file{regime/GUIDE\_REGIMI.md} è «come costruirne uno»;
\file{strategies/GUIDE\_REGIME.md} è «come usarlo dentro una strategia». Il
nome quasi identico è infelice, la distinzione è reale, e i due documenti si
rimandano a vicenda.
\end{nota}


\section{Le rotte della dashboard}\label{sez:rotte-tabella}

% ===========================================================================
%  GENERATO DA strumenti/genera-rotte.py --- non modificare a mano.
%  Per aggiornarlo:  python3 strumenti/genera-rotte.py /percorso/del/progetto
% ===========================================================================

\subsection{app.py --- il nucleo}
\begin{longtable}{@{}L{17mm}L{74mm}L{45mm}@{}}
\toprule
\textbf{Metodi} & \textbf{Percorso} & \textbf{Funzione} \\
\midrule\endfirsthead
\toprule
\textbf{Metodi} & \textbf{Percorso} & \textbf{Funzione} \\
\midrule\endhead
{\scriptsize\sffamily GET} & {\ttfamily\scriptsize /\allowbreak } & {\ttfamily\scriptsize index} \\
\rigaTenue {\scriptsize\sffamily GET} & {\ttfamily\scriptsize /\allowbreak admin} & {\ttfamily\scriptsize admin\_page} \\
{\scriptsize\sffamily GET POST} & {\ttfamily\scriptsize /\allowbreak api/\allowbreak account/\allowbreak security} & {\ttfamily\scriptsize api\_account\_security} \\
\rigaTenue {\scriptsize\sffamily GET POST} & {\ttfamily\scriptsize /\allowbreak api/\allowbreak admin/\allowbreak amministratori} & {\ttfamily\scriptsize api\_admin\_elenco, api\_admin\_modifica} \\
{\scriptsize\sffamily GET} & {\ttfamily\scriptsize /\allowbreak api/\allowbreak admin/\allowbreak artefatti} & {\ttfamily\scriptsize api\_admin\_artefatti} \\
\rigaTenue {\scriptsize\sffamily GET} & {\ttfamily\scriptsize /\allowbreak api/\allowbreak admin/\allowbreak artefatti/\allowbreak zip} & {\ttfamily\scriptsize api\_admin\_artefatti\_zip} \\
{\scriptsize\sffamily GET} & {\ttfamily\scriptsize /\allowbreak api/\allowbreak admin/\allowbreak beta} & {\ttfamily\scriptsize api\_admin\_beta} \\
\rigaTenue {\scriptsize\sffamily GET POST} & {\ttfamily\scriptsize /\allowbreak api/\allowbreak admin/\allowbreak beta-access} & {\ttfamily\scriptsize api\_admin\_beta\_access, api\_admin\_beta\_access\_modifica} \\
{\scriptsize\sffamily GET} & {\ttfamily\scriptsize /\allowbreak api/\allowbreak admin/\allowbreak feedback} & {\ttfamily\scriptsize api\_admin\_feedback} \\
\rigaTenue {\scriptsize\sffamily GET} & {\ttfamily\scriptsize /\allowbreak api/\allowbreak admin/\allowbreak feedback/\allowbreak export} & {\ttfamily\scriptsize api\_admin\_feedback\_export} \\
{\scriptsize\sffamily POST} & {\ttfamily\scriptsize /\allowbreak api/\allowbreak admin/\allowbreak feedback/\allowbreak group} & {\ttfamily\scriptsize api\_admin\_feedback\_group} \\
\rigaTenue {\scriptsize\sffamily GET POST} & {\ttfamily\scriptsize /\allowbreak api/\allowbreak admin/\allowbreak piani} & {\ttfamily\scriptsize api\_admin\_piani, api\_admin\_piani\_modifica} \\
{\scriptsize\sffamily GET} & {\ttfamily\scriptsize /\allowbreak api/\allowbreak admin/\allowbreak users} & {\ttfamily\scriptsize api\_admin\_users} \\
\rigaTenue {\scriptsize\sffamily GET POST} & {\ttfamily\scriptsize /\allowbreak api/\allowbreak admin/\allowbreak vendite} & {\ttfamily\scriptsize api\_admin\_vendite, api\_admin\_vendite\_modifica} \\
{\scriptsize\sffamily GET} & {\ttfamily\scriptsize /\allowbreak api/\allowbreak brokers/\allowbreak <broker>/\allowbreak check} & {\ttfamily\scriptsize api\_broker\_check} \\
\rigaTenue {\scriptsize\sffamily GET POST} & {\ttfamily\scriptsize /\allowbreak api/\allowbreak brokers/\allowbreak credentials} & {\ttfamily\scriptsize api\_broker\_credentials} \\
{\scriptsize\sffamily GET POST} & {\ttfamily\scriptsize /\allowbreak api/\allowbreak brokers/\allowbreak test/\allowbreak config} & {\ttfamily\scriptsize api\_test\_config} \\
\rigaTenue {\scriptsize\sffamily GET POST} & {\ttfamily\scriptsize /\allowbreak api/\allowbreak categorie} & {\ttfamily\scriptsize api\_categorie, api\_categorie\_crea} \\
{\scriptsize\sffamily POST} & {\ttfamily\scriptsize /\allowbreak api/\allowbreak categorie/\allowbreak <cat>} & {\ttfamily\scriptsize api\_categorie\_modifica} \\
\rigaTenue {\scriptsize\sffamily POST} & {\ttfamily\scriptsize /\allowbreak api/\allowbreak categorie/\allowbreak <cat>/\allowbreak elimina} & {\ttfamily\scriptsize api\_categorie\_elimina} \\
{\scriptsize\sffamily GET POST} & {\ttfamily\scriptsize /\allowbreak api/\allowbreak config/\allowbreak settings} & {\ttfamily\scriptsize api\_config\_settings\_get, api\_config\_settings\_save} \\
\rigaTenue {\scriptsize\sffamily GET} & {\ttfamily\scriptsize /\allowbreak api/\allowbreak config/\allowbreak yaml} & {\ttfamily\scriptsize api\_config\_yaml} \\
{\scriptsize\sffamily GET} & {\ttfamily\scriptsize /\allowbreak api/\allowbreak data/\allowbreak cache\_status} & {\ttfamily\scriptsize api\_data\_cache\_status} \\
\rigaTenue {\scriptsize\sffamily POST} & {\ttfamily\scriptsize /\allowbreak api/\allowbreak data/\allowbreak capabilities} & {\ttfamily\scriptsize api\_data\_capabilities} \\
{\scriptsize\sffamily GET POST} & {\ttfamily\scriptsize /\allowbreak api/\allowbreak data/\allowbreak categories} & {\ttfamily\scriptsize api\_data\_categories, api\_data\_category\_create} \\
\rigaTenue {\scriptsize\sffamily POST} & {\ttfamily\scriptsize /\allowbreak api/\allowbreak data/\allowbreak categories/\allowbreak <name>/\allowbreak delete} & {\ttfamily\scriptsize api\_data\_category\_delete} \\
{\scriptsize\sffamily POST} & {\ttfamily\scriptsize /\allowbreak api/\allowbreak data/\allowbreak categories/\allowbreak <name>/\allowbreak empty} & {\ttfamily\scriptsize api\_data\_category\_empty} \\
\rigaTenue {\scriptsize\sffamily GET} & {\ttfamily\scriptsize /\allowbreak api/\allowbreak data/\allowbreak categories/\allowbreak <name>/\allowbreak files} & {\ttfamily\scriptsize api\_data\_category\_files} \\
{\scriptsize\sffamily POST} & {\ttfamily\scriptsize /\allowbreak api/\allowbreak data/\allowbreak categories/\allowbreak <name>/\allowbreak kind} & {\ttfamily\scriptsize api\_data\_category\_kind} \\
\rigaTenue {\scriptsize\sffamily POST} & {\ttfamily\scriptsize /\allowbreak api/\allowbreak data/\allowbreak download/\allowbreak <mode>} & {\ttfamily\scriptsize api\_data\_download} \\
{\scriptsize\sffamily POST} & {\ttfamily\scriptsize /\allowbreak api/\allowbreak data/\allowbreak mms/\allowbreak generate} & {\ttfamily\scriptsize api\_data\_mms\_generate} \\
\rigaTenue {\scriptsize\sffamily POST} & {\ttfamily\scriptsize /\allowbreak api/\allowbreak data/\allowbreak upload} & {\ttfamily\scriptsize api\_data\_upload} \\
{\scriptsize\sffamily GET} & {\ttfamily\scriptsize /\allowbreak api/\allowbreak ensemble/\allowbreak <name>} & {\ttfamily\scriptsize api\_ensemble} \\
\rigaTenue {\scriptsize\sffamily POST} & {\ttfamily\scriptsize /\allowbreak api/\allowbreak ensemble\_custom} & {\ttfamily\scriptsize api\_ensemble\_custom} \\
{\scriptsize\sffamily GET} & {\ttfamily\scriptsize /\allowbreak api/\allowbreak health} & {\ttfamily\scriptsize api\_health} \\
\rigaTenue {\scriptsize\sffamily GET} & {\ttfamily\scriptsize /\allowbreak api/\allowbreak job/\allowbreak <job\_id>} & {\ttfamily\scriptsize api\_job} \\
{\scriptsize\sffamily POST} & {\ttfamily\scriptsize /\allowbreak api/\allowbreak job/\allowbreak <job\_id>/\allowbreak cancel} & {\ttfamily\scriptsize api\_job\_cancel} \\
\rigaTenue {\scriptsize\sffamily GET} & {\ttfamily\scriptsize /\allowbreak api/\allowbreak job/\allowbreak <job\_id>/\allowbreak stream} & {\ttfamily\scriptsize api\_job\_stream} \\
{\scriptsize\sffamily GET} & {\ttfamily\scriptsize /\allowbreak api/\allowbreak jobs} & {\ttfamily\scriptsize api\_jobs\_list} \\
\rigaTenue {\scriptsize\sffamily DELETE} & {\ttfamily\scriptsize /\allowbreak api/\allowbreak jobs/\allowbreak <job\_id>} & {\ttfamily\scriptsize api\_job\_rimuovi} \\
{\scriptsize\sffamily DELETE} & {\ttfamily\scriptsize /\allowbreak api/\allowbreak jobs/\allowbreak conclusi} & {\ttfamily\scriptsize api\_jobs\_svuota} \\
\rigaTenue {\scriptsize\sffamily GET} & {\ttfamily\scriptsize /\allowbreak api/\allowbreak live/\allowbreak approaches} & {\ttfamily\scriptsize api\_live\_approaches} \\
{\scriptsize\sffamily GET} & {\ttfamily\scriptsize /\allowbreak api/\allowbreak live/\allowbreak ensemble\_modes} & {\ttfamily\scriptsize api\_live\_ensemble\_modes} \\
\rigaTenue {\scriptsize\sffamily POST} & {\ttfamily\scriptsize /\allowbreak api/\allowbreak live/\allowbreak run\_cycle} & {\ttfamily\scriptsize api\_live\_run\_cycle} \\
{\scriptsize\sffamily POST} & {\ttfamily\scriptsize /\allowbreak api/\allowbreak live/\allowbreak start} & {\ttfamily\scriptsize api\_live\_start} \\
\rigaTenue {\scriptsize\sffamily GET} & {\ttfamily\scriptsize /\allowbreak api/\allowbreak live/\allowbreak state} & {\ttfamily\scriptsize api\_live\_state} \\
{\scriptsize\sffamily GET} & {\ttfamily\scriptsize /\allowbreak api/\allowbreak live/\allowbreak status} & {\ttfamily\scriptsize api\_live\_status} \\
\rigaTenue {\scriptsize\sffamily POST} & {\ttfamily\scriptsize /\allowbreak api/\allowbreak live/\allowbreak stop} & {\ttfamily\scriptsize api\_live\_stop} \\
{\scriptsize\sffamily GET} & {\ttfamily\scriptsize /\allowbreak api/\allowbreak live/\allowbreak strategies} & {\ttfamily\scriptsize api\_live\_strategies} \\
\rigaTenue {\scriptsize\sffamily GET POST} & {\ttfamily\scriptsize /\allowbreak api/\allowbreak logics} & {\ttfamily\scriptsize api\_logic\_create, api\_logics} \\
{\scriptsize\sffamily GET POST DELETE} & {\ttfamily\scriptsize /\allowbreak api/\allowbreak logics/\allowbreak <key>/\allowbreak file} & {\ttfamily\scriptsize api\_logic\_file} \\
\rigaTenue {\scriptsize\sffamily GET} & {\ttfamily\scriptsize /\allowbreak api/\allowbreak maintenance/\allowbreak plan} & {\ttfamily\scriptsize api\_maintenance\_plan} \\
{\scriptsize\sffamily GET} & {\ttfamily\scriptsize /\allowbreak api/\allowbreak markets} & {\ttfamily\scriptsize api\_markets} \\
\rigaTenue {\scriptsize\sffamily POST} & {\ttfamily\scriptsize /\allowbreak api/\allowbreak markets/\allowbreak <key>/\allowbreak config} & {\ttfamily\scriptsize api\_market\_config} \\
{\scriptsize\sffamily GET POST} & {\ttfamily\scriptsize /\allowbreak api/\allowbreak markets/\allowbreak <key>/\allowbreak file} & {\ttfamily\scriptsize api\_market\_file} \\
\rigaTenue {\scriptsize\sffamily GET POST DELETE} & {\ttfamily\scriptsize /\allowbreak api/\allowbreak metrics} & {\ttfamily\scriptsize api\_metrics} \\
{\scriptsize\sffamily POST} & {\ttfamily\scriptsize /\allowbreak api/\allowbreak metrics/\allowbreak compute} & {\ttfamily\scriptsize api\_metrics\_compute} \\
\rigaTenue {\scriptsize\sffamily GET} & {\ttfamily\scriptsize /\allowbreak api/\allowbreak notifications} & {\ttfamily\scriptsize api\_notifications} \\
{\scriptsize\sffamily GET POST} & {\ttfamily\scriptsize /\allowbreak api/\allowbreak overview/\allowbreak prefs} & {\ttfamily\scriptsize api\_overview\_prefs} \\
\rigaTenue {\scriptsize\sffamily GET} & {\ttfamily\scriptsize /\allowbreak api/\allowbreak plugins} & {\ttfamily\scriptsize api\_plugins} \\
{\scriptsize\sffamily GET POST} & {\ttfamily\scriptsize /\allowbreak api/\allowbreak projects} & {\ttfamily\scriptsize api\_projects} \\
\rigaTenue {\scriptsize\sffamily GET POST DELETE} & {\ttfamily\scriptsize /\allowbreak api/\allowbreak projects/\allowbreak <pid>} & {\ttfamily\scriptsize api\_project} \\
{\scriptsize\sffamily GET} & {\ttfamily\scriptsize /\allowbreak api/\allowbreak regime} & {\ttfamily\scriptsize api\_regime} \\
\rigaTenue {\scriptsize\sffamily POST} & {\ttfamily\scriptsize /\allowbreak api/\allowbreak reload} & {\ttfamily\scriptsize api\_reload} \\
{\scriptsize\sffamily POST} & {\ttfamily\scriptsize /\allowbreak api/\allowbreak retrain} & {\ttfamily\scriptsize api\_retrain} \\
\rigaTenue {\scriptsize\sffamily POST} & {\ttfamily\scriptsize /\allowbreak api/\allowbreak run/\allowbreak <strategy>} & {\ttfamily\scriptsize api\_run\_strategy} \\
{\scriptsize\sffamily POST} & {\ttfamily\scriptsize /\allowbreak api/\allowbreak run\_all} & {\ttfamily\scriptsize api\_run\_all} \\
\rigaTenue {\scriptsize\sffamily GET POST} & {\ttfamily\scriptsize /\allowbreak api/\allowbreak runners} & {\ttfamily\scriptsize api\_runners\_create, api\_runners\_list} \\
{\scriptsize\sffamily DELETE} & {\ttfamily\scriptsize /\allowbreak api/\allowbreak runners/\allowbreak <rid>} & {\ttfamily\scriptsize api\_runners\_delete} \\
\rigaTenue {\scriptsize\sffamily GET} & {\ttfamily\scriptsize /\allowbreak api/\allowbreak runners/\allowbreak <rid>/\allowbreak simulation} & {\ttfamily\scriptsize api\_runner\_simulation} \\
{\scriptsize\sffamily GET} & {\ttfamily\scriptsize /\allowbreak api/\allowbreak runners/\allowbreak <rid>/\allowbreak simulation/\allowbreak series/\allowbreak <symbol>} & {\ttfamily\scriptsize api\_runner\_simulation\_series} \\
\rigaTenue {\scriptsize\sffamily POST} & {\ttfamily\scriptsize /\allowbreak api/\allowbreak runners/\allowbreak <rid>/\allowbreak start} & {\ttfamily\scriptsize api\_runners\_start} \\
{\scriptsize\sffamily GET} & {\ttfamily\scriptsize /\allowbreak api/\allowbreak runners/\allowbreak <rid>/\allowbreak status} & {\ttfamily\scriptsize api\_runners\_status} \\
\rigaTenue {\scriptsize\sffamily POST} & {\ttfamily\scriptsize /\allowbreak api/\allowbreak runners/\allowbreak <rid>/\allowbreak stop} & {\ttfamily\scriptsize api\_runners\_stop} \\
{\scriptsize\sffamily GET} & {\ttfamily\scriptsize /\allowbreak api/\allowbreak strategies} & {\ttfamily\scriptsize api\_strategies} \\
\rigaTenue {\scriptsize\sffamily GET} & {\ttfamily\scriptsize /\allowbreak api/\allowbreak stream/\allowbreak <name>} & {\ttfamily\scriptsize api\_stream} \\
{\scriptsize\sffamily POST} & {\ttfamily\scriptsize /\allowbreak api/\allowbreak stream/\allowbreak <name>/\allowbreak delete} & {\ttfamily\scriptsize api\_stream\_delete} \\
\rigaTenue {\scriptsize\sffamily GET} & {\ttfamily\scriptsize /\allowbreak api/\allowbreak stream/\allowbreak <name>/\allowbreak files} & {\ttfamily\scriptsize api\_stream\_files} \\
{\scriptsize\sffamily GET} & {\ttfamily\scriptsize /\allowbreak api/\allowbreak streams} & {\ttfamily\scriptsize api\_streams\_list} \\
\rigaTenue {\scriptsize\sffamily GET} & {\ttfamily\scriptsize /\allowbreak api/\allowbreak streams/\allowbreak orfani} & {\ttfamily\scriptsize api\_streams\_orphans} \\
{\scriptsize\sffamily POST} & {\ttfamily\scriptsize /\allowbreak api/\allowbreak streams/\allowbreak orfani/\allowbreak pulisci} & {\ttfamily\scriptsize api\_streams\_orphans\_clean} \\
\rigaTenue {\scriptsize\sffamily GET} & {\ttfamily\scriptsize /\allowbreak api/\allowbreak universo} & {\ttfamily\scriptsize api\_universo} \\
{\scriptsize\sffamily POST} & {\ttfamily\scriptsize /\allowbreak api/\allowbreak universo/\allowbreak <symbol>/\allowbreak categoria} & {\ttfamily\scriptsize api\_universo\_categoria} \\
\rigaTenue {\scriptsize\sffamily GET POST} & {\ttfamily\scriptsize /\allowbreak api/\allowbreak user\_approaches} & {\ttfamily\scriptsize api\_user\_approach\_create, api\_user\_approaches\_list} \\
{\scriptsize\sffamily GET DELETE} & {\ttfamily\scriptsize /\allowbreak api/\allowbreak user\_approaches/\allowbreak <name>} & {\ttfamily\scriptsize api\_user\_approach\_delete, api\_user\_approach\_files} \\
\rigaTenue {\scriptsize\sffamily GET POST DELETE} & {\ttfamily\scriptsize /\allowbreak api/\allowbreak user\_approaches/\allowbreak <name>/\allowbreak <filename>} & {\ttfamily\scriptsize api\_user\_approach\_file\_delete, api\_user\_approach\_file\_get, api\_user\_approach\_file\_save} \\
{\scriptsize\sffamily GET POST} & {\ttfamily\scriptsize /\allowbreak api/\allowbreak user\_strategies} & {\ttfamily\scriptsize api\_user\_strategies\_list, api\_user\_strategy\_create} \\
\rigaTenue {\scriptsize\sffamily GET DELETE} & {\ttfamily\scriptsize /\allowbreak api/\allowbreak user\_strategies/\allowbreak <name>} & {\ttfamily\scriptsize api\_user\_strategy\_delete, api\_user\_strategy\_files} \\
{\scriptsize\sffamily GET POST DELETE} & {\ttfamily\scriptsize /\allowbreak api/\allowbreak user\_strategies/\allowbreak <name>/\allowbreak <filename>} & {\ttfamily\scriptsize api\_user\_file\_delete, api\_user\_file\_get, api\_user\_file\_save} \\
\rigaTenue {\scriptsize\sffamily GET} & {\ttfamily\scriptsize /\allowbreak api/\allowbreak visual/\allowbreak <strategy>/\allowbreak chart/\allowbreak <symbol>} & {\ttfamily\scriptsize api\_visual\_chart} \\
{\scriptsize\sffamily POST} & {\ttfamily\scriptsize /\allowbreak api/\allowbreak visual/\allowbreak <strategy>/\allowbreak custom\_indicator/\allowbreak <symbol>} & {\ttfamily\scriptsize api\_visual\_custom\_indicator} \\
\rigaTenue {\scriptsize\sffamily GET} & {\ttfamily\scriptsize /\allowbreak api/\allowbreak visual/\allowbreak <strategy>/\allowbreak indicators/\allowbreak <symbol>} & {\ttfamily\scriptsize api\_visual\_indicators} \\
{\scriptsize\sffamily GET} & {\ttfamily\scriptsize /\allowbreak api/\allowbreak visual/\allowbreak <strategy>/\allowbreak liquidity\_heatmap/\allowbreak <symbol>} & {\ttfamily\scriptsize api\_visual\_liquidity\_heatmap} \\
\rigaTenue {\scriptsize\sffamily GET} & {\ttfamily\scriptsize /\allowbreak api/\allowbreak visual/\allowbreak <strategy>/\allowbreak mms/\allowbreak <symbol>} & {\ttfamily\scriptsize api\_visual\_mms} \\
{\scriptsize\sffamily GET} & {\ttfamily\scriptsize /\allowbreak api/\allowbreak visual/\allowbreak <strategy>/\allowbreak symbols} & {\ttfamily\scriptsize api\_visual\_symbols} \\
\rigaTenue {\scriptsize\sffamily GET} & {\ttfamily\scriptsize /\allowbreak api/\allowbreak visual/\allowbreak <strategy>/\allowbreak zones/\allowbreak <symbol>} & {\ttfamily\scriptsize api\_visual\_zones} \\
{\scriptsize\sffamily GET} & {\ttfamily\scriptsize /\allowbreak api/\allowbreak visual/\allowbreak indicator\_template} & {\ttfamily\scriptsize api\_visual\_indicator\_template} \\
\rigaTenue {\scriptsize\sffamily GET POST} & {\ttfamily\scriptsize /\allowbreak api/\allowbreak visual/\allowbreak indicators} & {\ttfamily\scriptsize api\_visual\_indicators\_list, api\_visual\_indicators\_save} \\
{\scriptsize\sffamily POST} & {\ttfamily\scriptsize /\allowbreak api/\allowbreak visual/\allowbreak indicators/\allowbreak delete} & {\ttfamily\scriptsize api\_visual\_indicators\_delete} \\
\rigaTenue {\scriptsize\sffamily GET} & {\ttfamily\scriptsize /\allowbreak config} & {\ttfamily\scriptsize config\_page} \\
{\scriptsize\sffamily GET} & {\ttfamily\scriptsize /\allowbreak data} & {\ttfamily\scriptsize data\_page} \\
\rigaTenue {\scriptsize\sffamily GET} & {\ttfamily\scriptsize /\allowbreak editor} & {\ttfamily\scriptsize editor\_page} \\
{\scriptsize\sffamily GET} & {\ttfamily\scriptsize /\allowbreak ensemble/\allowbreak <name>} & {\ttfamily\scriptsize ensemble\_page} \\
\rigaTenue {\scriptsize\sffamily GET} & {\ttfamily\scriptsize /\allowbreak ensemble\_custom} & {\ttfamily\scriptsize ensemble\_custom\_page} \\
{\scriptsize\sffamily GET} & {\ttfamily\scriptsize /\allowbreak favicon.ico} & {\ttfamily\scriptsize favicon} \\
\rigaTenue {\scriptsize\sffamily GET} & {\ttfamily\scriptsize /\allowbreak live/\allowbreak test/\allowbreak <rid>} & {\ttfamily\scriptsize live\_test\_page} \\
{\scriptsize\sffamily GET} & {\ttfamily\scriptsize /\allowbreak profile} & {\ttfamily\scriptsize profile\_page} \\
\rigaTenue {\scriptsize\sffamily GET} & {\ttfamily\scriptsize /\allowbreak projects} & {\ttfamily\scriptsize projects\_page} \\
{\scriptsize\sffamily GET} & {\ttfamily\scriptsize /\allowbreak regime} & {\ttfamily\scriptsize regime\_page} \\
\rigaTenue {\scriptsize\sffamily GET} & {\ttfamily\scriptsize /\allowbreak run} & {\ttfamily\scriptsize run\_page} \\
{\scriptsize\sffamily GET} & {\ttfamily\scriptsize /\allowbreak security} & {\ttfamily\scriptsize security\_page} \\
\rigaTenue {\scriptsize\sffamily GET} & {\ttfamily\scriptsize /\allowbreak stream/\allowbreak <name>} & {\ttfamily\scriptsize stream\_page} \\
{\scriptsize\sffamily GET} & {\ttfamily\scriptsize /\allowbreak streams} & {\ttfamily\scriptsize streams\_page} \\
\rigaTenue {\scriptsize\sffamily GET} & {\ttfamily\scriptsize /\allowbreak subscription} & {\ttfamily\scriptsize subscription\_page} \\
{\scriptsize\sffamily GET} & {\ttfamily\scriptsize /\allowbreak visual} & {\ttfamily\scriptsize visual\_page} \\
\rigaTenue {\scriptsize\sffamily GET} & {\ttfamily\scriptsize /\allowbreak visual/\allowbreak <strategy>} & {\ttfamily\scriptsize visual\_page} \\
\bottomrule
\end{longtable}

\subsection{studio.py --- blueprint \cod{studio}}
\begin{longtable}{@{}L{17mm}L{74mm}L{45mm}@{}}
\toprule
\textbf{Metodi} & \textbf{Percorso} & \textbf{Funzione} \\
\midrule\endfirsthead
\toprule
\textbf{Metodi} & \textbf{Percorso} & \textbf{Funzione} \\
\midrule\endhead
{\scriptsize\sffamily GET} & {\ttfamily\scriptsize /\allowbreak studio} & {\ttfamily\scriptsize page} \\
\rigaTenue {\scriptsize\sffamily POST} & {\ttfamily\scriptsize /\allowbreak studio/\allowbreak api/\allowbreak approach/\allowbreak install} & {\ttfamily\scriptsize api\_approach\_install} \\
{\scriptsize\sffamily GET} & {\ttfamily\scriptsize /\allowbreak studio/\allowbreak api/\allowbreak available/\allowbreak approaches} & {\ttfamily\scriptsize api\_avail\_approaches} \\
\rigaTenue {\scriptsize\sffamily GET} & {\ttfamily\scriptsize /\allowbreak studio/\allowbreak api/\allowbreak available/\allowbreak metrics} & {\ttfamily\scriptsize api\_avail\_metrics} \\
{\scriptsize\sffamily GET} & {\ttfamily\scriptsize /\allowbreak studio/\allowbreak api/\allowbreak available/\allowbreak overlays} & {\ttfamily\scriptsize api\_avail\_overlays} \\
\rigaTenue {\scriptsize\sffamily GET} & {\ttfamily\scriptsize /\allowbreak studio/\allowbreak api/\allowbreak available/\allowbreak strategies} & {\ttfamily\scriptsize api\_avail\_strategies} \\
{\scriptsize\sffamily GET} & {\ttfamily\scriptsize /\allowbreak studio/\allowbreak api/\allowbreak available/\allowbreak weightmodes} & {\ttfamily\scriptsize api\_avail\_weightmodes} \\
\rigaTenue {\scriptsize\sffamily POST} & {\ttfamily\scriptsize /\allowbreak studio/\allowbreak api/\allowbreak compile} & {\ttfamily\scriptsize api\_compile} \\
{\scriptsize\sffamily POST} & {\ttfamily\scriptsize /\allowbreak studio/\allowbreak api/\allowbreak data/\allowbreak adapter/\allowbreak apply} & {\ttfamily\scriptsize api\_data\_apply} \\
\rigaTenue {\scriptsize\sffamily POST} & {\ttfamily\scriptsize /\allowbreak studio/\allowbreak api/\allowbreak data/\allowbreak adapter/\allowbreak generate} & {\ttfamily\scriptsize api\_data\_generate} \\
{\scriptsize\sffamily POST} & {\ttfamily\scriptsize /\allowbreak studio/\allowbreak api/\allowbreak data/\allowbreak analyze} & {\ttfamily\scriptsize api\_data\_analyze} \\
\rigaTenue {\scriptsize\sffamily POST} & {\ttfamily\scriptsize /\allowbreak studio/\allowbreak api/\allowbreak fidelity/\allowbreak codify} & {\ttfamily\scriptsize api\_fid\_codify} \\
{\scriptsize\sffamily POST} & {\ttfamily\scriptsize /\allowbreak studio/\allowbreak api/\allowbreak fidelity/\allowbreak score} & {\ttfamily\scriptsize api\_fid\_score} \\
\rigaTenue {\scriptsize\sffamily POST} & {\ttfamily\scriptsize /\allowbreak studio/\allowbreak api/\allowbreak generate} & {\ttfamily\scriptsize api\_generate} \\
{\scriptsize\sffamily POST} & {\ttfamily\scriptsize /\allowbreak studio/\allowbreak api/\allowbreak indicator/\allowbreak compile} & {\ttfamily\scriptsize api\_ind\_compile} \\
\rigaTenue {\scriptsize\sffamily POST} & {\ttfamily\scriptsize /\allowbreak studio/\allowbreak api/\allowbreak indicator/\allowbreak generate} & {\ttfamily\scriptsize api\_ind\_generate} \\
{\scriptsize\sffamily POST} & {\ttfamily\scriptsize /\allowbreak studio/\allowbreak api/\allowbreak install} & {\ttfamily\scriptsize api\_install} \\
\rigaTenue {\scriptsize\sffamily GET} & {\ttfamily\scriptsize /\allowbreak studio/\allowbreak api/\allowbreak installati/\allowbreak <categoria>} & {\ttfamily\scriptsize api\_installati} \\
{\scriptsize\sffamily GET} & {\ttfamily\scriptsize /\allowbreak studio/\allowbreak api/\allowbreak installati/\allowbreak <categoria>/\allowbreak <nome>} & {\ttfamily\scriptsize api\_installato\_apri} \\
\rigaTenue {\scriptsize\sffamily POST} & {\ttfamily\scriptsize /\allowbreak studio/\allowbreak api/\allowbreak installati/\allowbreak <categoria>/\allowbreak <nome>/\allowbreak elimina} & {\ttfamily\scriptsize api\_installato\_elimina} \\
{\scriptsize\sffamily GET POST} & {\ttfamily\scriptsize /\allowbreak studio/\allowbreak api/\allowbreak manual-trades} & {\ttfamily\scriptsize api\_manual\_trades} \\
\rigaTenue {\scriptsize\sffamily POST} & {\ttfamily\scriptsize /\allowbreak studio/\allowbreak api/\allowbreak metric/\allowbreak compile} & {\ttfamily\scriptsize api\_metric\_compile} \\
{\scriptsize\sffamily GET} & {\ttfamily\scriptsize /\allowbreak studio/\allowbreak api/\allowbreak metric/\allowbreak doc} & {\ttfamily\scriptsize api\_metric\_doc} \\
\rigaTenue {\scriptsize\sffamily POST} & {\ttfamily\scriptsize /\allowbreak studio/\allowbreak api/\allowbreak metric/\allowbreak generate} & {\ttfamily\scriptsize api\_metric\_generate} \\
{\scriptsize\sffamily GET} & {\ttfamily\scriptsize /\allowbreak studio/\allowbreak api/\allowbreak metric/\allowbreak template} & {\ttfamily\scriptsize api\_metric\_template} \\
\rigaTenue {\scriptsize\sffamily POST} & {\ttfamily\scriptsize /\allowbreak studio/\allowbreak api/\allowbreak metric/\allowbreak validate} & {\ttfamily\scriptsize api\_metric\_validate} \\
{\scriptsize\sffamily GET} & {\ttfamily\scriptsize /\allowbreak studio/\allowbreak api/\allowbreak models} & {\ttfamily\scriptsize api\_models} \\
\rigaTenue {\scriptsize\sffamily POST} & {\ttfamily\scriptsize /\allowbreak studio/\allowbreak api/\allowbreak overlay/\allowbreak compile} & {\ttfamily\scriptsize api\_overlay\_compile} \\
{\scriptsize\sffamily POST} & {\ttfamily\scriptsize /\allowbreak studio/\allowbreak api/\allowbreak overlay/\allowbreak generate} & {\ttfamily\scriptsize api\_overlay\_generate} \\
\rigaTenue {\scriptsize\sffamily POST} & {\ttfamily\scriptsize /\allowbreak studio/\allowbreak api/\allowbreak overlay/\allowbreak install} & {\ttfamily\scriptsize api\_overlay\_install} \\
{\scriptsize\sffamily GET POST DELETE} & {\ttfamily\scriptsize /\allowbreak studio/\allowbreak api/\allowbreak overlay/\allowbreak port} & {\ttfamily\scriptsize api\_overlay\_port} \\
\rigaTenue {\scriptsize\sffamily GET} & {\ttfamily\scriptsize /\allowbreak studio/\allowbreak api/\allowbreak overlay/\allowbreak template} & {\ttfamily\scriptsize api\_overlay\_template} \\
{\scriptsize\sffamily GET} & {\ttfamily\scriptsize /\allowbreak studio/\allowbreak api/\allowbreak runner/\allowbreak get/\allowbreak <name>} & {\ttfamily\scriptsize api\_runner\_get} \\
\rigaTenue {\scriptsize\sffamily GET} & {\ttfamily\scriptsize /\allowbreak studio/\allowbreak api/\allowbreak runner/\allowbreak list} & {\ttfamily\scriptsize api\_runner\_list} \\
{\scriptsize\sffamily POST} & {\ttfamily\scriptsize /\allowbreak studio/\allowbreak api/\allowbreak runner/\allowbreak save} & {\ttfamily\scriptsize api\_runner\_save} \\
\rigaTenue {\scriptsize\sffamily POST} & {\ttfamily\scriptsize /\allowbreak studio/\allowbreak api/\allowbreak runner/\allowbreak validate} & {\ttfamily\scriptsize api\_runner\_validate} \\
{\scriptsize\sffamily POST} & {\ttfamily\scriptsize /\allowbreak studio/\allowbreak api/\allowbreak weightmode/\allowbreak compile} & {\ttfamily\scriptsize api\_wm\_compile} \\
\rigaTenue {\scriptsize\sffamily POST} & {\ttfamily\scriptsize /\allowbreak studio/\allowbreak api/\allowbreak weightmode/\allowbreak generate} & {\ttfamily\scriptsize api\_wm\_generate} \\
{\scriptsize\sffamily POST} & {\ttfamily\scriptsize /\allowbreak studio/\allowbreak api/\allowbreak weightmode/\allowbreak install} & {\ttfamily\scriptsize api\_wm\_install} \\
\rigaTenue {\scriptsize\sffamily GET POST DELETE} & {\ttfamily\scriptsize /\allowbreak studio/\allowbreak api/\allowbreak weightmode/\allowbreak port} & {\ttfamily\scriptsize api\_wm\_port} \\
{\scriptsize\sffamily GET} & {\ttfamily\scriptsize /\allowbreak studio/\allowbreak api/\allowbreak weightmode/\allowbreak template} & {\ttfamily\scriptsize api\_wm\_template} \\
\rigaTenue {\scriptsize\sffamily GET} & {\ttfamily\scriptsize /\allowbreak studio/\allowbreak dati-esterni} & {\ttfamily\scriptsize data\_esterni\_page} \\
{\scriptsize\sffamily GET} & {\ttfamily\scriptsize /\allowbreak studio/\allowbreak runner-live} & {\ttfamily\scriptsize runner\_live\_page} \\
\bottomrule
\end{longtable}

\subsection{community.py --- blueprint \cod{community}}
\begin{longtable}{@{}L{17mm}L{74mm}L{45mm}@{}}
\toprule
\textbf{Metodi} & \textbf{Percorso} & \textbf{Funzione} \\
\midrule\endfirsthead
\toprule
\textbf{Metodi} & \textbf{Percorso} & \textbf{Funzione} \\
\midrule\endhead
{\scriptsize\sffamily GET} & {\ttfamily\scriptsize /\allowbreak account} & {\ttfamily\scriptsize account} \\
\rigaTenue {\scriptsize\sffamily GET} & {\ttfamily\scriptsize /\allowbreak api/\allowbreak abbonamento} & {\ttfamily\scriptsize abbonamento\_stato} \\
{\scriptsize\sffamily POST} & {\ttfamily\scriptsize /\allowbreak api/\allowbreak abbonamento/\allowbreak checkout} & {\ttfamily\scriptsize abbonamento\_checkout} \\
\rigaTenue {\scriptsize\sffamily POST} & {\ttfamily\scriptsize /\allowbreak api/\allowbreak abbonamento/\allowbreak portale} & {\ttfamily\scriptsize abbonamento\_portale} \\
{\scriptsize\sffamily POST} & {\ttfamily\scriptsize /\allowbreak api/\allowbreak account} & {\ttfamily\scriptsize account\_update} \\
\rigaTenue {\scriptsize\sffamily POST} & {\ttfamily\scriptsize /\allowbreak api/\allowbreak community/\allowbreak feedback} & {\ttfamily\scriptsize feedback\_create} \\
{\scriptsize\sffamily POST} & {\ttfamily\scriptsize /\allowbreak api/\allowbreak feedback} & {\ttfamily\scriptsize api\_submit\_feedback} \\
\rigaTenue {\scriptsize\sffamily GET} & {\ttfamily\scriptsize /\allowbreak api/\allowbreak marketplace/\allowbreak <artifact\_id>} & {\ttfamily\scriptsize marketplace\_detail} \\
{\scriptsize\sffamily POST} & {\ttfamily\scriptsize /\allowbreak api/\allowbreak marketplace/\allowbreak <artifact\_id>/\allowbreak buy} & {\ttfamily\scriptsize marketplace\_buy} \\
\rigaTenue {\scriptsize\sffamily POST} & {\ttfamily\scriptsize /\allowbreak api/\allowbreak marketplace/\allowbreak <artifact\_id>/\allowbreak comments} & {\ttfamily\scriptsize marketplace\_comment} \\
{\scriptsize\sffamily GET} & {\ttfamily\scriptsize /\allowbreak api/\allowbreak marketplace/\allowbreak <artifact\_id>/\allowbreak conflicts} & {\ttfamily\scriptsize marketplace\_conflicts} \\
\rigaTenue {\scriptsize\sffamily POST} & {\ttfamily\scriptsize /\allowbreak api/\allowbreak marketplace/\allowbreak <artifact\_id>/\allowbreak contesta} & {\ttfamily\scriptsize marketplace\_contesta} \\
{\scriptsize\sffamily POST} & {\ttfamily\scriptsize /\allowbreak api/\allowbreak marketplace/\allowbreak <artifact\_id>/\allowbreak install} & {\ttfamily\scriptsize marketplace\_install} \\
\rigaTenue {\scriptsize\sffamily POST} & {\ttfamily\scriptsize /\allowbreak api/\allowbreak marketplace/\allowbreak <artifact\_id>/\allowbreak installa-admin} & {\ttfamily\scriptsize marketplace\_installa\_admin} \\
{\scriptsize\sffamily GET} & {\ttfamily\scriptsize /\allowbreak api/\allowbreak marketplace/\allowbreak <artifact\_id>/\allowbreak issue-stato} & {\ttfamily\scriptsize marketplace\_issue\_stato} \\
\rigaTenue {\scriptsize\sffamily GET} & {\ttfamily\scriptsize /\allowbreak api/\allowbreak marketplace/\allowbreak <artifact\_id>/\allowbreak perche-non-mio} & {\ttfamily\scriptsize marketplace\_perche\_non\_mio} \\
{\scriptsize\sffamily POST} & {\ttfamily\scriptsize /\allowbreak api/\allowbreak marketplace/\allowbreak <artifact\_id>/\allowbreak remove} & {\ttfamily\scriptsize marketplace\_remove} \\
\rigaTenue {\scriptsize\sffamily POST} & {\ttfamily\scriptsize /\allowbreak api/\allowbreak marketplace/\allowbreak <artifact\_id>/\allowbreak riconcilia} & {\ttfamily\scriptsize marketplace\_riconcilia} \\
{\scriptsize\sffamily POST} & {\ttfamily\scriptsize /\allowbreak api/\allowbreak marketplace/\allowbreak <artifact\_id>/\allowbreak scheda} & {\ttfamily\scriptsize marketplace\_modifica\_scheda} \\
\rigaTenue {\scriptsize\sffamily GET} & {\ttfamily\scriptsize /\allowbreak api/\allowbreak marketplace/\allowbreak acquisti} & {\ttfamily\scriptsize marketplace\_acquisti} \\
{\scriptsize\sffamily GET} & {\ttfamily\scriptsize /\allowbreak api/\allowbreak marketplace/\allowbreak claim/\allowbreak <claim\_id>} & {\ttfamily\scriptsize marketplace\_claim} \\
\rigaTenue {\scriptsize\sffamily POST} & {\ttfamily\scriptsize /\allowbreak api/\allowbreak marketplace/\allowbreak claim/\allowbreak <claim\_id>/\allowbreak chiudi} & {\ttfamily\scriptsize marketplace\_claim\_chiudi} \\
{\scriptsize\sffamily POST} & {\ttfamily\scriptsize /\allowbreak api/\allowbreak marketplace/\allowbreak claim/\allowbreak <claim\_id>/\allowbreak messaggio} & {\ttfamily\scriptsize marketplace\_claim\_messaggio} \\
\rigaTenue {\scriptsize\sffamily POST} & {\ttfamily\scriptsize /\allowbreak api/\allowbreak marketplace/\allowbreak claim/\allowbreak <claim\_id>/\allowbreak verifica-integrita} & {\ttfamily\scriptsize marketplace\_claim\_integrita} \\
{\scriptsize\sffamily POST} & {\ttfamily\scriptsize /\allowbreak api/\allowbreak marketplace/\allowbreak claim/\allowbreak <claim\_id>/\allowbreak verifica-performance} & {\ttfamily\scriptsize marketplace\_claim\_performance} \\
\rigaTenue {\scriptsize\sffamily GET} & {\ttfamily\scriptsize /\allowbreak api/\allowbreak marketplace/\allowbreak claim/\allowbreak aperte} & {\ttfamily\scriptsize marketplace\_claim\_aperte} \\
{\scriptsize\sffamily GET} & {\ttfamily\scriptsize /\allowbreak api/\allowbreak marketplace/\allowbreak contestazioni} & {\ttfamily\scriptsize marketplace\_contestazioni} \\
\rigaTenue {\scriptsize\sffamily GET} & {\ttfamily\scriptsize /\allowbreak api/\allowbreak marketplace/\allowbreak diagnostica} & {\ttfamily\scriptsize marketplace\_diagnostica} \\
{\scriptsize\sffamily GET} & {\ttfamily\scriptsize /\allowbreak api/\allowbreak marketplace/\allowbreak items} & {\ttfamily\scriptsize marketplace\_items\_api} \\
\rigaTenue {\scriptsize\sffamily GET} & {\ttfamily\scriptsize /\allowbreak api/\allowbreak marketplace/\allowbreak mine} & {\ttfamily\scriptsize marketplace\_mine} \\
{\scriptsize\sffamily POST} & {\ttfamily\scriptsize /\allowbreak api/\allowbreak marketplace/\allowbreak publish} & {\ttfamily\scriptsize marketplace\_publish} \\
\rigaTenue {\scriptsize\sffamily GET} & {\ttfamily\scriptsize /\allowbreak api/\allowbreak marketplace/\allowbreak system-artifacts} & {\ttfamily\scriptsize marketplace\_system\_artifacts} \\
{\scriptsize\sffamily GET} & {\ttfamily\scriptsize /\allowbreak api/\allowbreak marketplace/\allowbreak system-artifacts/\allowbreak <kind>/\allowbreak <path:name>} & {\ttfamily\scriptsize marketplace\_system\_artifact} \\
\rigaTenue {\scriptsize\sffamily GET} & {\ttfamily\scriptsize /\allowbreak api/\allowbreak marketplace/\allowbreak video-check} & {\ttfamily\scriptsize marketplace\_video\_check} \\
{\scriptsize\sffamily GET POST} & {\ttfamily\scriptsize /\allowbreak api/\allowbreak messages/\allowbreak <conversation\_id>} & {\ttfamily\scriptsize messages\_send, messages\_thread} \\
\rigaTenue {\scriptsize\sffamily GET} & {\ttfamily\scriptsize /\allowbreak feedback} & {\ttfamily\scriptsize feedback} \\
{\scriptsize\sffamily GET} & {\ttfamily\scriptsize /\allowbreak marketplace} & {\ttfamily\scriptsize marketplace} \\
\rigaTenue {\scriptsize\sffamily GET} & {\ttfamily\scriptsize /\allowbreak marketplace/\allowbreak <artifact\_id>} & {\ttfamily\scriptsize marketplace\_item} \\
{\scriptsize\sffamily GET} & {\ttfamily\scriptsize /\allowbreak marketplace/\allowbreak acquisto/\allowbreak <artifact\_id>} & {\ttfamily\scriptsize marketplace\_purchase\_return} \\
\rigaTenue {\scriptsize\sffamily GET} & {\ttfamily\scriptsize /\allowbreak marketplace/\allowbreak rimborso/\allowbreak <claim\_id>} & {\ttfamily\scriptsize claim\_page} \\
{\scriptsize\sffamily GET} & {\ttfamily\scriptsize /\allowbreak marketplace/\allowbreak sell} & {\ttfamily\scriptsize marketplace\_sell} \\
\rigaTenue {\scriptsize\sffamily GET} & {\ttfamily\scriptsize /\allowbreak messages} & {\ttfamily\scriptsize messages} \\
\bottomrule
\end{longtable}

\subsection{marketplace\_routes.py --- sul blueprint \cod{community}}
\begin{longtable}{@{}L{17mm}L{74mm}L{45mm}@{}}
\toprule
\textbf{Metodi} & \textbf{Percorso} & \textbf{Funzione} \\
\midrule\endfirsthead
\toprule
\textbf{Metodi} & \textbf{Percorso} & \textbf{Funzione} \\
\midrule\endhead
{\scriptsize\sffamily GET} & {\ttfamily\scriptsize /\allowbreak api/\allowbreak market/\allowbreak admin/\allowbreak payments} & {\ttfamily\scriptsize market\_admin\_payments} \\
\rigaTenue {\scriptsize\sffamily POST} & {\ttfamily\scriptsize /\allowbreak api/\allowbreak market/\allowbreak admin/\allowbreak refund} & {\ttfamily\scriptsize market\_admin\_refund} \\
{\scriptsize\sffamily POST} & {\ttfamily\scriptsize /\allowbreak api/\allowbreak market/\allowbreak buy/\allowbreak <artifact\_id>} & {\ttfamily\scriptsize market\_buy} \\
\rigaTenue {\scriptsize\sffamily GET} & {\ttfamily\scriptsize /\allowbreak api/\allowbreak market/\allowbreak catalog} & {\ttfamily\scriptsize market\_catalog} \\
{\scriptsize\sffamily GET} & {\ttfamily\scriptsize /\allowbreak api/\allowbreak market/\allowbreak earnings} & {\ttfamily\scriptsize market\_earnings} \\
\rigaTenue {\scriptsize\sffamily POST} & {\ttfamily\scriptsize /\allowbreak api/\allowbreak market/\allowbreak install/\allowbreak <artifact\_id>} & {\ttfamily\scriptsize market\_install} \\
{\scriptsize\sffamily GET} & {\ttfamily\scriptsize /\allowbreak api/\allowbreak market/\allowbreak item/\allowbreak <artifact\_id>} & {\ttfamily\scriptsize market\_item} \\
\rigaTenue {\scriptsize\sffamily GET} & {\ttfamily\scriptsize /\allowbreak api/\allowbreak market/\allowbreak license/\allowbreak <artifact\_id>} & {\ttfamily\scriptsize market\_license} \\
{\scriptsize\sffamily POST} & {\ttfamily\scriptsize /\allowbreak api/\allowbreak market/\allowbreak publish} & {\ttfamily\scriptsize market\_publish} \\
\rigaTenue {\scriptsize\sffamily GET POST} & {\ttfamily\scriptsize /\allowbreak api/\allowbreak market/\allowbreak seller} & {\ttfamily\scriptsize market\_seller, market\_seller\_azione} \\
{\scriptsize\sffamily POST} & {\ttfamily\scriptsize /\allowbreak api/\allowbreak market/\allowbreak seller/\allowbreak connect} & {\ttfamily\scriptsize market\_seller\_connect} \\
\bottomrule
\end{longtable}

\subsection{agent.py --- blueprint \cod{agent}}
\begin{longtable}{@{}L{17mm}L{74mm}L{45mm}@{}}
\toprule
\textbf{Metodi} & \textbf{Percorso} & \textbf{Funzione} \\
\midrule\endfirsthead
\toprule
\textbf{Metodi} & \textbf{Percorso} & \textbf{Funzione} \\
\midrule\endhead
{\scriptsize\sffamily GET} & {\ttfamily\scriptsize /\allowbreak agent} & {\ttfamily\scriptsize page} \\
\rigaTenue {\scriptsize\sffamily GET POST} & {\ttfamily\scriptsize /\allowbreak agent/\allowbreak api/\allowbreak chats} & {\ttfamily\scriptsize api\_chat\_create, api\_chats} \\
{\scriptsize\sffamily POST} & {\ttfamily\scriptsize /\allowbreak agent/\allowbreak api/\allowbreak chats/\allowbreak <chat\_id>/\allowbreak delete} & {\ttfamily\scriptsize api\_chat\_delete} \\
\rigaTenue {\scriptsize\sffamily POST} & {\ttfamily\scriptsize /\allowbreak agent/\allowbreak api/\allowbreak chats/\allowbreak <chat\_id>/\allowbreak message} & {\ttfamily\scriptsize api\_chat\_message} \\
{\scriptsize\sffamily GET} & {\ttfamily\scriptsize /\allowbreak agent/\allowbreak api/\allowbreak chats/\allowbreak <chat\_id>/\allowbreak messages} & {\ttfamily\scriptsize api\_chat\_messages} \\
\rigaTenue {\scriptsize\sffamily POST} & {\ttfamily\scriptsize /\allowbreak agent/\allowbreak api/\allowbreak chats/\allowbreak <chat\_id>/\allowbreak prepare} & {\ttfamily\scriptsize api\_chat\_prepare} \\
{\scriptsize\sffamily GET} & {\ttfamily\scriptsize /\allowbreak agent/\allowbreak api/\allowbreak chats/\allowbreak <chat\_id>/\allowbreak progress} & {\ttfamily\scriptsize api\_chat\_progress} \\
\rigaTenue {\scriptsize\sffamily POST} & {\ttfamily\scriptsize /\allowbreak agent/\allowbreak api/\allowbreak chats/\allowbreak <chat\_id>/\allowbreak rename} & {\ttfamily\scriptsize api\_chat\_rename} \\
{\scriptsize\sffamily GET} & {\ttfamily\scriptsize /\allowbreak agent/\allowbreak api/\allowbreak contexts} & {\ttfamily\scriptsize api\_contexts} \\
\bottomrule
\end{longtable}

\subsection{genetic.py --- blueprint \cod{genetic}}
\begin{longtable}{@{}L{17mm}L{74mm}L{45mm}@{}}
\toprule
\textbf{Metodi} & \textbf{Percorso} & \textbf{Funzione} \\
\midrule\endfirsthead
\toprule
\textbf{Metodi} & \textbf{Percorso} & \textbf{Funzione} \\
\midrule\endhead
{\scriptsize\sffamily GET} & {\ttfamily\scriptsize /\allowbreak genetic} & {\ttfamily\scriptsize page} \\
\rigaTenue {\scriptsize\sffamily GET} & {\ttfamily\scriptsize /\allowbreak genetic/\allowbreak api/\allowbreak config} & {\ttfamily\scriptsize api\_config} \\
{\scriptsize\sffamily POST} & {\ttfamily\scriptsize /\allowbreak genetic/\allowbreak api/\allowbreak job/\allowbreak <job\_id>/\allowbreak installa} & {\ttfamily\scriptsize api\_installa\_respinta} \\
\rigaTenue {\scriptsize\sffamily GET} & {\ttfamily\scriptsize /\allowbreak genetic/\allowbreak api/\allowbreak jobs} & {\ttfamily\scriptsize api\_jobs} \\
{\scriptsize\sffamily GET} & {\ttfamily\scriptsize /\allowbreak genetic/\allowbreak api/\allowbreak jobs/\allowbreak <job\_id>} & {\ttfamily\scriptsize api\_job} \\
\rigaTenue {\scriptsize\sffamily POST} & {\ttfamily\scriptsize /\allowbreak genetic/\allowbreak api/\allowbreak jobs/\allowbreak <job\_id>/\allowbreak cancel} & {\ttfamily\scriptsize api\_job\_cancel} \\
{\scriptsize\sffamily POST} & {\ttfamily\scriptsize /\allowbreak genetic/\allowbreak api/\allowbreak run} & {\ttfamily\scriptsize api\_run} \\
\rigaTenue {\scriptsize\sffamily GET} & {\ttfamily\scriptsize /\allowbreak genetic/\allowbreak api/\allowbreak strategies} & {\ttfamily\scriptsize api\_strategies} \\
\bottomrule
\end{longtable}

\subsection{auth.py --- dentro \cod{init\_auth(app)}}
\begin{longtable}{@{}L{17mm}L{74mm}L{45mm}@{}}
\toprule
\textbf{Metodi} & \textbf{Percorso} & \textbf{Funzione} \\
\midrule\endfirsthead
\toprule
\textbf{Metodi} & \textbf{Percorso} & \textbf{Funzione} \\
\midrule\endhead
{\scriptsize\sffamily GET} & {\ttfamily\scriptsize /\allowbreak auth/\allowbreak google/\allowbreak callback} & {\ttfamily\scriptsize auth\_google\_callback} \\
\rigaTenue {\scriptsize\sffamily GET POST} & {\ttfamily\scriptsize /\allowbreak login} & {\ttfamily\scriptsize auth\_login} \\
{\scriptsize\sffamily GET} & {\ttfamily\scriptsize /\allowbreak login/\allowbreak google} & {\ttfamily\scriptsize auth\_google\_start} \\
\rigaTenue {\scriptsize\sffamily GET} & {\ttfamily\scriptsize /\allowbreak login/\allowbreak google/\allowbreak pronto} & {\ttfamily\scriptsize auth\_google\_pronto} \\
{\scriptsize\sffamily GET} & {\ttfamily\scriptsize /\allowbreak logout} & {\ttfamily\scriptsize auth\_logout} \\
\rigaTenue {\scriptsize\sffamily GET POST} & {\ttfamily\scriptsize /\allowbreak register} & {\ttfamily\scriptsize auth\_register} \\
{\scriptsize\sffamily GET POST} & {\ttfamily\scriptsize /\allowbreak reset} & {\ttfamily\scriptsize auth\_reset} \\
\rigaTenue {\scriptsize\sffamily GET POST} & {\ttfamily\scriptsize /\allowbreak verify} & {\ttfamily\scriptsize auth\_verify} \\
\bottomrule
\end{longtable}

\subsection{ui\_ext\_a.py --- estensione v2}
\begin{longtable}{@{}L{17mm}L{74mm}L{45mm}@{}}
\toprule
\textbf{Metodi} & \textbf{Percorso} & \textbf{Funzione} \\
\midrule\endfirsthead
\toprule
\textbf{Metodi} & \textbf{Percorso} & \textbf{Funzione} \\
\midrule\endhead
{\scriptsize\sffamily POST} & {\ttfamily\scriptsize /\allowbreak api/\allowbreak a/\allowbreak metric\_preview} & {\ttfamily\scriptsize a\_metric\_preview} \\
\rigaTenue {\scriptsize\sffamily GET} & {\ttfamily\scriptsize /\allowbreak api/\allowbreak a/\allowbreak overview} & {\ttfamily\scriptsize a\_overview} \\
{\scriptsize\sffamily GET} & {\ttfamily\scriptsize /\allowbreak api/\allowbreak a/\allowbreak stream/\allowbreak <name>/\allowbreak meta} & {\ttfamily\scriptsize a\_stream\_meta} \\
\bottomrule
\end{longtable}

\subsection{ui\_ext\_b.py --- estensione v2}
\begin{longtable}{@{}L{17mm}L{74mm}L{45mm}@{}}
\toprule
\textbf{Metodi} & \textbf{Percorso} & \textbf{Funzione} \\
\midrule\endfirsthead
\toprule
\textbf{Metodi} & \textbf{Percorso} & \textbf{Funzione} \\
\midrule\endhead
{\scriptsize\sffamily GET} & {\ttfamily\scriptsize /\allowbreak ensemble-custom} & {\ttfamily\scriptsize b\_ensemble\_custom\_alias} \\
\rigaTenue {\scriptsize\sffamily GET} & {\ttfamily\scriptsize /\allowbreak projects/\allowbreak <path:pid>} & {\ttfamily\scriptsize b\_project\_detail} \\
\bottomrule
\end{longtable}

\subsection{ui\_ext\_e.py --- estensione v2}
\begin{longtable}{@{}L{17mm}L{74mm}L{45mm}@{}}
\toprule
\textbf{Metodi} & \textbf{Percorso} & \textbf{Funzione} \\
\midrule\endfirsthead
\toprule
\textbf{Metodi} & \textbf{Percorso} & \textbf{Funzione} \\
\midrule\endhead
{\scriptsize\sffamily GET} & {\ttfamily\scriptsize /\allowbreak api/\allowbreak e/\allowbreak account-data} & {\ttfamily\scriptsize e\_account\_data} \\
\rigaTenue {\scriptsize\sffamily GET} & {\ttfamily\scriptsize /\allowbreak api/\allowbreak e/\allowbreak sell-status} & {\ttfamily\scriptsize e\_sell\_status} \\
\bottomrule
\end{longtable}

\subsection{ui\_ext\_f.py --- estensione v2}
\begin{longtable}{@{}L{17mm}L{74mm}L{45mm}@{}}
\toprule
\textbf{Metodi} & \textbf{Percorso} & \textbf{Funzione} \\
\midrule\endfirsthead
\toprule
\textbf{Metodi} & \textbf{Percorso} & \textbf{Funzione} \\
\midrule\endhead
{\scriptsize\sffamily GET} & {\ttfamily\scriptsize /\allowbreak api/\allowbreak f/\allowbreak account} & {\ttfamily\scriptsize f\_account} \\
\rigaTenue {\scriptsize\sffamily GET} & {\ttfamily\scriptsize /\allowbreak api/\allowbreak f/\allowbreak subscription/\allowbreak invoices} & {\ttfamily\scriptsize f\_subscription\_invoices} \\
\bottomrule
\end{longtable}

\begin{nota}{Quanti sono}
Questa tabella elenca \num{246} percorsi distinti, da \num{267} decoratori \cod{@route} su \num{11} moduli. Se i numeri non coincidono con quelli del progetto che si ha davanti, la tabella è vecchia: si rigenera con \file{strumenti/genera-rotte.py}.
\end{nota}



\backmatter
\printindex

\end{document}
